\documentclass[12pt,a4paper]{article}

\usepackage{mathtools}
\usepackage{mathrsfs}
\usepackage{float}
\usepackage{caption}
\usepackage{tikz}
\usepackage{tikz-cd}
\usetikzlibrary{arrows,positioning,calc,decorations.pathmorphing}

%% Current project v3.070 source-reference bindings; previous macro bodies unchanged.
% Current II-D reference root for PapersMinimal v3.063; mathematical macros otherwise unchanged.
% Macros-v2_188: re-based current paper identities and A--G family; no change to geometry.
% Macros-v2_184: active Paper III-E v1.200/III-F v1.267 binding for framework v3.031; mathematical macros unchanged.
% Macros-v2_183: updated shared cross-reference bindings for Paper1B/III-E/III-F v3.030; mathematical macros unchanged.
% Macros-v2_182: active Paper II-F/V-A cross-references for framework v3.029; definitions unchanged.
% Paper-series v2.381: new I-B torsion/modular results, bibliography and current source roots.
% Paper-series v2.380: new exceptional bibliography and synchronized paper references.
% Paper-series v2.379 current 31-paper macro/aux registry.
% Paper-series v2.378: current 31-paper AUX/bibliography registry updated for Paper I-B v1.2 literature additions.
% Framework v2.371: current-source reference registry refreshed; old theorem labels retained.
% Framework v2.370: current 30-paper source/AUX root rebinding; native mathematical macros unchanged.
% Framework v2.368: synchronized merged 30-root registry; mathematical notation unchanged.
% Framework v2.358: active 30-paper source registry synchronized; existing mathematical macros unchanged.
% Framework v2.357: selective new source-owned certificates; active reference roots updated for five revised papers.
% Framework v2.356: active 30-root registry synchronized to the 30 papers in this release.
% Framework v2.275: exact positive dense-twin word cubature, finite repeatability and eighth-moment classification, nested sampling frames, joint action and POVM-jet exports; preserves v2.274.
% Framework v2.274: filename and root-source synchronization.
% Framework v2.273: restored missing Problem 3 theorems from v2.268; exact exceptional frames, twin geometry, convolution repeatability and projector/occurrence firewalls; merged on v2.272 without overwriting independent branches.
% Framework v2.272: merged v2.270 covariant native-readout mathematics with v2.271 Problem 4 finite-action Hopf hypercharge/holonomy corrections; no universal bulk-gauge uniqueness claim.
% Independent branches: v2.270 readout/guard exports and v2.270 K16 action-consumed charge closure; all prior scoped assumptions retained.
% Framework v2.269: synchronized 30 active AUX roots; mathematics and bibliography preserved from Macros-v2_161.
% Framework v2.266: Paper II-C converse-derived outcome resources; II-E, III-C, IV-A/B, V-A and VI-B downstream mathematical closures; 30-root registry synchronized.
% Framework v2.265: II-C coordinate-stage constitution and exact first-symbol host inheritance; prior mathematical domains retained.
% Framework v2.264: paper II-B reverse/audit and IV-A observable-minimal exchange; synchronized external references to 30 active roots.
% Framework v2.263: Paper I-C adds bridge-controlled root-carry parity and source-relative old-cap cancellation; other root mathematical bodies retained; active 30-root references synchronized.
% Framework v2.262: integrates Paper II-F radius, Casimir and readout results; active 30-root references synchronized.
% Framework v2.261: merged Paper II-B GP2-level/geometric repairs; preserves v2.260 II-D repairs; active references synchronized.
% Framework v2.260: Paper II-D signed RS, finite-symbol typing, and conditional curvature/Jordan repairs; references synchronized.
% Framework v2.259: conditional native SU(2) atlas and explicit later discharge; v2.258 mathematical bodies retained elsewhere.
% Framework v2.258: preserves parallel II-A/dynamics update and adds complete readout/frequency closures.
% Framework v2.257: exact minimal complete position readout and arbitrary-preparation frequency law.
% Framework v2.256: exact native inertia and sharp bounded ordered-source closure; preserves concurrent v2.255 work.
% Framework v2.255: preserves parallel detector integration and adds exact finite readout closures.
% Framework v2.254: Pasterski bibliography and synchronized 30-paper external-reference registry.
% Framework v2.253: typed raw coordinate defect in I-A and synchronized VI-A import; active references synchronized.
% Framework v2.252: synchronized source registry; optional III-G shell-frame calculation.
% Framework v2.251: forward common-action execution, joint source jets and scoped coupled control.
% Framework v2.251: exact full-heat variations, fermion source jets and explicit physical completion frontier.
% Framework v2.251: certified interacting source jets, pole stability and scoped completion proofs.
% Framework v2.251: synchronized active source registry; mathematical notation unchanged.
% Framework v2.251: synchronized active source registry; mathematical notation unchanged.
% Framework v2.251: synchronized active source registry; mathematical notation unchanged.
% Framework v2.251: synchronized source registry; no mathematical notation changes.
% Framework v2.244: synchronized 30-paper AUX registry and rooted phase-source action-order tests.
% Framework v2.242: signed rooted phase lift, complete transverse return and exact family cubic.
% Framework v2.241: proved rooted conormal geometry, conditional cap realization and scoped action tests.
% Framework v2.239: all 30 active source roots synchronized; no mathematical macro changes.
% Framework v2.238: Paper IV A/B/C split and 30-root reference registry.
% Framework v2.237: Paper VII-D technical split and 28-root reference registry.
% Framework v2.236: Paper VI-C technical-volume reference added; all current root bindings advanced.
% Framework v2.235: 26-root AUX and bibliography bindings synchronized; mathematical macros unchanged.
% BEGIN INLINE SOURCE: Paper00-TeXBuildCompatibility-v1_0.tex
% Rendering-only repairs for inherited long appendices and bitmap text fonts.
% Integer counters through 26 retain the original labels; later values use AA,
% AB, ... . This file changes no action coefficient, state or theorem premise.
\RequirePackage{alphalph}
\makeatletter
\let\BHPrimitiveAlph\@Alph
\newalphalph{\BHExtendedAlph}{\BHPrimitiveAlph}{26}
\let\@Alph\BHExtendedAlph
\AtBeginDocument{%
  \@ifpackageloaded{microtype}{\microtypesetup{expansion=false}}{}%
}
\makeatother

% END INLINE SOURCE: Paper00-TeXBuildCompatibility-v1_0.tex

% Framework v2.164: current reference/bibliography profile; mathematical definitions retained.
% Framework v2.162: finite metric, functional Hilbert and scoped Ward integration.
% Framework v2.161: inherits Macros-v2_57.tex exactly except the current external-reference binding.
% Macros v2.57 / framework v2.156: current reference registry only; mathematical definitions unchanged.
% Framework v2.143 / Macros v2.55: current external-reference bindings; mathematical notation unchanged.
% Framework v2.110 / Macros v2.54: current external-reference registry and bibliography; mathematical notation inherited.
% Framework v2.67 / Macros-v2_53: adds Born-action-cell and joint-EH readout symbols; absolute unit realization is metrological; MATCH_{J/scal}^{EH} remains the sole G24 normalization bridge.
% Framework v2.35 / Macros-v2_52: grouped paper numbering, active AUX registry, and Born-v6_20 bibliography binding synchronized; notation unchanged.
% Framework v2.26: active-source AUX and bibliography bindings synchronized; notation unchanged.
% Framework v2.09: reviewed early-reciprocity integration; scoped physical imports.
% Framework v2.08: Born--Euler mathematics and conditional complete-action potential criterion.
% Framework v2.07: active AUX registry includes Paper I-B v1.2; shared bibliography Born-v6_18.
% Macros-v2_46 / framework v2.04: mathematical atlas and downstream realization correction.
% Macros-v2_46 / framework v2.03: scoped structure integration and current dependency bindings.
% Macros-v2_46 / framework v2.01: shared source bindings; prior revision comments are historical.
% v2.42 / framework v1.99: Synchronized RG/AS owner and bibliography bindings.
% v2.41 / framework v1.99: active manuscript AUX bindings and geometric-closure bibliography synchronized.
% Release-v1.80 centralized source synchronization for scalar-circle/comparison-band integration.
% Release v1.79: Hurwitz hierarchy and kinetic-normalized action response integration; physical mass matching remains conditional.
% Release v1.78: approved M8 closure-batch integration; physical and analytic domains retained.
% Current v1.78 registry: all active TeX roots are synchronized below.
% Earlier release comments are historical scope records, not current status.
% Release-v1.77 synchronization: native same-owner Noether closure, generator--flow/history typing, analytic-certification firewalls, scalaron/K16, LF/QCD, quantum-Ward, and RAR boundary-owner transport are registered; shared mathematical notation is unchanged.
% Release-v1.76 synchronization: registry points to the v1.76 active paper roots.
% Release-v1.76 synchronization: registers the local A+Phi+g Noether/constraint closure and dynamics-frontier propagation; shared mathematical notation is unchanged.
% Release-v1.75 synchronization: registry points to the v1.75 active paper roots.
% Release-v1.75 synchronization: integrates the joint-Dirac Einstein bulk real form and the local odd-scalar heat-coefficient action on explicit analytic first-jet/interior domains; v1.74 Yang--Mills closure is retained. Fermionic reality, pure-metric curvature-squared/exact-nonlocal terms, noncanonical boundaries, physical vacuum selection, and actual solution-history tokens remain open.
% Release-v1.74 synchronization: registry points to the v1.74 active paper roots and retains the v1.73 theorem labels.
% Release-v1.74 synchronization: canonical Macros-v2_34.tex; generated local nonlinear Yang--Mills action, constrained connection-field phase space, source-free Gauss momentum map, and Legendre closure are integrated on their proved domains. The remaining non-gauge ACT_cl sectors and actual solution-history selection remain open.
% Release-v1.71 synchronization: closes the former JOINT_x/JOINT_5
% connection-occurrence gap by canonical structure-group extension of the
% nonflat Born connection, faithful associated representation, and the local
% cotangent--Clifford operator constructor.  Actual fluctuation members,
% actions, vacua, solutions, calibrations, and analytic QFT domains remain
% separately typed.  The active-AUX registry is advanced once.
% Release-v1.71 canonical macro and active-AUX registry.
% ======================================================================
% Macros-v2_39.tex -- canonical publication macro file
% Release v1.71 retains the v1.70 null-ray/QCD boundaries and additionally
% removes the duplicate structural connection actuality premise.  It does not
% claim global bundle topology, select a gauge/Higgs member, or promote the
% finite construction to an interacting continuum theory.
% Release v1.68 constructs the finite light-front total-fraction operators and
% parton number/momentum operator measures, proves their simplex and Hausdorff-limit
% theorems, and separates positive cut measures from scheme-renormalized
% light-ray distributions.  It also separates finite real Schur poles, stable
% norm-resolvent/Riesz pole-residue limits, and continued resonance
% pole/residue limits.  The Paper-VII-C scattering package is claim-indexed:
% local QFT and spectral claims import no S-matrix or infrared premise, while
% asymptotic amplitudes, LSZ, S-matrix poles, completeness, and cross sections
% import only their named massive or massless/long-range clauses.
% Release v1.67 factorizes BND48_H into a same-action Hessian/Schur
% descendant and still-open typed boundary-trace, critical-lift, and
% refinement-limit conditions.  It proves the finite-collar Dirichlet-to-
% Neumann law sqrt(y)tanh(L sqrt(y)), its nonuniform critical limit, and the
% finite-return ultra-infrared turnover.  The Paper-I-A Hardy constant mode is
% an exact geometric near-hit, not yet the physical metric boundary mode.
% The current local f(R) Hessian remains massive and therefore still refutes
% a RAR prediction on the minimal action branch.
% Release v1.66 derives the transverse depth-six normalizer bridge from the
% already generated C6 Reynolds image and Paper-I-A binary-dihedral normalizer.
% The bridge is the unique least 2D3-stable scalar-containing realization of
% the faithful quotient 2D3/C6=C2, has trace defect six and hence K=48, but is
% transverse to the nested active carrier.  Exact K=44 and K=52 contraction,
% semigroup, and energy-midpoint identities are registered as structural
% readouts rather than projector events.  The active E48 carrier gap remains
% unchanged.  A same-action boundary-symbol/Calderon intertwiner is isolated
% as the necessary and still unproved owner bridge to the static DtN sort;
% representation equivalence alone neither derives RAR nor creates a K=48
% action junction, pole, occupation, or cosmological transition.
% Release v1.65 derives the native E_48-centered carrier mass gap from the
% GP1--GP3-generated active spectrum, distinguishes its carrier-free interval
% from a particle located at E_48, and withdraws the unsupported promotion to
% a positive continuum Yang--Mills mass gap.  The latter remains a valid
% conditional transfer theorem on an enlarged continuum branch, but it is not
% an active framework prediction or an unfinished native-gap proof.
% Release v1.64 synchronizes Derivation Batches 14--22.  It registers the
% GP1--GP3 heat-profile orbit, same-Hessian electroweak response and pole
% ownership, K16/K48 action-response corrections, perturbation state and
% non-Gaussianity typing, finite-vacuum/gap/QCD limit repairs, factorized
% position-area ownership, the actual static metric response, and the ordered
% Wilson/string-tension limits.  Strict continuum convergence, positive
% physical gap, surviving higher cumulants, and positive tension remain open.
% Release v1.61 closes TIMEWIDTH16 as an independent direct-production datum.
% JUNC_{2|7} continues to own the finite intrinsic microscopic profile/width,
% while the actual same-action background/production history determines its
% physical-time pullback, support, and duration.  The current Paper-I-A--VII-C AUX
% registry is advanced once; shared mathematical notation is unchanged.
% Release v1.60 derives the final-octet primitive-remnant depletion, canonical
% Hopf-current representative, exact Hopf descent to dOmega/(4pi), and
% closed-FLRW a^{-3} neutral-current shape.  It retypes K=48 as the onset
% of the final one-unit geometric transfer and records cDM/RAR as distinct
% physical projections of that common neutral parent while preserving
% occurrence, Hilbert-owner, stress-matching, and static-kernel premises.
% Release v1.59 transports the v1.58 exceptional-shell/K16 closure through
% Papers III-F and VI-B--VII-C.  It replaces Paper VI-B's obsolete mixed-Hessian
% output typing by the third-variation pair vertex and Bogoliubov occupation
% support, derives the post-seed Higgs-pair feedback vertex and its exact
% alpha--delta kernel factors, propagates the unique K16/K48 action-junction
% split, identifies the QCD delta count with the upstream 2O shell atom only
% as a common-scalar numerical compatibility, and narrows THERM16 to dynamical
% approach/beta selection because the finite Gibbs/KMS target is already
% unique on a supplied simple observable algebra once its Hamiltonian and
% inverse temperature are supplied; direct-sum central weights remain typed.
% Release v1.58 registers the exceptional nearest-shell closure, the derived
% transverse K=16 action event, the unique interior action junction, the
% type-correct ordered pair vertex, the thin-layer squeeze law, and the
% radiation-like/matter-like kinematic split.  REG_0, transfer magnitude,
% physical layer width, and thermalization remain explicitly conditional.
% Release v1.57 registers the active-depth-six obstruction, the native
% K=48 carrier desert, the typed 7|9 owner reduction, and the dynamical
% reheating handoff across Papers III-B--VI-B.
% Release v1.56 registers the least-support n_2 inflation dressing of the
% persistent n_7 Laurent response, the analytic global-convexity and positive
% endpoint-measure proofs, the derived inflation core and K=16 coefficient
% freeze, and the no-particle dark-sector split into static RAR plus the
% conditional homogeneous effective-curvature response at reflected K=48.
% There is no n_6 action, scalaron-DM mass, or second Einstein--Hilbert term.
% Release v1.55 registers the protected-octet resolved/residual reflection
% q_sp+Z_W=8, its j<->8-j and K<->64-K scale covariance, and the derived
% rank-one-exit quadratic grade C_1^2/U=C_2 in Papers III-C, III-E, and VIII.
% The reflection is a typed trace/scale theorem, not a time reversal or a
% physical event-occurrence theorem; reheating transfer remains conditional.
% Release v1.54 registers the exact positive-part Laurent indices, the
% active-domain/terminal-endpoint action split, the no-second-EH firewall, and
% the q_DM=6/K_DM=48 scalaron readout typing across Papers III-D/E/F, VI,
% VIII, and IX.
% Release v1.53 registers the common nonnegative Laurent obstacle projection
% for n_1^+ and n_7^+, discharges the three former exit/freeze assumptions,
% and advances the Paper-III-D/E/F, VIII, and IX consumers.  Shared notation
% and bibliography ownership are unchanged.
% Release v1.51 registers the PQTI derivation batch and its updated
% Paper-III-A/C/D/E/F, VIII, and IX consumers.  Shared notation and
% bibliography ownership are unchanged.
% Release v1.50 advances the Paper-III-B AUX target to v1.40 after the
% rigorous audit repair; shared notation and bibliography ownership are unchanged.
% Release v1.49 advances the Paper-II-E/F AUX target to v3.14 after the
% minimal audit repair; shared notation and bibliography ownership are unchanged.
% Release v1.46 synchronizes Paper II-A's minimal frozen constructor-origin
% theorem and its five dependent AUX targets.  Shared notation is unchanged.
% Release v1.45 updates the Paper-II-B AUX target to v2.23 after the
% GP3-rigidity and appendix-placement repair; shared notation is unchanged.
% Release v1.44 advances the centralized Paper-I-A--VII-C AUX registry after the
% Paper-I-A v5.16 protected-completion scope repair.  Every current
% paper imports this one file; no paper-level opt-in or \externaldocument list
% is permitted.  The bibliography remains centralized here through Born-v6_13.
% Shared bibliography, theorem infrastructure, colors, notation alignment,
% and current Paper-I-A--VII-C external-reference registration.
% This file centralizes common imports, bibliography sorting, colors, theorem
% infrastructure, and the notation shared by the complete series.
% v2.9 (release v1.43) made the Paper-I-A--VII-C AUX registry automatic and
% centralized the bibliography and all external-reference ownership.
% v2.8 release-registry refresh (v1.33) registers the rooted generated-parent
% cascade and the POL/REG/free-law/11+4/Born-group closure revisions.  No
% shared mathematical notation or Paper-II-A fundamental principle is changed.
% v2.8 release-registry refresh (v1.32) registers the universal generated-
% continuation revisions in Papers II-A, III-A, III-C, IV, and VIII.  No
% shared mathematical notation or Paper-II-A fundamental principle is changed.
% v2.8 release-registry refresh (v1.30) registers the CRP4/projective-state
% source revisions and their direct registry-only consumers.  No shared
% mathematical notation or Paper-II-A axiom is changed.
% v2.7 release-registry refresh (v1.27) updates only the current external
% Paper-I-A--VII-C basenames after the placement/assumption synchronization; no
% shared mathematical notation or Paper-II-A axiom is changed.
% v2.7 synchronizes the geometric defect-repair revision: Paper-II-A v3.42,
% Paper-III-C v1.26, Paper-III-D v1.27, Paper-III-E v1.32, and Paper-III-F
% v1.34.  It registers the common trace-defect/spectral-filtration theorem,
% its owner-typed response pullback, and the reduced GRED and TERM interfaces.
% No shared mathematical notation is changed.
% v2.6 synchronizes Paper-III-E v1.31's cone-type repair, minimal-response
% classification, and explicit classical action-role boundary, together with
% Paper-III-F v1.32's downstream response-atlas synchronization.  No shared
% mathematical notation is changed.
% v2.5 synchronizes the Paper-III-D v1.26 faithful-trace-duality proof,
% exact ordered-owner joint locus, and typed alternative-presentation
% bookkeeping repair, together with its Paper-III-E and closure-frontier
% references.  No shared mathematical notation is changed.
% v2.4 synchronizes the Paper-III-C v1.25 orbit/process/owner repair, its
% downstream Papers III-B--F, V, VI, VIII, and IX revisions, and Born-v6_13.
% v2.3 synchronizes the Paper-II-A v3.41 witness/support/principal inheritance
% layer, the Paper-III-A v1.24 joint-realizability/readout firewall, all dependent
% host-rate revisions, and the current Paper-I-A--VII-B source registry.
% v2.2 synchronizes the generated active-realization, determinant-depth,
% dimension-weighted hierarchy, and typed cosmological-event interfaces.
% v2.1 synchronizes the exceptional-support interface revisions: Paper I-A
% v5.9, Paper II-A v3.30, Paper II-E/F v3.7, and Paper IV v4.16.
% v2.0 introduces the stage-safe G0 occurrence notation, synchronizes the
% active Paper-I-A--IV source registry (current rigorous Paper-III closure:
% III-A v1.17, III-B v1.22, III-C--E v1.17, III-F v1.14, IV v4.14), and makes
% external-reference self-skip depend on logical paper role rather than a
% versioned job name.
% ======================================================================

% ======================================================================
% macros.tex — Shared macros for the Born Geometry series
% ======================================================================

% ======================================================================
% 1. PACKAGES, BIBLIOGRAPHY, COLORS, AND UNIT DEFINITIONS
% ======================================================================
% Core packages shared across the complete Paper-I-A--VII-C series.  Loading them here makes
% the macro file a genuine series preamble while remaining compatible with existing
% paper preambles that load the same packages before \input.
\usepackage{amsmath}
\usepackage{amssymb}
\usepackage{amsthm}
\usepackage{mathrsfs}
\usepackage{array}
\usepackage{booktabs}
\usepackage{xcolor}
\makeatletter
\@ifpackageloaded{geometry}{}{\usepackage{geometry}\geometry{margin=1in}}
\makeatother
\usepackage{mdframed}

% Shared support packages required by the macros below and by most papers.
\usepackage{ragged2e}
\usepackage{longtable}
\usepackage{tabularx}
\usepackage{enumitem}
\usepackage[most]{tcolorbox}

% Series-wide numeric bibliography convention using the portable BibTeX path.
% All Papers I-A--VII-C use the ordinary \cite interface.
\usepackage[numbers,sort&compress]{natbib}
\providecommand{\printbibliography}{%
  \bibliographystyle{plainnat}%
  \bibliography{Born-v6_59}%
}

% Shared colors.  \providecolor avoids clobbering a paper-specific override.
\providecolor{postulatecolor}{RGB}{230,242,255}
\providecolor{compare}{RGB}{230,245,230}
\providecolor{definitioncolor}{RGB}{230,242,255}

\providecolor{openriskcolor}{RGB}{255,235,235}
\providecolor{bridgecolor}{RGB}{240,240,255}
\providecolor{gold}{RGB}{255,215,0}

% ======================================================================
% 2. CUSTOM ENVIRONMENTS
% ======================================================================
% Common table column types used in all papers.  Guarded so that a paper may
% define its own versions before loading this file.
\makeatletter
\@ifundefined{NC@find@L}{%
  \newcolumntype{L}[1]{>{\RaggedRight\arraybackslash}p{#1}}%
}{}
\@ifundefined{NC@find@C}{%
  \newcolumntype{C}[1]{>{\Centering\arraybackslash}p{#1}}%
}{}
\@ifundefined{NC@find@Y}{%
  \newcolumntype{Y}{>{\RaggedRight\arraybackslash}X}%
}{}
\makeatother

% ======================================================================
% 3. THEOREM STYLES AND ENVIRONMENTS
% ======================================================================
% Custom underlined theorem style (useful for remarks or notes)
\newtheoremstyle{underlined}
  {3pt}   % Space above
  {3pt}   % Space below
  {}      % Body font (empty means normal/roman text, typical for remarks)
  {}      % Indent amount
  {}      % Theorem head font (can be \itshape, \bfseries, etc.)
  {.}     % Punctuation after theorem head
  { }     % Space after theorem head
  {\underline{\thmname{#1}\thmnumber{ #2}}\thmnote{ (#3)}} % Theorem head formatting

% Standard plain theorem style environments.  Each environment is guarded so
% the shared macro file can be loaded after a paper-specific theorem setup.
\theoremstyle{plain}
\makeatletter
\@ifundefined{theorem}{\newtheorem{theorem}{Theorem}[section]}{}
\@ifundefined{lemma}{\newtheorem{lemma}{Lemma}[section]}{}
\@ifundefined{proposition}{\newtheorem{proposition}{Proposition}[section]}{}
\@ifundefined{corollary}{\newtheorem{corollary}{Corollary}[section]}{}
\makeatother

% Definition style environments (typically non-italicized body)
\theoremstyle{definition}
\makeatletter
\@ifundefined{definition}{\newtheorem{definition}{Definition}[section]}{}
\@ifundefined{postulate}{\newtheorem{postulate}{Postulate}[section]}{}
\@ifundefined{closure}{\newtheorem{closure}{Closure}[section]}{}
\makeatother

% Applying the custom underlined style
\theoremstyle{underlined}
\makeatletter
\@ifundefined{remark}{\newtheorem{remark}{Remark}[section]}{}
\@ifundefined{convention}{\newtheorem{convention}{Convention}[section]}{}
\@ifundefined{principle}{\newtheorem{principle}{Principle}[section]}{}
\makeatother

% Framed proof environment.
% Every \begin{proof} ... \end{proof} is automatically framed/shaded.
% This avoids mdframed/trivlist nesting, which can drop the top border after
% theorem-like environments.
\makeatletter
\renewenvironment{proof}[1][\proofname]{%
  \par\addvspace{6pt}%
  \noindent
  \begin{tcolorbox}[%
    enhanced,
    breakable,
    colback=black!4,
    colframe=black!35,
    boxrule=0.35pt,
    arc=1pt,
    left=6pt,
    right=6pt,
    top=4pt,
    bottom=4pt,
    boxsep=0pt,
    before skip=0pt,
    after skip=6pt,
    parbox=false
  ]%
  \pushQED{\qed}%
  \normalfont\itshape #1\@addpunct{.}\enspace
  \normalfont\ignorespaces
}{%
  \unskip\nobreak\hfill\popQED
  \end{tcolorbox}%
  \par\addvspace{2pt}%
}
\makeatother

\mdfdefinestyle{postulatebox}{
  backgroundcolor=yellow!10,
  linecolor=black,
  linewidth=0.4pt,
  skipabove=\baselineskip,
  skipbelow=\baselineskip,
  nobreak=true
}

% ======================================================================
% PROOF-READABILITY HELPERS
% ======================================================================

\newcommand{\ProofPlan}[1]{%
  \par\smallskip
  \noindent\textbf{Proof plan. }#1%
  \par\smallskip
}

\newlist{proofsteps}{enumerate}{2}
\setlist[proofsteps,1]{
  label=\textbf{(\alph*)},
  ref=\alph*,
  leftmargin=2.2em,
  itemsep=0.45\baselineskip,
  topsep=0.35\baselineskip
}
\setlist[proofsteps,2]{
  label=\textbf{(\roman*)},
  ref=\roman*,
  leftmargin=2.0em,
  itemsep=0.25\baselineskip,
  topsep=0.2\baselineskip
}


% ======================================================================
% 4. NUMBER SYSTEMS AND ALPHABETS
% ======================================================================
\providecommand{\RR}{\mathbb{R}} % Real numbers
\providecommand{\PP}{\mathbb{P}} % Projective space / Probability
\providecommand{\CC}{\mathbb{C}} % Complex numbers
\providecommand{\ZZ}{\mathbb{Z}} % Integers
\providecommand{\NN}{\mathbb{N}} % Natural numbers
\providecommand{\HH}{\mathbb{H}} % Quaternions / Hyperbolic space
\providecommand{\LFRayGeom}{\mathsf{LF\!\mbox{-}RAY\!\mbox{-}GEOM}}
\providecommand{\LFRep}{\mathsf{LF\!\mbox{-}REP}}
\providecommand{\LFState}{\mathsf{LF\!\mbox{-}STATE}}
\providecommand{\LFLim}{\mathsf{LF\!\mbox{-}LIM}}
\providecommand{\LFScat}{\mathsf{LF\!\mbox{-}SCAT}}
% Number systems / projective spaces
\providecommand{\CP}{\mathbb{CP}}
\providecommand{\RP}{\mathbb{RP}}

% ======================================================================
% 5. GROUPS AND LIE THEORY
% ======================================================================
\providecommand{\SO}{\ensuremath{\mathrm{SO}}}     % Special Orthogonal group
\providecommand{\SU}{\ensuremath{\mathrm{SU}}}     % Special Unitary group
\providecommand{\Spin}{\ensuremath{\mathrm{Spin}}} % Spin group
\providecommand{\CS}{\ensuremath{\mathrm{CS}}}     % Chern-Simons (capitalized)
\providecommand{\cs}{\ensuremath{\mathrm{cs}}}     % Chern-Simons (lowercase)

% ======================================================================
% 6. COMMON MANIFOLDS AND TOPOLOGY
% ======================================================================
\providecommand{\Sone}{S^1}                 % Standard circle

% ======================================================================
% 7. MATH OPERATORS (Declared via \DeclareMathOperator)
% ======================================================================
\DeclareMathOperator{\End}{End}   % Endomorphisms
\DeclareMathOperator{\Aut}{Aut}   % Automorphisms
\DeclareMathOperator{\spec}{spec} % Lowercase spectrum
\DeclareMathOperator{\tr}{tr}     % Lowercase trace (if you already have \Tr)
\DeclareMathOperator{\supp}{supp} % Support
\DeclareMathOperator{\TrA}{Tr_{A}} % Partial trace over subsystem A
\DeclareMathOperator{\TrB}{Tr_{B}} % Partial trace over subsystem B
% Lie/algebra operators
\DeclareMathOperator{\Ad}{Ad}
\DeclareMathOperator{\Sym}{Sym}
\DeclareMathOperator{\su}{su}
% ======================================================================
% 8. STANDARD OPERATOR-LIKE SYMBOLS (Safe definitions via \providecommand)
% ======================================================================
\providecommand{\vol}{\operatorname{vol}}   % Volume form
\providecommand{\Vol}{\operatorname{Vol}}   % Total volume
\providecommand{\Ric}{\operatorname{Ric}}   % Ricci tensor
\providecommand{\Rm}{\operatorname{Rm}}     % Riemann tensor
\providecommand{\Tr}{\operatorname{Tr}}     % Uppercase trace
\providecommand{\Spec}{\operatorname{Spec}} % Uppercase spectrum
\providecommand{\diag}{\operatorname{diag}} % Diagonal matrix
\providecommand{\rank}{\operatorname{rank}} % Rank
\providecommand{\scal}{\mathcal R} % Scalar curvature
\providecommand{\Cl}{\operatorname{Cl}}     % Clifford algebra
\providecommand{\diam}{\operatorname{diam}} % Diameter

% ======================================================================
% 9. FUNCTIONS, TRIGONOMETRY, AND CALCULUS
% ======================================================================

% Codifferential naming conventions
% (Avoid \ddag: commonly defined by physics.sty and symbol packages)
\providecommand{\codiff}{\mathrm{d}^{\dagger}} % Codifferential operator

% Quantum mechanics bra-ket notation (needed for basis tables)
\providecommand{\ket}[1]{\lvert #1 \rangle}
\providecommand{\bra}[1]{\langle #1 \rvert}
\providecommand{\braket}[2]{\langle #1 \mid #2 \rangle}

% ======================================================================
% 10. PHYSICAL SCALES AND CONSTANTS (Paper 1c / Born Geometry)
% ======================================================================
\providecommand{\Rp}{\ensuremath{R_{\mathrm P}}}     % matched common EH curvature scale (not defined from observed G)
\providecommand{\hBorn}{\ensuremath{h_{\mathrm B}}} % physical representative 2\pi\hbar of one primitive Born scalar action cell
\providecommand{\UEHjoint}{\ensuremath{U_{J}}} % joint-Dirac EH curvature-coordinate scale 8 kappa_J
\providecommand{\Gread}{\ensuremath{G_{\rm scal}}} % scalar-action coordinate c_op^3/(h_B U); physical matching explicit
\providecommand{\GreadJoint}{\ensuremath{G_{J,\mathrm{raw}}}} % raw ordinary-trace coordinate = G_scal/96
\providecommand{\Rds}{R_{\mathrm{dS}}}    % de Sitter Curvature
\providecommand{\Mpl}{M_{\mathrm{Pl}}}    % Planck mass
\providecommand{\lp}{\ell_{\mathrm{Pl}}}  % Planck length
\providecommand{\Pmax}{P_{\mathrm{max}}}  % Maximum momentum
\providecommand{\Mplbar}{\bar{M}_{\mathrm{Pl}}} % Reduced Planck mass
\providecommand{\eff}{\mathrm{eff}}       % Subscript for effective quantities
\providecommand{\Cosmo}{\ensuremath{\Lambda}}                % Cosmological constant (use \Cosmo in source)
\providecommand{\LCDM}{\ensuremath{\Cosmo\mathrm{CDM}}}  % \LambdaCDM (typesets with \Cosmo)
\providecommand{\Lcut}{\ensuremath{\Lambda_{\mathrm{cut}}}}  % Cutoff scale (brace when used as script)
\providecommand{\Rcs}{\ensuremath{R_{\mathrm{cs}}}}      % Configuration-space radius
\providecommand{\CosmoEff}{\ensuremath{\Cosmo_{\eff}}}   % Effective cosmological constant

% Disambiguation: \Rmom is the momentum-space (S^3) radius. 
% Reserve plain R for the 4D curvature scalar in f(R).
\providecommand{\Rmom}{\ensuremath{R_{\mathrm{mom}}}} 
\providecommand{\Punit}{P_{0}} % Momentum unit set by \Rmom (= \hbar/\Rmom)
\providecommand{\Rvec}{\ensuremath{R_{\mathrm{vec}}}} % invariant vector-space radius in the reduced round sector
% Specific numeric values and expressions
\providecommand{\Vlogcap}{L_{\mathrm{Hopf}}^{(1)}}
\providecommand{\Hopf}{\operatorname{Hopf}}
% Scalar phase (S^1) sector + universal cover notation (Paper I-A)
\providecommand{\tper}{T_\star}      % Time period/scale
\providecommand{\Eper}{E_\star}      % Energy period/scale
\providecommand{\ttilde}{\tilde t}   % Universal cover time
\providecommand{\Etilde}{\tilde E}   % Universal cover energy

% ======================================================================
% 10A. SERIES-WIDE LOGICAL / NOTATIONAL DISAMBIGUATIONS
% ======================================================================
% Use these symbols to avoid collisions among: (i) the Paper--III
% no-insertion/geometric-closure class, (ii) Paper--I-A trimer diagonals,
% (iii) Paper--V-A/V-C family-space diagonals, (iv) scalar-index running,
% and (v) protected shell-clock labels.
\providecommand{\Rprim}{\ensuremath{\mathsf R_{\mathrm{prim}}^{\Sigma}}} % primitive finite structural readouts
\providecommand{\Rloop}{\ensuremath{\mathsf R_{\circ}^{\Sigma}}} % closed-loop readout role
\providecommand{\Rbdry}{\ensuremath{\mathsf R_{\partial}^{\Sigma}}} % boundary-conditioned readout role
\providecommand{\Rgen}{\ensuremath{\mathsf R_{\mathrm{fin}}^{\Sigma}}} % finite generated closure of primitive roles
\providecommand{\PhysOcc}{\operatorname{PhysOcc}}
\providecommand{\Actual}{\operatorname{Actual}}
\providecommand{\GeomLic}{\operatorname{GeomLic}}
\providecommand{\InvDist}{\operatorname{InvDist}}
\providecommand{\Dcl}{\ensuremath{\mathsf{D}_1}}      % unmarked positive Dirac-even closed spectral-readout form class
\providecommand{\Dheat}{\ensuremath{\mathsf D_{\mathrm{heat}}}} % one-parameter heat-semigroup subclass
\providecommand{\DtriI}{\ensuremath{D^{\mathrm{tri}}_1}} % Paper--I-A cyclic-trimer diagonal
\providecommand{\DtriII}{\ensuremath{D^{\mathrm{tri}}_2}}
\providecommand{\DfamI}{\ensuremath{D^{\mathrm{fam}}_1}} % Family-space diagonal selector
\providecommand{\DfamII}{\ensuremath{D^{\mathrm{fam}}_2}}
\providecommand{\alpharun}{\ensuremath{\alpha_{\mathrm{run}}}} % scalar spectral-index running
\providecommand{\Ksh}{\ensuremath{K_{\mathrm{sh}}}} % protected shell-clock coordinate
\providecommand{\ddef}{\ensuremath{\delta_{\mathrm{def}}}} % universal defect scalar after Paper--III realization

% ======================================================================
% 10B. SERIES-WIDE LOGICAL STATUS TAGS
% ======================================================================
% Upright small-caps wrapper for audit/status labels.  It prevents LaTeX from
% trying to use unavailable small-caps italic font shapes inside theorem bodies
% such as TU/lmr/m/scit.
\providecommand{\AuditLabel}[1]{}
\renewcommand{\AuditLabel}[1]{{\normalfont\scshape #1}}

% These tags are bookkeeping labels for consistency audits. They do not
% create new mathematical hypotheses; they record how a datum is being used
% in the proof chain.
\providecommand{\StatusPrimitive}{}
\renewcommand{\StatusPrimitive}{\AuditLabel{Primitive}}
\providecommand{\StatusDerived}{}
\renewcommand{\StatusDerived}{\AuditLabel{Derived}}
\providecommand{\StatusImported}{}
\renewcommand{\StatusImported}{\AuditLabel{Imported}}
\providecommand{\StatusConditional}{}
\renewcommand{\StatusConditional}{\AuditLabel{Conditional}}
\providecommand{\StatusBridge}{}
\renewcommand{\StatusBridge}{\AuditLabel{Bridge-derived}}
\providecommand{\StatusComparison}{}
\renewcommand{\StatusComparison}{\AuditLabel{Comparison-only}}

% ======================================================================
% 10C. SERIES-WIDE CLAIM-WORDING AUDIT LABELS
% ======================================================================
% These are textual audit labels used to keep modal language consistent
% across the series.  They introduce no new mathematical assumptions.
\providecommand{\ClaimForced}{}
\renewcommand{\ClaimForced}{\AuditLabel{forced}}
\providecommand{\ClaimUnique}{}
\renewcommand{\ClaimUnique}{\AuditLabel{unique}}
\providecommand{\ClaimCanonical}{}
\renewcommand{\ClaimCanonical}{\AuditLabel{canonical}}
\providecommand{\ClaimMinimal}{}
\renewcommand{\ClaimMinimal}{\AuditLabel{minimal}}
\providecommand{\ClaimClosed}{}
\renewcommand{\ClaimClosed}{\AuditLabel{closed}}
\providecommand{\ClaimPrediction}{}
\renewcommand{\ClaimPrediction}{\AuditLabel{prediction}}

% ======================================================================
% 10D. SERIES-WIDE THEOREM-CHAIN AUDIT LABELS
% ======================================================================
% These labels mark the direction of load-bearing theorem chains. They are
% audit annotations only; they introduce no new mathematical assumptions.
\providecommand{\ChainInput}{}
\renewcommand{\ChainInput}{\AuditLabel{Chain input}}
\providecommand{\ChainOutput}{}
\renewcommand{\ChainOutput}{\AuditLabel{Chain output}}
\providecommand{\ChainGuard}{}
\renewcommand{\ChainGuard}{\AuditLabel{Chain guard}}
\providecommand{\ChainNoBackflow}{}
\renewcommand{\ChainNoBackflow}{\AuditLabel{No backflow}}

% ======================================================================
% 10E. SERIES-WIDE LOCAL THEOREM-AUDIT LABELS
% ======================================================================
% These labels mark a statement-level dependency audit around individual
% theorem/definition nodes. They are bookkeeping labels only.
\providecommand{\LocalAudit}{}
\renewcommand{\LocalAudit}{\AuditLabel{Local audit}}
\providecommand{\UsesOnly}{}
\renewcommand{\UsesOnly}{\AuditLabel{Uses only}}
\providecommand{\Exports}{}
\renewcommand{\Exports}{\AuditLabel{Exports}}
\providecommand{\DoesNotUse}{}
\renewcommand{\DoesNotUse}{\AuditLabel{Does not use}}

% ======================================================================
% 10F. SERIES-WIDE NUMERICAL-ROW AUDIT LABELS
% ======================================================================
% These labels classify individual rows in numerical, benchmark, status,
% and comparison tables.  They are audit annotations only and do not change
% any theorem, value, formula, or comparison number.
\providecommand{\RowAudit}{}
\renewcommand{\RowAudit}{\AuditLabel{Numerical-row audit}}
\providecommand{\RowTheorem}{}
\renewcommand{\RowTheorem}{\AuditLabel{Theorem-fixed row}}
\providecommand{\RowBridge}{}
\renewcommand{\RowBridge}{\AuditLabel{Bridge-fixed row}}
\providecommand{\RowRepresentative}{}
\renewcommand{\RowRepresentative}{\AuditLabel{Representative row}}
\providecommand{\RowComparison}{}
\renewcommand{\RowComparison}{\AuditLabel{Comparison row}}
\providecommand{\RowCalibrated}{}
\renewcommand{\RowCalibrated}{\AuditLabel{Calibrated row}}
\providecommand{\RowEstimate}{}
\renewcommand{\RowEstimate}{\AuditLabel{Estimate row}}
\providecommand{\RowOpen}{}
\renewcommand{\RowOpen}{\AuditLabel{Open row}}
\providecommand{\RowNoSelector}{}
\renewcommand{\RowNoSelector}{\AuditLabel{No selector}}

% ======================================================================
% 10G. TABLE HELPERS
% ======================================================================
% Table helpers used by Paper V-C.  The unused readout-label preset has been
% removed; only active table helpers remain.
\providecommand{\gpsection}[1]{%
  \addlinespace[0.55em]
  \multicolumn{3}{@{}l}{\textsc{#1}}\\[-0.25em]
  \midrule
}
\providecommand{\mathcell}[1]{%
  \(\begin{gathered}#1\end{gathered}\)%
}

\newcommand{\DerBy}[1]{%
  \begin{tabular}[t]{@{}l@{}}#1\end{tabular}%
}

\newcommand{\TableNotes}[1]{%
  \par\medskip
  \begin{minipage}{0.94\linewidth}
  \footnotesize\emph{Notes.} #1
  \end{minipage}\par
}
% ======================================================================
% 10H. PAPER VI-A / QCD NOTATION
% ======================================================================
% Shared symbols used by Paper VI-A and downstream QCD/Yang--Mills handoffs.
\providecommand{\Pqcd}{\mathfrak P_{\rm VIII}^{\rm QCD}}
\providecommand{\Iqcd}{\mathfrak I_{\rm QCD}^{(7)}}
\providecommand{\muQCD}{\mu_{\rm QCD}^{\rm int}}
\providecommand{\LambdaQCD}{\Lambda_{\rm QCD}}
\providecommand{\LambdaBH}{\Lambda_{\rm BH}^{\rm QCD}}
\providecommand{\HQCD}{\mathcal H_{\rm QCD}^{\rm BH}}
\providecommand{\HQCDfin}{H_{\rm QCD}^{\rm fin}}
\providecommand{\Hraw}{\mathcal H_{\rm raw}}
\providecommand{\HphysQCD}{\mathcal H_{\rm phys}^{\rm QCD}}
\providecommand{\Vcol}{V_C}
\providecommand{\Ps}{P_{\rm sing}}
\providecommand{\PiOct}{\Pi_8}
\providecommand{\MSbar}{\overline{\rm MS}}
\providecommand{\YM}{\mathrm{YM}}
\providecommand{\Ktri}{K_{\triangle}}
\providecommand{\KYM}{K_{\YM}^{\rm BH}}
\providecommand{\aBH}{a_{\rm BH}}
\providecommand{\epsBH}{\varepsilon_{\rm BH}}
\newcommand{\RowConstructive}{\textnormal{\textsc{Constructive}}}
\newcommand{\RowObligation}{\textnormal{\textsc{Bridge obligation}}}
\newcommand{\RowTarget}{\textnormal{\textsc{Target row}}}
\newcommand{\RowValidation}{\textnormal{\textsc{Validation row}}}
\providecommand{\slashed}[1]{\not{#1}}

% ======================================================================
% 10I. PAPER VI-B / CONTINUOUS-QFT NOTATION
% ======================================================================
% Shared symbols used by Paper VI-B and any downstream continuous-QFT or
% Yang--Mills handoff paper.  All definitions are guarded so individual
% papers may override them locally before loading this file if needed.
\providecommand{\A}{\mathfrak A}
\providecommand{\Hphys}{\mathcal H_{\rm phys}}
\providecommand{\QCD}{\mathrm{QCD}}
\providecommand{\BH}{\mathrm{BH}}
\providecommand{\loc}{\mathrm{loc}}

\providecommand{\RBH}{\mathcal R_{\rm BH}}
\providecommand{\DG}{\mathfrak D_G}
\providecommand{\PW}{\mathcal{P\!W}} % Peter--Weyl spectral-symbol subspace
\providecommand{\Toe}{\operatorname{Toe}} % Toeplitz compression/representative
\providecommand{\Hol}{\operatorname{Hol}} % Holonomy

% ======================================================================
% 11. LATE-LOADED PACKAGES (Must be at the end of the preamble)
% ======================================================================
% Load late packages only if a paper has not already loaded them, avoiding
% option clashes in paper-specific preambles.
\makeatletter
\@ifpackageloaded{hyperref}{}{\usepackage{hyperref}}
\makeatother
% --- Series-wide disambiguations (Pass 3) ---
\providecommand{\diracdelta}{\delta_{\rm D}} % Dirac delta distribution (delta reserved for hierarchy)
\providecommand{\Hg}{H_{\mathrm g}} % Hodge harmonic projector (keep H for Hubble)
\providecommand{\Gr}{G_{\mathrm r}} % Green operator (keep G for Newton)
\providecommand{\FD}{\mathrm{D}} % functional/variational derivative symbol (delta reserved)
\providecommand{\Id}{\mathbf{1}} % identity/Kronecker tensor (avoid \delta_{ij})


% BEGIN INLINE SOURCE: Paper00-ActiveExternalReferences-v2_130.tex
% Framework v2.191: current reference binding; existing definitions retained.
% Framework v2.190: current22-root registry with native clock-evolution proofs.
% Framework v2.189: exact current22-root external registry, including retained v2.188 hardening owners.
% Framework v2.187: current registry synchronization only; source release v2.186 preserved.
% Framework v2.186: authoritative current 22-paper registry.
% Framework v2.182: all current 22 logical paper AUX bindings.
% Framework v2.170: current I-B v1.43 external AUX binding.
% Framework v2.167: exact current 22-root registry after CP--BL2 integration.
% Framework v2.166: exact current 22-root compact-release AUX registry.
% Framework v2.165: authoritative current root bindings.
% Framework v2.164: exact current 22 manuscript AUX bindings.
% Framework v2.162: finite metric, functional Hilbert and scoped Ward integration.
% Framework v2.161: current 22 manuscript AUX bindings; source only.
% Framework v2.156 / registry v2.113: synchronized current 22 manuscript roots; source-only registry.
% Framework v2.143 / registry v2.111: active manuscript AUX basenames; no generated AUX or PDF included.
% ======================================================================
% Central external-reference imports, numbers only, no PDF jumping
% ======================================================================

% xr-hyper is loaded centrally so that papers may retain their existing
% \usepackage{xr-hyper} lines without redefining \externaldocument after this
% registry is installed.  If plain xr was loaded earlier, retain it.
\makeatletter
\@ifpackageloaded{xr-hyper}{}{%
  \@ifpackageloaded{xr}{}{\usepackage{xr-hyper}}%
}
\makeatother

% Make external-document registration idempotent.  The registry identifies a
% logical paper by its Paper<number><optional letter> prefix.  Prefix-qualified
% imports remain independent.  Current paper roots contain no local
% \externaldocument calls; this guard also makes accidental duplicates harmless.
\ExplSyntaxOn
\prop_new:N \g_bh_external_refs_prop
\cs_new_eq:NN \BHOriginalExternalDocument \externaldocument

\cs_new_protected:Npn \bh_external_document:nnnn #1#2#3#4
 {
  \tl_set:Nn \l_tmpa_tl {#3}
  \regex_replace_once:nnN
   { \A (?: xref/ )? (Paper[0-9]+[A-G]?) .* \Z }
   { \1 }
   \l_tmpa_tl
  \tl_set:Nx \l_tmpb_tl { #1 | \tl_use:N \l_tmpa_tl }
  \prop_if_in:NVTF \g_bh_external_refs_prop \l_tmpb_tl
   {
    \PackageInfo{BornHopf}
     {Skipping~duplicate~external~refs~for~#3}
   }
   {
    \prop_gput:NVn \g_bh_external_refs_prop \l_tmpb_tl {#3}
    \IfNoValueTF{#4}
     { \BHOriginalExternalDocument[#1][#2]{#3} }
     { \BHOriginalExternalDocument[#1][#2]{#3}[#4] }
   }
 }

\RenewDocumentCommand{\externaldocument}{O{} O{} m o}
 { \bh_external_document:nnnn {#1}{#2}{#3}{#4} }

\NewDocumentCommand{\BHExternalRefs}{m}
 {
  % Do not import the aux file of the logical paper currently being compiled,
  % even when the job and registry use different version suffixes.
  \tl_set:Nx \l_tmpa_tl { \c_sys_jobname_str }
  \regex_replace_once:nnN
   { \A (Paper[0-9]+[A-G]?) .* \Z }
   { \1 }
   \l_tmpa_tl
  \tl_set:Nn \l_tmpb_tl { #1 }
  \regex_replace_once:nnN
   { \A (Paper[0-9]+[A-G]?) .* \Z }
   { \1 }
   \l_tmpb_tl
  \str_if_eq:eeTF { \tl_use:N \l_tmpa_tl } { \tl_use:N \l_tmpb_tl }
   {
    \PackageInfo{BornHopf}
     {Skipping~external~refs~for~#1~because~this~is~the~current~document}
   }
   {
    \IfFileExists{xref/#1.aux}
     { \externaldocument[][nocite]{xref/#1}[#1.pdf] }
     {
      \IfFileExists{#1.aux}
       { \externaldocument[][nocite]{#1} }
       {
        \PackageInfo{BornHopf}
         {External~aux~file~#1.aux~not~found;~no~labels~imported}
       }
     }
   }
 }
\ExplSyntaxOff

% External references are automatic.  This file is the sole owner of both the
% bibliography and the current Paper-I-A--VII-D AUX registry (30 roots including I-C).  Missing auxiliaries
% are informational on the first pass and are populated by the source-only
% two-pass series build.  Basenames deliberately omit temporary download/copy
% suffixes such as ``(1)'' or ``-1''.
  % Paper I-A: round S3 compendium.
  \BHExternalRefs{Paper1A-Geometry-v5_181}

  % Paper I-B: complete exceptional geometry and finite spectral structures.
  \BHExternalRefs{Paper1B-ExceptionalGeometry-v1_25}
  % Paper I-D: technical companion of I-A through I-C.
  \BHExternalRefs{Paper1D-TechnicalMaterial-v1_69}
  % Paper I-C: reciprocal M8 compendium.
  \BHExternalRefs{Paper1C-M8Geometry-v1_146}

\ifdefined\BHMathOnly
  % Papers I-A and I-B import only the mathematical atlas labels.
\else
  % Paper II: five physical-mathematical main papers plus one technical proof volume.
  \BHExternalRefs{Paper2A-GeometricConstitution-v3_227}
  \BHExternalRefs{Paper2B-MomentumCarrier-v2_147}
  \BHExternalRefs{Paper2C-ReciprocalLocalization-v2_176}
  \BHExternalRefs{Paper2D-SymmetrySpectra-v2_145}
  % The former II-F export labels remain in II-E; current II-F is the technical volume.
  \BHExternalRefs{Paper2E-KinematicExports-v3_140}
  % Paper II-F: complete technical proof and calculation owner for II-A--II-E.
  \BHExternalRefs{Paper2F-TechnicalMaterial-v1_71}

  % Paper III: Dynamics and Intrinsic Geometric Constitution (six components).
  \BHExternalRefs{Paper3A-GeneratedDynamics-v1_181}
  \BHExternalRefs{Paper3B-SpectralAction-v1_192}
  \BHExternalRefs{Paper3C-ProtectedHierarchy-v1_196}
  \BHExternalRefs{Paper3D-RunningResponse-v1_167}
  \BHExternalRefs{Paper3E-GravityAction-v1_203}
  \BHExternalRefs{Paper3F-CompleteDynamics-v1_279}
  % Paper III-G: complete preserved technical proof owner.
  \BHExternalRefs{Paper3G-TechnicalMaterial-v1_57}

  % Paper IV-A: generated finite quantum mechanics.
  \BHExternalRefs{Paper4A-FiniteQuantumMechanics-v4_217}
  % Paper IV-B: complete finite quantum processes and physical response.
  \BHExternalRefs{Paper4B-QuantumResponse-v1_55}
  % Paper IV-C: extended proofs and finite technical calculations.
  \BHExternalRefs{Paper4C-TechnicalMaterial-v1_50}

  % Paper V-A: The Standard Model.
  \BHExternalRefs{Paper5A-FiniteMatterGeometry-v3_218}

  % Paper V-B: Finite Spectrum and Flavor Closure.
  \BHExternalRefs{Paper5B-TerminalSpectrum-v5_202}

  % Paper V-C: Finite Interaction Boundaries of the Standard Model.
  \BHExternalRefs{Paper5C-InteractionResponse-v2_184}
  % Paper V-D: preserved technical proof and calculation owner.
  \BHExternalRefs{Paper5D-TechnicalMaterial-v1_58}

  % Paper VI-A: Gravity and Closed Cosmology.
  \BHExternalRefs{Paper6A-GravityCosmology-v2_268}

  % Paper VI-B: Geometric Dark Sector and Cosmological Readouts.
  \BHExternalRefs{Paper6B-CosmologyReadouts-v1_215}

  % Paper VI-C: complete technical proof and comparison owner.
  \BHExternalRefs{Paper6C-TechnicalMaterial-v1_53}

  % Paper VII-A: Finite QCD Kernel and Nonperturbative Closure.
\BHExternalRefs{Paper7A-QCDKernel-v1_183}

  % Paper VII-B: Hadronic, Thermal, and Nuclear QCD Readouts.
  \BHExternalRefs{Paper7B-QCDReadouts-v1_167}

  % Paper VII-C: Continuous Quantum Field Theory.
  \BHExternalRefs{Paper7C-ContinuousQFT-v2_208}

  % Paper VII-D: complete technical material for VII-A through VII-C.
  \BHExternalRefs{Paper7D-TechnicalMaterial-v1_60}

\fi

% END INLINE SOURCE: Paper00-ActiveExternalReferences-v2_130.tex


\hypersetup{hidelinks}


\usepackage{xr-hyper}

\title{\textbf{Born--Hopf Geometry III-C: Protected Hopf--Clifford Hierarchy and Determinants}\\[0.5em]
\large Fixed-root septet, determinant characters, and boundary carriers}

\author{Peter Kugelmann\thanks{\texttt{pkugelmann@proton.me}}}
\date{Version 1.196 --- 20 November 2026 \\ Framework v3.071}
\setlength{\emergencystretch}{3em}


\begin{document}
\maketitle
\begin{abstract}
Paper III-C derives protected Hopf--Clifford hierarchy structure,
finite-carrier ancestry, determinant characters, and their compatible
boundary-response types. The seven-level rank statements concern geometric
availability and rooted carrier constructions; observable occurrence requires
the explicitly stated dynamical and source owners. Extended exceptional-sector
and finite-shell calculations are carried in Paper~III-G.
\end{abstract}
\tableofcontents
\clearpage

\section{Protected carriers and geometric prerequisites}
\label{paper3c:sec:organized-scope}
\subsection{Introduction and Paper-III series allocation}
\label{p3c:sec:series-allocation}
Paper~III-C imports Paper~II-B's derived six-type one-copy \emph{generated}
momentum seed \(\mathcal H_1\) and its associated level-one comparisons,
not a universal full-physical-endpoint multiplicity-one certificate.  The
latter still requires a generated whole-candidate GP3 comparison.  The
present protected hierarchy does not require that stronger conclusion.
It also imports
Paper~III-B's loop-derived finite spectral form class \(\Dcl\).  The current
selector contract inherited from Paper~III-B distinguishes generated
role-level constituents are distinguished from G0 token occurrences, unresolved
families are retained until an independently constituted later discriminator acts,
and a later discharge is never rewritten as upstream uniqueness.  Paper~III-B
now supplies the constructor-relative closed-loop theorem.  The present
determinant, corner, and hierarchy readouts remain decorated, regularized, or
postprocessed branches with their displayed additional data.  Paper~III-C owns the
Hopf/contact residue calculation, the \((\mathrm{REG})\)-specific smooth
determinant readout, the derived active-support chain and protected
\(1\oplus7\) endpoint, and the
finite-window Dirac threshold formulas.  The hierarchy's neutral content is
the generated dual determinant-line ratio.  Its action on an external positive
reference readout is inherited only after a generated same-owner ancestry map
is present.  The completed neutral number is labelled by the fixed-shell
coordinate \(\sigma_{\rm fin}\); monotone decrease holds exactly on the
positive one-step locus, and the dimensionful reference magnitude and units
belong to later calibration.

Physical reciprocal-role statements consume the typed Paper-II/III-A
reciprocal-role theorem
\[
 \begin{gathered}
  \text{Paper-II-C intrinsic }E_1\text{-evaluation}+\mathrm{GP2}
  \\[2pt]
  \Longrightarrow
  \\[2pt]
  \text{distinct Born/GP2-generated reciprocal roles, relation-free product,}
  \\
  \text{and Born involution}.
 \end{gathered}
\]
The certification labels \( (\mathrm{BR}_{\rm ind}) \) and
\( (\mathrm{IND}_{xp}) \) record certification of the GP2 antecedents and
are not independent assumptions at this stage.  Neither reciprocity nor III-C
changes the already fixed momentum level.

\subparagraph{Imported G0 actualization boundary.}
All physical interfaces below specialize Paper~II-A's stage-safe G0.  Each
asserted occurrence requires a contingent token of a geometrically licensed,
claim-relative invariant distinction with result-independent provenance; its
tag alone supplies no witness or independent actualization principle.  No
instance reopens frozen upstream geometry; a new role or standard enters only
through a declared downstream frontier.  On the active branch, G0 supplies
the occurring unmarked Hurwitz/Hopf token; equivariant finite functional
calculus then co-realizes the complete orbit of active chains and the common
terminal scalar constituent.  A particular body diagonal remains a marked
representative.  Thus
\((\mathrm{IR}_{O(3)})\) records a derived certification and carries no
second occurrence witness.  An actual rank-one refinement
\((\mathrm{FLAG})\), a marked member, and any projector law other than the
canonical dimensionless upper spectral filtration derived below remain
optional G0 process data.  The optional
\((\mathrm{ACT}_{\rm Cl})\) tag is retained only for such an enlarged or
physically actualized profile, not for the canonical filtration itself.  The
unmarked common-rate family \(\mathcal V(P_0)\) is different: an explicit
derived constructor transports the scalar/non-scalar and positive-decay
witnesses to the declared bare-split process sort.  Thus
\((\mathrm{NEUT}_7)\) is only its certification label and carries no second
occurrence witness.  The hierarchy-transfer label \((\mathrm{HT})\)
likewise certifies generated neutral dual-line postprocessing and supplies no
new normalization.  A physical reference scale still requires a generated
ancestry map into that owner; merely occurring in the same larger composite
does not prove the map.  The Clifford-corner tag
\((\mathrm{CORN})\)
instead records the derived constituent theorem of Paper~III-B once its
occurring ambient protected branch is present.  The
determinant-shell scheme \((\mathrm{REG})\), by contrast, remains a separately
declared analytic hypothesis: G0 types its physical use but does not derive
its regularity or uniqueness assumptions.  The generated momentum seed is already co-realized upstream and needs no
additional role witness here; a claim about a full physical endpoint beyond
that seed is governed by the separate GP3 multiplicity comparison.

\subparagraph{Standing momentum-seed boundary.}
The theorem-level \emph{generated one-copy seed} is
\[
        \mathcal H_1,\qquad m_{\rm seed}=1.
\]
Its structural role is co-realized with the upstream transfer-operation
root.  Identifying every complete physical momentum endpoint with precisely
this carrier is conditional on the unresolved GP3 role-faithful comparison;
no III-C theorem uses that stronger assertion.  Particular
tests, processes, and readouts remain G0 tokens.
Generic modules \(\mathcal H_m\) and associated spinor windows
\(\mathcal S_m\) remain auxiliary mathematical families unless a separate
role interface realizes them.  The alternative
\(\mathsf P_{\rm sli}+(\mathrm{COORD}_A)\) level-one route is an alternative
diagnostic branch and is recorded in the Paper~II-B constitutional momentum theorem.

\subparagraph{Protected Clifford and infrared boundary.}
On the larger raw completion atlas, the normalized protected Clifford
identification requires the nondegenerate protected-coordinate stratum
\[
 (\mathrm{ND}_{\rm sp}):\qquad \varrho_{\rm sp}>0.
\]
On the generated physical \(\mathsf C_D\) branch this antecedent is already
discharged by Corollary~\ref{p3c:cor:v261-physical-ndsp-discharge}.
Paper~II-C proves on this whole stratum that the normalized even generators
satisfy exactly \(Y_iY_4=C_iC_4\), so the represented protected even algebra
is the same \(\Cl(0,3)\) independently of the nonzero completion scale and
relative chiral phase.  The stronger representative \(\mathsf C_{\rm PI}\)
remains relevant only to completion-dependent numerical coordinate data and
is not used in the active Clifford chain below.

\begin{corollary}[Physical-branch discharge of protected nondegeneracy]
\label{p3c:cor:v261-physical-ndsp-discharge}
On the generated fixed-radius physical momentum branch, Paper~II-E,
Corollary~\ref{p2e:cor:generated-fixed-radius-protected-branch}, fixes the
finite-coordinate essential image to the singleton \(\{\mathsf C_D\}\) with
\[
 \boxed{\varrho_{\rm sp}(\mathsf C_D)=\frac{\sqrt3\Rmom}{4}>0.}
\]
Hence \((\mathrm{ND}_{\rm sp})\) is automatic on the physical \(\mathsf C_D\)
branch and supplies no additional occurrence, calibration, or selector there.
Every theorem below whose only extra coordinate-completion antecedent is
\((\mathrm{ND}_{\rm sp})\) therefore specializes to the generated physical
branch without a separate nondegeneracy premise.  The predicate remains
necessary on the larger raw completion atlas, where \(\varrho_{\rm sp}=0\)
continues to mark a genuine degeneracy boundary.
\end{corollary}
  Likewise,
\(\mathsf{Supp}_{2T}\) denotes the unmarked core-support output
after an actual Paper-I-B exceptional-core constructor application or a
generated bridge from occurring metric point resolution to that constructor's
domain.  Its output then needs no second witness, and labelled Hurwitz data are
calculation representatives.  Point resolution alone does not prove the
exceptional-domain bridge.  Paper~III-B proves that the only unital,
parity-even scalar projector is
\[
        P_0(a)=\langle a\rangle_0\mathbf1.
\]
The active-orbit theorem below derives this map from the grade-even projector
and the complete equivariant family of cyclic averages.  Every representative
gives the same map, so no occurring labelled Hopf line is required.  Hence
\[
        \mathbb E_{\rm IR}:\Cl(0,3)\longrightarrow\mathbb R\mathbf1
\]
is the terminal spectral projector and equals \(P_0\).  Its parity evenness,
full \(O(3)\)-equivariance, and seven-dimensional kernel \(W\) are
conclusions.  We retain \((\mathrm{IR}_{O(3)})\) only as the certification
label for these conclusions.  Write \((\mathrm{CORN})\) for Paper~III-B's
derived certification that the ambient rank-one scalar constituent is
co-realized and has normalized real trace \(1/8\); it remains distinct from
the uninserted complex-\(\Dcl\) Hilbert trace.

\subparagraph{Protected refinement boundary.}
Write \((\mathrm{FLAG})\) for an explicit nested rank-one refinement of the
seven-dimensional complement, or equivalent simple-spectrum data.  Full
\(O(7)\) symmetry does not supply \((\mathrm{FLAG})\).  Without it, only the
endpoint rank drop \(8\to1\), its rank seven, and neutral rank bookkeeping are
canonical.  No chronology of seven physical events is inferred from the
common-rate semigroup.

\subparagraph{Selector-status boundary for the active hierarchy.}
The complete orbit of discrete active chains, their invariant ranks
\(8,4,3,1\), their scalar rank-trace depth coordinate, and the common terminal
projector \(P_0\) are generated theorem-level structure.  A particular
rank-three carrier or cyclic embedding requires a marked diagonal.  An
ordered rank-one projector chronology requires the separately supplied
\(\mathrm{FLAG}\) refinement.  The active-depth operator itself has a unique
upper spectral filtration, derived below, and its pullback along the smooth
trace-defect clock needs no \(\mathrm{ACT}_{\rm Cl}\).  A marked member, an
alternative threshold law not obtained from that spectral family, or an
external physical scale/process assignment remains additional data.
The associated-character family \(\{L^q:q\in\mathbb Z\}\) is retained
upstream; only an independently constituted primitive oriented associated-line
role can later discriminate the member \(q=+1\).  The fixed-shell
unanchored regularization family \(\mathrm{REG}\) remains at this stage; the occurring Hurwitz/Hopf branch is anchored by Theorem~\ref{p3c:thm:occurring-hurwitz-reg-origin}.

\subparagraph{Forward Paper-III-D ledger.}
The flow rows citing
Theorems~\ref{p3:thm:reeb-universal},
\ref{p3:thm:eight-mode-harmonic-response},
\ref{p3:prop:one_parameter_reduction},
\ref{p3:lem:reduced-born-reciprocity},
\ref{p3:thm:geff_from_reciprocity},
\ref{p3:cor:fully_derived_one_parameter_trajectory}, and
\ref{p3:thm:exact-reciprocal-freeze-interface}
are forward-conditional Paper-III-D roadmap entries.  None is used to prove
the native contact residue, determinant anomaly, protected rank drop, or
finite-window Dirac results below.

\label{p3:part:hierarchy-reduction}

\subsection{Quadratic-parent provenance of the protected septet}
\label{p3c:sec:v189-quadratic-parent-septet}
Paper~1B, Theorem~\ref{p1b:thm:protected},
Theorem~\ref{p1b:thm:quadratic}, and
Proposition~\ref{p1b:prop:septet}~\cite{Kugelmann2026Paper1B}, give the
protected Clifford identification and the seven-dimensional Clifford-visible quotient
$\mathbb R\oplus\Lambda^2V$ of the complete rank-four quadratic parent
and its explicit oriented identification with the scalar-free septet.
Paper~III-B supplies the reciprocal realization.  After
the already frozen orientation and the owner-preserving protected Clifford
identification, this quotient is carried to the protected scalar-free septet
$W_7$ used by the rank-depth construction.  Thus the positive provenance is
\[
 \boxed{M_8\to V\otimes V\to
 (\mathbb R\oplus\Lambda^2V)\to W_7\to q\to K=8q.}
\]
The map is the existing Clifford owner map, not dimension matching.

\begin{remark}[Three sixteens remain differently typed]
\label{p3c:rem:v189-three-sixteens-firewall}
The identities
\[
 \dim(V\otimes V)=16,
 \qquad 8\,\mathscr Q''=16,
 \qquad K(q=2)=16
\]
share a parent provenance but not an owner identity.  In particular
$16_{\rm dim}$ does not select the $K=16$ action junction.  Likewise the
$(1,1)$ symmetric-traceless branch exists at the parent level but is excluded
from the generated level-one finite momentum model; this is carrier truncation,
not a dynamical suppression theorem.
\end{remark}

\begin{theorem}[Current generated local G25 boundary parent]
\label{p3c:thm:v285-main-chain-g25-boundary-parent}
The depth-six Clifford bivector generates, through its oriented grade dual and
the reciprocal first-symbol pairing, a canonical regular local geometric
boundary germ.  The resulting owner-preserving PCAD relation is
\[
 \boxed{\beta_{48,u}(1)=d_u,\qquad \beta_{48,u}(u)=n_u.}
\]
This closes the local \emph{geometric/PCAD structural port}.  It does not by
itself supply a physical action trace, flux, stiffness, Schur complement, or
boundary history; those remain same-action descendants once the corresponding
action/trace owner is supplied.  The complete regularity and collar proof is
Theorem~\ref{p3c:thm:v187-b48-boundary-parent} in the mathematical body below.
\end{theorem}

\begin{remark}[Scope of the boundary obstruction]
Theorem~\ref{p3c:thm:v158-b48-boundary-owner-nogo} correctly excludes a
physical boundary map constructed merely from the two already separated
\(C_2\) modules.  The theorem above does not reverse that argument: it
generates the local geometric/PCAD parent from the common Clifford-normal
ancestor.  Promotion from that geometric parent to a physical action trace
remains separately typed.
\end{remark}


\begin{corollary}[Dyadic Hopf trace is a structural null spectral response]
\label{p3c:cor:v191-dyadic-null-owner-typing}
The exact object
\[
 \mathcal S_H(x)=\Tr(x^{\mathcal D_H}-I)
\]
of Theorem~\ref{p3c:thm:dyadic-hopf-null-trace} is a scalar-free oriented
Hopf spectral response.  Its operator, multiplicities and dyadic eigenvalues
are derived.  Identifying a particular evaluation
$\mathcal S_H(\alpha/7)$ with the dimensionless numerical readout of the
physical operational null speed requires a separately typed physical-owner
and metrology bridge.  No recoil denominator, SI unit, or photon source is
generated by the trace theorem itself.
\end{corollary}


\subsection{First constituted nonzero carrier and its dressing port}
\label{p3c:sec:v2219-first-carrier}
\begin{proposition}[First actual nonzero active depth, not first formal Laurent index]
\label{p3c:prop:v2219-first-actual-carrier}
The previously proved active spectral resolution satisfies
\[
 \operatorname{spec}\mathcal Q_{\rm act}=\{0,4,5,7\},\quad
 P_2=P_3=0,\quad P_4=\mathbf1_{\{4\}}(\mathcal Q_{\rm act})\ne0.
\]
Here $P_j$ denotes an \emph{actual} spectral projector, not the
positive-part exponent $n_j^+$ of the Laurent comparison family.  Thus the
first nonzero constituted carrier above grade zero is uniquely grade four.
Its native energy is $E_{32}$ at the structural carrier label $K_{\rm sh}=32$.
Neither the formal ranks two/three nor the balanced structural midpoint
$K_{\rm sh}=16$ are newly occupied particle carriers.
\end{proposition}
\begin{proof}
The earlier protected-chain spectral decomposition has exactly the grades
$0,4,5,7$.  Spectral calculus gives $P_2=P_3=0$ and $P_4\ne0$.
The native scale dictionary sends carrier grade $4$ to $K_{\rm sh}=8\cdot4$.
No additional constituent is introduced by taking a scalar response midpoint.
\end{proof}
\begin{remark}[Typed coefficient-source incidence]
Selecting a gravitational Laurent dressing by this actual carrier is the
\emph{carrier-alignment rule} implemented in Paper~III-D.  It is a
cross-sector selection premise at the common-action source port; the existing
finite spectral theorem alone does not evaluate the mixed gravitational
coupling or the physical time history.
\end{remark}


\subsection{GP1--GP3 closure of the native carrier gap}
\label{p3c:sec:v165-native-carrier-gap}

\begin{theorem}[GP1--GP3 protected carrier-gap theorem]
\label{p3c:thm:v165-gp123-protected-carrier-gap}
On the occurring active branch, the GP1-finite, GP2-exact and host-faithful,
GP3-least carrier realization has
\begin{equation}
 \label{p3c:eq:v165-active-depth-support}
 \operatorname{spec}\mathcal Q_{\rm act}(u)=\{0,4,5,7\},
 \qquad
 \boxed{\mathbf1_{(5,7)}(\mathcal Q_{\rm act}(u))=0.}
\end{equation}
Equivalently, its protected shell support is
\begin{equation}
 \label{p3c:eq:v165-active-shell-support}
 \boxed{8\,\operatorname{spec}\mathcal Q_{\rm act}(u)
 =\{0,32,40,56\}.}
\end{equation}
Thus \(K=48\) is the unique point of the allowed shell lattice
\(8\mathbb Z\) strictly between \(40\) and \(56\), and
\begin{equation}
 \label{p3c:eq:v165-k48-projector-zero}
 \boxed{P_{48}^{\rm act}(u)
 =\mathbf1_{\{6\}}(\mathcal Q_{\rm act}(u))=0.}
\end{equation}
This is a hard gap in generated carrier support, not a missing profile
choice and not a declaration that a hidden carrier remains to be found.
\end{theorem}

\begin{proof}
GP1 supplies the finite functional-calculus arena in which every spectral
projector of \(\mathcal Q_{\rm act}\) is an exact finite object.  The generated
active-support theorem and the equivariant rank-two classification give the
four spectral values in \eqref{p3c:eq:v165-active-depth-support}: the only
rank-two \(C_3\)-equivariant plane is \(u^\perp\), and it excludes the scalar
line, so it cannot be inserted between the rank-three active support and the
rank-one scalar endpoint.  GP2 host faithfulness and no-insertion therefore
forbid replacing the zero spectral projector on \((5,7)\) by an externally
adjoined carrier.  GP3 compares only complete GP2-exact realizations and
cannot enlarge the least active chain by that forbidden stage.  Multiplying
the depth spectrum by the generated valuation \(K=8q\) proves
\eqref{p3c:eq:v165-active-shell-support} and the unique-lattice-point claim.
\end{proof}

\begin{corollary}[Orbit and generated-profile invariance]
\label{p3c:cor:v165-carrier-gap-profile-invariance}
The zero projector in
Theorem~\ref{p3c:thm:v165-gp123-protected-carrier-gap} is invariant under
the unmarked \(2O\)-conjugation orbit and under the generated heat-profile
rate/scale congruence of Paper~III-B.  In particular, changing
\((a,\Lambda)\) to \((\rho a,\sqrt\rho\Lambda)\) cannot create a carrier at
depth six: it leaves the exact process operator unchanged and does not alter
\(\mathcal Q_{\rm act}\) or any of its spectral projectors.
\end{corollary}

\begin{remark}[Finite carrier theorem, not a continuum theorem]
\label{p3c:rem:v165-carrier-gap-not-continuum}
Equation~\eqref{p3c:eq:v165-active-depth-support} is a theorem about the
generated finite carrier.  It supplies no refinement-limit Hamiltonian and
therefore neither asserts nor requires a continuous spectral mass gap.  The
energy realization and the precise \(E_{48}\)-centered gap are given in
Paper~V-B.
\end{remark}

\subsection{M8-rooted Clifford hierarchy and local boundary parent}
\label{p3c:sec:v187-m8-rooted-hierarchy-boundary}

\begin{theorem}[Top-down origin of the protected octet and rank-seven endpoint]
\label{p3c:thm:v187-three-eight-seven}
The vector child of the $M_8$ relative solder has rank three.  By the exact
transvection--Clifford identification of Paper~II-C it generates the protected
real Clifford algebra
\[
 \Cl(0,3),\qquad \dim_{\mathbb R}\Cl(0,3)=2^3=8.
\]
The already derived invariant support/projector reductions
\[
 \boxed{8\longrightarrow4\longrightarrow3\longrightarrow1}
\]
are therefore reductions of one $M_8$-descended Clifford owner.  Their
codimensions are
\[
 \boxed{q=0,4,5,7},
\]
and with the already fixed Clifford block size eight the hierarchy indices are
\[
 \boxed{K=8q=(0,32,40,56).}
\]
The rank-seven endpoint is the scalar-free part of the protected octet.  This
is a provenance theorem for the hierarchy owner; numerical mass calibration
or occurrence of a selected physical member is not implied.
\end{theorem}

\begin{lemma}[Automatic regularity of the canonical G25 boundary harmonic]
\label{p3c:lem:v176-g25-canonical-boundary-regularity}
Let $v_u$ be the unit position first-symbol vector obtained from the canonical
Clifford grade dual $\nu_u=-u\omega$, and let the reciprocal position parent
be the round sphere $S_x^3(R_x)$.  Up to a harmless nonzero scalar, its
canonical degree-one harmonic is
\[
 \phi_u(x)=\langle v_u,x\rangle .
\]
Its spherical gradient satisfies
\[
 \nabla_{S_x^3}\phi_u
 =v_u-\frac{\phi_u}{R_x^2}x,
 \qquad
 \|\nabla_{S_x^3}\phi_u\|^2
 =1-\frac{\phi_u^2}{R_x^2}.
\]
Hence on $\Sigma_u=\phi_u^{-1}(0)$ one has
\[
 \boxed{\|\nabla_{S_x^3}\phi_u\|=1},
\]
so $0$ is automatically a regular value and $\Sigma_u$ is the oriented great
$S^2$.  An explicit collar is
\[
 C_u(y,t)=\cos(t/R_x)y+R_x\sin(t/R_x)v_u,
 \qquad
 \phi_u(C_u(y,t))=R_x\sin(t/R_x),
\]
with $\partial_tC_u(y,0)=v_u$.  Thus regularity and local-collar existence for
the canonical generated G25 harmonic are theorem outputs, not additional
physical or analytic premises.
\end{lemma}

\begin{proof}
For the restriction of the ambient linear coordinate
$x\mapsto\langle v_u,x\rangle$ to a round sphere, the tangential projection of
$v_u$ is $v_u-(\phi_u/R_x^2)x$.  Taking its norm and using
$\|x\|^2=R_x^2$ gives the displayed identity.  The collar formulas are direct
and have nonzero normal derivative at $t=0$.
\end{proof}

\begin{theorem}[Generated local boundary parent from the depth-six bivector]
\label{p3c:thm:v187-b48-boundary-parent}
Let
\[
 B_{48}(u)=\mathbb R1\oplus\mathbb Ru
\]
be the generated depth-six carrier with unit Clifford bivector $u$, and let
$\omega=e_1e_2e_3$ be the retained orientation pseudoscalar.  Its canonical
oriented grade dual
\[
 \boxed{\nu_u=-u\omega}
\]
is a unit Clifford vector.  Through the exact protected
transvection--Clifford identity and the already generated reciprocal
first-symbol pairing it determines a unique position first-symbol covector.
By Lemma~\ref{p3c:lem:v176-g25-canonical-boundary-regularity}, its canonical
first-harmonic representative $\phi_u$ automatically has a regular zero set
and therefore generates the oriented boundary germ
\[
 \Sigma_u=\phi_u^{-1}(0)
\]
with unit conormal $n_u$, local collar and zero-horizontal-mode boundary jet
$d_u$.  PCAD then retains
\[
 \mathcal R_{48,\partial,u}
 =\{(a1+bu,\,a d_u+b n_u):a,b\in\mathbb R\},
\]
which is the graph of
\[
 \boxed{\beta_{48,u}(1)=d_u,\qquad
 \beta_{48,u}(u)=n_u.}
\]
Thus the geometric odd normalization is one.  Physical flux/stiffness
coefficients belong to the same-action Hessian response and are not free
coefficients of this parent relation.
\end{theorem}

\begin{remark}[Scope of the earlier G25 no-go]
\label{p3c:rem:v187-g25-nogo-scope}
The earlier theorem excluding a canonical element of
$\operatorname{Hom}(B_{48},J_\partial)$ from the two already separated
$C_2$ modules remains correct on that antecedent.  It now serves as a
late-descendant firewall.  The primitive construction is instead the common
Clifford-normal parent relation above.  This closes only the local structural
G25 port; the nonlocal critical $\sqrt y$ law and its refinement limit are
unchanged.
\end{remark}


\begin{corollary}[Dynamic provenance of the $B_{48}$ sign carrier]
\label{p3c:cor:v188-b48-abelianized-normalizer}
For the generated normalizer branch
\[
 2T\triangleleft2O,
 \qquad [2O,2O]=2T,
\]
the quotient that reverses the oriented diagonal is exactly the abelianization
\[
 \boxed{2O/[2O,2O]=2O/2T\cong C_2.}
\]
Consequently the $C_2$ action on
\[
 B_{48}(u)=\mathbb R1\oplus\mathbb Ru
\]
is the restriction of the real sign representation of the abelianized outer
residue of the restricted $M_8$ order-swap defect.  Among the three exceptional
supports this real sign residue is unique to $2O$.

Moreover $u\in\Cl^{(2)}(0,3)$ lies in the $M_8$-descended curvature-bivector
module.  Once the existing $2O$ orbit has generated the member $u$, the
already used dual $\nu_u=-u\omega$ is therefore the Hodge dual of a parent
curvature plane, not an independently inserted Clifford map.  The local
boundary-parent theorem then descends this vector to the conormal germ.
\end{corollary}

\begin{proof}
Paper~I-B gives the diagonal normalizer
$N_u=N_{2O}(C_6(u))\cong2D_3$ and
$N_u/C_6(u)\cong C_2$.  Projectively this is the inclusion
$D_3\cong S_3\le S_4\cong2O/\{\pm1\}$, whose oriented diagonal stabilizer is
its even subgroup $C_3=D_3\cap A_4$.  Pulling back through the double cover
therefore gives
\[
 N_u\cap2T=C_6(u).
\]
Hence the global quotient character $2O\to2O/2T\cong C_2$ restricts to the
local quotient character
\[
 N_u\longrightarrow N_u/(N_u\cap2T)
 =N_u/C_6(u)\cong C_2,
\]
which is exactly the sign action used in the construction of $B_{48}(u)$.
The curvature-bivector statement follows from
$F_{\mathcal A}=-\mathcal E\wedge\mathcal E$ together with the protected
transvection--Clifford identity and
$\Lambda^2\mathbb R^3\cong\Cl^{(2)}(0,3)$; the displayed Hodge dual is
Corollary~\ref{p2c:cor:v188-curvature-bivector-normal}.
\end{proof}

\begin{remark}[Quadratic relation closure is even under the normalizer sign]
\label{p3c:rem:v188-b48-square-even}
On the abelian quotient $C_2\cong2O/2T$, the Cartan square is trivial:
$x\mapsto x^2=1$.  Hence the first quadratic parent-relation layer is blind to
the $u\mapsto-u$ normalizer sign.  The $B_{48}$ odd line is therefore a
first-order/oriented residue and must be retained through its parent
correspondence; it cannot be reconstructed from an even quadratic invariant
alone.
\end{remark}

\begingroup
\small
\RaggedRight
\setlength{\emergencystretch}{3em}

\endgroup


\paragraph{Technical proof ownership and constituted scope.}
Every extended auxiliary derivation, explicit matrix verification,
and optional comparison moved during this reorganization is
preserved in Paper~III-G, Section~\ref{p3g:sec:hierarchy}.
The generated level-one momentum seed is inherited from Paper~II-B:
$\mathcal H_1$ with seed multiplicity one.  Neither later scalar,
action, regulator, nor source data reselect this seed.  Full physical
endpoint multiplicity-one remains conditional on its GP3 comparison.

\section{Hopf--Clifford hierarchy scale and finite seed}
\label{paper3c:sec:organized-scale}
\subsection{Intrinsic Hierarchy Constants and Reduced Spectral Factors}
\label{p3:sec:constants}

Table~\ref{tab:p3-hierarchy-flow-ledger} gives the proof ledger for the hierarchy constants and the reduced one-parameter flow.

\begin{table}[H]
\centering
\small
\begingroup
\setlength{\tabcolsep}{3pt}
\renewcommand{\arraystretch}{1.15}
\begin{tabular}{@{}
  L{3.0cm}
  L{7.2cm}
  L{\dimexpr\linewidth-10.2cm-4\tabcolsep\relax}
@{}}
\toprule
\textbf{Ledger row}
& \textbf{Claim}
& \textbf{Discharged by} \\
\midrule
Topological step
& The residue is \(\Vlogcap=(2\pi)^2\); the positive determinant character derives \(\alpha_{\rm top}=e^{4\pi^2}\).
& Definition~\ref{p3:def:alpha-two}, Theorem~\ref{p3:prop:contact-residue-normalization}, Theorem~\ref{p3:thm:Vlog} \\

Rank-one smooth value
& Assuming \(({\rm REG})\), the normalized readout is \(\delta_{\rm pr}=1/(6\sqrt2)\).
& Theorem~\ref{p3:thm:delta_rigorous}, Theorem~\ref{p3:thm:D1-primitive-smooth-uniqueness} \\

REG scheme boundary
& Alternative finite-part prescriptions are different regularization schemes.
& Corollary~\ref{p3:cor:non-D1-smooth-completions-are-insertions} \\

Seven protected ranks
& The generated active orbit has invariant carrier ranks \(8\to4\to3\to1\)
and terminal split \(1\oplus7\); a particular rank-three member remains marked,
while \((\mathrm{IR}_{O(3)})\) certifies the common endpoint.
& Theorem~\ref{p3:thm:derived-active-support}, Corollary~\ref{p3:cor:derived-infrared-projector}, Theorem~\ref{p3:thm:hopf-decomposition} \\

Protected bookkeeping clock
& The common-rate semigroup gives the unique invariant additive trace defect.
Functional calculus of the active-depth operator gives its canonical
left-continuous upper spectral filtration.  This is not an
operator-continuous rank-changing path; seven rank-one events still require
\((\mathrm{FLAG})\).
& Definition~\ref{p3:def:cumulative-rank-depth}, Theorem~\ref{p3:thm:basis-free-sevenfold-shell-clock}, Theorem~\ref{p3:thm:smooth-protected-resolution-semigroup}, Theorem~\ref{p3c:thm:common-trace-defect-spectral-filtration} \\

Universal Reeb step
& Forward Paper-III-D branch: rank valuation makes each protected rank unit contribute the same Hopf step.
& Theorem~\ref{p3:thm:reeb-universal}, Theorem~\ref{p3:thm:eight-mode-harmonic-response} \\

One-parameter reduction
& Forward Paper-III-D branch: the reduced Born flow has no hidden transverse beta system.
& Theorem~\ref{p3:prop:one_parameter_reduction}, Lemma~\ref{p3:lem:reduced-born-reciprocity} \\

Effective reciprocal coupling
& Forward Paper-III-D branch: visible reciprocal rigidity fixes \(g_{\eff}\) and the UV-normalized trajectory.
& Theorem~\ref{p3:thm:geff_from_reciprocity}, Corollary~\ref{p3:cor:fully_derived_one_parameter_trajectory} \\

Freeze interface
& Forward Paper-III-D branch: the exact reciprocal freeze equation follows from the reduced coupling.
& Theorem~\ref{p3:thm:exact-reciprocal-freeze-interface} \\
\bottomrule
\end{tabular}
\caption[Hierarchy and flow ledger]{Hierarchy and flow ledger.}
\label{tab:p3-hierarchy-flow-ledger}
\TableNotes{The table separates the proved Paper-III-C contact and rank
statements from forward Paper-III-D roadmap rows.  The integer
\(q_{\rm sp}\)-crossings are rank-trace-coordinate values, not symmetry-forced
spectral-projector jumps.  Every physical reciprocal role-space or Born-product
row inherits the premise-minimal Paper-II-C/GP2 reciprocal-role closure;
\(\mathsf P_{\rm mc}\) is only an optional G0 full-profile test realization.
The labels \( (\mathrm{BR}_{\rm ind}) \) and
\( (\mathrm{IND}_{xp}) \) are certification labels rather than added
hypotheses.}
\endgroup
\end{table}

\begin{table}[H]
\centering
\small
\begingroup
\setlength{\tabcolsep}{3pt}
\renewcommand{\arraystretch}{1.15}
\begin{tabular}{@{}
  L{2.0cm}
  L{4.0cm}
  L{3.2cm}
  L{\dimexpr\linewidth-9.2cm-6\tabcolsep\relax}
@{}}
\toprule
\textbf{Hierarchy datum}
& \textbf{Formula}
& \textbf{Value}
& \textbf{Status} \\
\midrule

\(\delta_{\rm loc}=\delta_{\rm pr}\)
& \(1/(6\sqrt{2})\)
& \(0.1178\)
& Normalized determinant-shell readout \\

\((2\pi)^2=4\pi^2\)
& Hopf contact residue
& \(39.478\)
& Derived contact residue \\

\((2\pi)^2+\delta_{\rm pr}\)
& Completed shell log-step
& \(4\pi^2+1/(6\sqrt2)\) under \(({\rm REG})\)
& Scheme-dependent completed readout \\

\(\alpha_{\rm top}\)
& \(e^{4\pi^2}\)
& \(1.397\times 10^{17}\)
& Derived neutral determinant character \\

\(\alpha_{\rm hier}(\sigma_{\rm fin})^7\)
& Seven protected non-scalar rank contributions
& Derived character value
& Scheme-labelled; dimensionful normalization remains a calibration \\

\bottomrule
\end{tabular}
\caption[Hierarchy numerical data]{Hierarchy numerical data.}
\label{tab:p3-hierarchy-data}
\TableNotes{The data table is a results roadmap.  It distinguishes the residue
theorem, the generated neutral dual-character ratios, and the
regularization-specific smooth value.}
\endgroup
\end{table}

\begin{theorem}[Unique positive determinant character]
\label{p3:thm:determinant-rank-character}
Let \(H\) be a positive isomorphism-invariant character on finite-dimensional
real multiplicity spaces, satisfying
\[
 H(0)=1,\qquad H(Q_1\oplus Q_2)=H(Q_1)H(Q_2),\qquad
 \log H(\mathbb R)=L_{\rm top}^{(1)}.
\]
Then, with \(\alpha_{\rm top}:=e^{L_{\rm top}^{(1)}}\),
\begin{equation}
\label{p3:eq:HT-rank-valuation}
 \boxed{H(Q)=
 \exp\!\bigl(\dim_{\mathbb R}(Q)L_{\rm top}^{(1)}\bigr)
 =\alpha_{\rm top}^{\,\dim_{\mathbb R}Q}.}
\end{equation}
Thus the neutral exponential rank valuation is theorem-level once the
rank-one determinant normalization is fixed; it requires no flag or physical
hierarchy postulate.
\end{theorem}

\begin{proof}
Every finite-dimensional real vector space is isomorphic to \(\mathbb R^n\).
Direct-sum multiplicativity gives
\(H(Q)=H(\mathbb R)^n\), and the rank-one normalization gives
\(H(\mathbb R)=e^{L_{\rm top}^{(1)}}\).
\end{proof}

\begin{theorem}[Abstract continuous depth character]
\label{p3c:thm:continuous-depth-character}
Let \(M:[0,7]\to\mathbb R_{>0}\) be continuous, with
\[
 M(0)=1,\qquad
 M(q+r)=M(q)M(r)
 \quad(q,r,q+r\in[0,7]),
 \qquad M(1)=\alpha>0.
\]
Then
\[
 \boxed{M(q)=\alpha^q\qquad(0\le q\le7).}
\]
This classifies any abstract continuous depth character.  No such law is
added as a premise on the physical branch: there the explicit generated
quantity \(D_\sigma(\tau)=
\alpha_\sigma^{q_{\rm sp}(\tau)}\) is constructed directly in
Theorem~\ref{p3c:thm:generated-dual-depth-hierarchy}.
\end{theorem}

\begin{proof}
Apply the continuous positive-character classification in Paper~1B,
Proposition~\ref{p1b:prop:trace-coordinate}~\cite{Kugelmann2026Paper1B},
on the additive interval \([0,7]\).  The unit fixes the constant factor
and \(M(1)=\alpha\) fixes the logarithmic slope to \(\log\alpha\),
so \(M(q)=\alpha^q\), including the endpoints by continuity.
\end{proof}

\begin{definition}[Hierarchy-transfer certification]
\label{p3:def:contact-to-scale-transfer}
The symbols \((\mathrm{HT}_{\rm phys})\) and \((\mathrm{HT})\) no longer
denote an independent occurrence or transfer hypothesis.  They are retained
only as certifications that a generated ancestry morphism has already
placed a physical readout in the dual determinant owner with its declared
weight.  The neutral dual ratio is generated in
Theorem~\ref{p3c:thm:generated-dual-depth-hierarchy}, and its owner-typed
physical inheritance is Corollary~\ref{p3c:cor:owner-typed-physical-dual-scale}.
The label neither proves ancestry nor creates the reference magnitude, chooses
physical units, or selects a weight not already fixed by the readout type.
\end{definition}

\begin{definition}[Determinant-shell regularization \(({\rm REG})\)]
\label{p3:def:determinant-shell-regularization}
On the single fixed unmarked rank-one scalar comparison shell, let a natural additive
finite prescription differ from the zeta determinant by a shell difference
\(\Delta C_{\rm fin}\), and define
\[
 \sigma_{\rm fin}:=\frac{\Delta C_{\rm fin}}{\sqrt2\,\log2},
 \qquad
 \delta_\sigma:=\frac1{6\sqrt2}+\sigma_{\rm fin}.
\]
All geometric and spectral factors in the first term---the reduced scalar
operator, absence of zero modes, one-octave interval, and Levi-to-round
covector factor---are derived below.  The surviving fixed-shell realization family is
\[
 \boxed{
 \mathfrak A_{\rm REG}
 :=\{\mathrm{REG}_{\sigma}:\sigma_{\rm fin}\in\mathbb R\},
 \qquad
 \delta_{\sigma}=\frac1{6\sqrt2}+\sigma_{\rm fin}.}
\]
The source prescription is the declared member
\[
 (\mathrm{REG}_0):\qquad \sigma_{\rm fin}=0.
\]
For statements that additionally require a decreasing determinant dual, use
the sign-restricted subfamily
\[
 \mathfrak A_{\rm REG}^{+}
 :=\left\{\mathrm{REG}_{\sigma}\in\mathfrak A_{\rm REG}:
 4\pi^2+\frac1{6\sqrt2}+\sigma_{\rm fin}>0\right\}.
\]
This is the exact domain of that ordering statement, not a GP selection of a
scheme member; \(\mathrm{REG}_0\in\mathfrak A_{\rm REG}^{+}\).
This one-parameter coordinate classifies only the observed finite difference
of this fixed scalar shell.  It is not a classification of counterterms on
arbitrary backgrounds, fields, boundaries, potentials, or profiles.  The
all-charge contact residue remains separately typed, and neither GP closure
nor the finite-Hilbert classification selects \(\sigma_{\rm fin}=0\).
\end{definition}



\subsubsection{\texorpdfstring{The Hierarchy Step $\alpha_{\rm hier}$}{The Hierarchy Step alpha hier}}

The hierarchy step has two logically separate parts.  The topological principal unit-Chern step
\[
\alpha_{\rm top}:=\exp(\Vlogcap)
\]
is the neutral determinant character of
Theorem~\ref{p3:thm:determinant-rank-character} and equals
\(\exp(4\pi^2)\).  A completed smooth step may be written
\[
\alpha(\delta_{\rm pr}):=\exp\!\big(\Vlogcap+\delta_{\rm pr}\big),
\]
where \(\delta_{\rm pr}\) denotes the finite rank-one smooth completion of the same contact
atom.  The contact theorem proves the common rank-one residue \(4\pi^2\).
The determinant theorem supplies its multiplicative rank character.  The
neutral dual ratio is generated without a physical scale premise.  On an
owner-typed positive scale role it is inherited without an additional
normalization, and it is decreasing exactly when the completed one-step
logarithm is positive; no additional \((\mathrm{HT}_{\rm phys})\) occurrence
is used.  The reduced zeta-determinant anomaly and metric conversion supply
\(\delta_{\rm loc}=1/(6\sqrt2)\); the local heat coefficient is an independent
consistency check.  Theorem~
\ref{p3:thm:D1-primitive-smooth-uniqueness} identifies this value with the
rank-one zeta-scheme completion.  A different finite shell counterterm shifts
it by \(\sigma_{\rm fin}\); this is a different regularization scheme, not
something forbidden by GP closure.


\begin{definition}[Topological and completed determinant-threshold factors]\label{p3:def:alpha-two}
We define the topological contact-residue one-step factor by
\[
\alpha_{\rm top}:=\exp(\Vlogcap).
\]
For a formal rank-one smooth-completion parameter \(\delta_{\rm pr}\), define the completed one-step
factor by
\[
\alpha(\delta_{\rm pr}) := \exp\!\big(\Vlogcap+\delta_{\rm pr}\big).
\]
Paper~III denotes the completed, scheme-labelled determinant factor by
\(\alpha_{\rm hier}(\sigma_{\rm fin})\), not by a bare \(\alpha\).  The
unconditional analytic value is the residue \(\Vlogcap=4\pi^2\), and the
source prescription \(({\rm REG}_0)\) gives
\[
\alpha_{\rm hier}(0):=\alpha(\delta_{\rm loc})=
\exp\!\left(4\pi^2+\frac{1}{6\sqrt2}\right).
\]
More generally the fixed-shell coordinate of
Definition~\ref{p3:def:determinant-shell-regularization} gives
\[
 \boxed{\alpha_{\rm hier}(\sigma_{\rm fin})
 =\exp\!\left(4\pi^2+\frac{1}{6\sqrt2}+\sigma_{\rm fin}\right).}
\]
These are neutral determinant factors.  Theorem~
\ref{p3c:thm:generated-dual-depth-hierarchy} generates their dual ratios.
Their action on a physical scale role is Corollary~
\ref{p3c:cor:owner-typed-physical-dual-scale} and additionally requires the
generated owner-ancestry morphism; the reference magnitude, unit convention,
and role-fixed weight are not selected here.
Finite/fuzzy comparison branches not identified with the smooth Hopf/contact calculation are outside the proof path, and no theorem in
Paper~III depends on them.
\end{definition}


\paragraph{The Principal Unit-Chern Hopf Atom and Primitive Associated-Line Determinant Cell}

On the relation-exact source form \(SU(2)\simeq S^3\), the Hopf fibration
\(\Sone\hookrightarrow S^3\to S^2\) is a principal \(U(1)\)-bundle.
Paper~II, Theorem~\ref{p2:thm:forced-hopf-chern-unit}, proves that any such
representative on the already derived simply connected total space has
\(|c_1|=1\); the Paper-I-A orientation convention writes the chosen
representative as \(c_1=+1\).  The \(SU(2)\) global form is an upstream
GP1--GP3 theorem.  We use the normalization
\[
\eta_0:=\eta/\Rmom,\qquad
\oint_{\Sone}\eta_0=2\pi,
\qquad
\frac{1}{2\pi}\int_{S^2}\omega=c_1(P)=1,\qquad d\eta_0=\pi^*\omega.
\]
Equivalently, the metric-scaled contact form is \(\eta=\Rmom\,\eta_0\), with
\[
\eta\wedge d\eta=\frac{2}{\Rmom}\,\mathrm{vol}_{S^3(\Rmom)}.
\]


\begin{definition}[Canonical principal unit-Chern Hopf contact-residue atom]
\label{p3:def:primitive-contact-residue}
Let \(\Delta_{b,0}\) be the positive scalar horizontal sub-Laplacian on the
standard normalized Hopf contact three-sphere, restricted off its scalar zero
mode.  More explicitly, let
\(\gamma_0(X,Y):=d\eta_0(X,JY)\) on \(H=\ker\eta_0\), choose any local
\(\gamma_0\)-orthonormal oriented frame \(X_1,X_2=JX_1\), and use the contact
density \(\eta_0\wedge d\eta_0\).  Then
\[
 \Delta_{b,0}:=d_{b}^{\,*}d_b
 =-\bigl(X_1^2+X_2^2\bigr)
\]
in the left-invariant Hopf frame, with its nonnegative Friedrichs extension
on \(L^2(S^3,\eta_0\wedge d\eta_0)\), and the inverse powers below act on
\((\ker\Delta_{b,0})^\perp\).  This fixes the Levi metric, density, sign,
and zero-mode convention used in the Heisenberg residue.  Its Heisenberg
order is two and the homogeneous dimension is \(Q=4\).
Let \(\gamma_{10}\) denote the universal scalar coefficient in
\[
 \operatorname{Res}_{H}(\Delta_{b,0}^{-2})
 =\gamma_{10}\int_{S^3}\eta_0\wedge d\eta_0.
\]
Define the residue-normalized critical contact operator
\begin{equation}
\label{p3:eq:primitive-contact-residue}
 A_{\rm H}:=\gamma_{10}^{-1}\Delta_{b,0}^{-2},
 \qquad
 L_{\rm top}^{(1)}:=\operatorname{Res}_{H}(A_{\rm H}).
\end{equation}
This is a contact spectral residue.  It is not identified with an ordinary
finite-window determinant variation \(D^{-1}\delta D\), whose residue depends
on its full Heisenberg symbol of degree \(-4\).
\end{definition}

\begin{theorem}[Canonical Heisenberg-residue normalization]
\label{p3:prop:contact-residue-normalization}
For the normalized Hopf contact form \(\eta_0\),
\begin{equation}
\label{p3:eq:canonical-hopf-residue-normalization}
 \operatorname{Res}_{H}(\Delta_{b,0}^{-2})
 =\frac18\int_{S^3}\eta_0\wedge d\eta_0
 =\frac{\pi^2}{2},
 \qquad
 \operatorname{Res}_{H}(A_{\rm H})
 =\int_{S^3}\eta_0\wedge d\eta_0.
\end{equation}
Thus the coefficient one is fixed by the canonical critical operator
\(A_{\rm H}=\gamma_{10}^{-1}\Delta_{b,0}^{-2}
=8\Delta_{b,0}^{-2}\), not by assuming a universal determinant
variation coefficient.
\end{theorem}

\begin{proof}
Use the operator, contact density, sign and zero-mode convention fixed in
Definition~\ref{p3:def:primitive-contact-residue}.  The exact scalar Hopf
contact spectrum, equivalently the Peter--Weyl horizontal spectrum recorded
upstream in Paper~I-A, gives in this normalization
\[
 \operatorname{Tr}'e^{-t\Delta_{b,0}}
 \sim \frac{\pi^2}{4t^2}
 \qquad(t\downarrow0).
\]
For an order-two Heisenberg operator in homogeneous dimension four, the
residue convention used here is
\[
 \operatorname{Res}_{H}(\Delta_{b,0}^{-2})
 =2\operatorname*{res}_{s=2}
 \operatorname{Tr}'\Delta_{b,0}^{-s}.
\]
The Mellin transform of the heat trace therefore gives
\[
 \operatorname*{res}_{s=2}\operatorname{Tr}'\Delta_{b,0}^{-s}
 =\frac{\pi^2}{4},
 \qquad
 \operatorname{Res}_{H}(\Delta_{b,0}^{-2})=\frac{\pi^2}{2}.
\]
Since
\(\int_{S^3}\eta_0\wedge d\eta_0=4\pi^2\), it follows that
\(\gamma_{10}=1/8\), and multiplication by eight gives the normalized
residue \(4\pi^2\).

Ponge's general Heisenberg-residue formula is retained as a comparison
reference only.  Its printed prefactor is not identified term-by-term with
the present normalized operator/density convention; in particular this proof
does not attribute a \(2^{-n}\) factor to Proposition~5.2 of that reference.
\end{proof}

\begin{remark}[Independent heat-trace check of the contact coefficient]
\label{p3c:rem:contact-coefficient-spectrum-check}
The exact scalar contact spectrum gives
\[
 \operatorname{Tr}(e^{-t\Delta_{b,0}})
 \sim \frac{\pi^2}{4t^2}
 \qquad(t\downarrow0),
\]
which yields
\(\operatorname{Res}_{H}(\Delta_{b,0}^{-2})=\pi^2/2\) in the same
normalization.  This independently excludes the different coefficient \(1/2\).
\end{remark}


\begin{lemma}[No hidden channel in the principal unit-Chern contact residue]
\label{p3:lem:no-hidden-hopf-channel}
For the constituted principal unit-Chern \(c_1(P)=1\) scalar Hopf contact operator, every extra
multiplicative residue factor would require an additional finite-rank tensor
factor, projector, endomorphism, support label, or profile not present in the
contact operator datum.  Hence the scalar contact residue has no hidden
channel multiplicity.
\end{lemma}

\begin{proof}
The Heisenberg residue depends on the critical symbol of the displayed scalar
operator.  No finite channel factor is present until a tensor factor is
explicitly introduced, as in Theorem~\ref{p3:thm:rank-one-hopf-localization}.
Thus the rank-one scalar operator carries exactly the residue already computed.
\end{proof}

\begin{theorem}[Matched tensor extension and uniform rank residue]
\label{p3:thm:rank-one-hopf-localization}
Let \(Q\) be a finite-dimensional real multiplicity space of rank \(r\).
Protected quotients of \(W\) are a later specialization.  Define the rank-matched contact operator
by
\[
 A_Q:=A_{\rm H}\otimes I_Q.
\]
Equivalently, require that the scalar coefficient agree with the rank-one
principal unit-Chern atom.  Then
\begin{equation}
\label{p3:eq:critical-symbol-tensor-identity}
 \sigma_{-4}(A_Q)
 =\sigma_{-4}(A_{\rm H})\otimes I_Q.
\end{equation}
Consequently
\begin{equation}
\label{p3:eq:uniform-rank-contact-residue}
 \boxed{
 \operatorname{Res}_{H}(A_Q)
 =r\operatorname{Res}_{H}(A_{\rm H})
 =4\pi^2 r.}
\end{equation}
In particular every rank-one protected atom carries the same contact
spectral residue \(4\pi^2\).  The absolute matching
\(A_Q=A_{\rm H}\otimes I_Q\) is part of this tensor extension; it is not a
consequence of \(O(Q)\)-naturality alone.
\end{theorem}

\begin{proof}
A change of orthonormal presentation of the real multiplicity space acts by
\(U\in\mathrm O(Q)\).  Naturality requires
\(U\sigma_{-4}(A_Q)U^{-1}=\sigma_{-4}(A_Q)\) for every such \(U\).  The
commutant of the full orthogonal presentation action is scalar, so the critical
symbol has the form \(a(x,\xi)I_Q\).  Symmetry alone leaves the overall scalar
function \(a\) undetermined: \(cA_{\rm H}\otimes I_Q\) is equally natural for
every constant \(c\).  The displayed rank-one matching definition fixes
\(a=\sigma_{-4}(A_{\rm H})\), which gives
\eqref{p3:eq:critical-symbol-tensor-identity}.  The Heisenberg residue includes the ordinary real finite-fibre trace, hence
\[
 \operatorname{Res}_{H}(A_Q)
 =\Tr_{\mathbb R,Q}(I_Q)\operatorname{Res}_{H}(A_{\rm H})
 =r\operatorname{Res}_{H}(A_{\rm H}).
\]
Theorem~\ref{p3:prop:contact-residue-normalization} and the Hopf integral then
give \eqref{p3:eq:uniform-rank-contact-residue}.
\end{proof}

\begin{lemma}[Determinant-line representative of the contact step]
\label{p3:lem:sharp-threshold-determinant-jump}
The principal unit-Chern logarithmic hierarchy step is the contact spectral value
\(L_{\rm top}^{(1)}=\operatorname{Res}_{H}(A_{\rm H})\).  When the reduced
shell is represented on its one-dimensional positive determinant line, the
precision character \(K_{\rm top}\) is normalized by
\begin{equation}
\label{p3:eq:det-jump-residue}
 \log\det K_{\rm top}=L_{\rm top}^{(1)}.
\end{equation}
This is the determinant-line readout of the already computed contact residue;
it is not the assertion that a generic ordinary variation
\(D^{-1}\delta D\) has residue \(4\pi^2\).
\end{lemma}

\begin{proof}
The positive determinant line is one-dimensional.  Its continuous
multiplicative character is fixed by its additive logarithmic weight, so the
unique normalized representative is
\(K_{\rm top}=\exp(L_{\rm top}^{(1)})I\), which gives
\eqref{p3:eq:det-jump-residue}.
\end{proof}

\begin{theorem}[Unit-charge Hopf contact-residue theorem: \texorpdfstring{$L_{\rm top}^{(1)}=(2\pi)^2$}{Ltop=(2pi)^2}]
\label{p3:thm:Vlog}
Within the principal unit-Chern \(c_1(P)=1\) Hopf branch,
\begin{equation}
\boxed{
 \Vlogcap
 :=L_{\rm top}^{(1)}
 =\operatorname{Res}_{H}(A_{\rm H})
 =\int_{S^3}\eta_0\wedge d\eta_0
 =(2\pi)^2
 =4\pi^2 .}
\end{equation}
\end{theorem}

\begin{proof}
Theorem~\ref{p3:prop:contact-residue-normalization} identifies the canonical
critical residue with the contact integral.  Since
\(d\eta_0=\pi^*\omega\),
\((2\pi)^{-1}[\omega]=c_1(P)=1\), and
\(\oint_{S^1}\eta_0=2\pi\), fibre integration gives
\[
 \int_{S^3}\eta_0\wedge d\eta_0
 =\left(\oint_{S^1}\eta_0\right)
  \left(\int_{S^2}\omega\right)
 =(2\pi)(2\pi)=4\pi^2.
\]
\end{proof}

\begin{corollary}[Additive Hopf-residue weight versus determinant character]
\label{p3:cor:v2122-hopf-residue-target-typing}
On the primitive rank-one Hopf atom, the generated scalar residue is
\[
 \boxed{L_H:=\Vlogcap=4\pi^2.}
\]
If a downstream generated owner is an additive one-dimensional scalar
normalization line and an owner-preserving constructor consumes this primitive
residue, its residue endomorphism is the unique scalar map
\[
 \boxed{\mathscr R_H=L_H\,\operatorname{id}.}
\]
By contrast, on the positive determinant line the corresponding multiplicative
character is
\[
 \boxed{K_{\rm top}=e^{L_H}\operatorname{id}.}
\]
These are different target-sort representations of the same generated Hopf
datum.  The corollary does not by itself assert that an arbitrary action
coefficient consumes the residue; that ancestry must be established by the
downstream owner-preserving constructor.
\end{corollary}

\begin{proof}
Theorem~\ref{p3:thm:Vlog} fixes the primitive rank-one residue to
\(L_H=4\pi^2\), while Lemma~\ref{p3:lem:no-hidden-hopf-channel} excludes any
additional primitive multiplicity.  The endomorphism algebra of a real
one-dimensional additive normalization line is scalar, so once an
owner-preserving constructor has placed \(L_H\) in that target, multiplication
by \(L_H\) is the unique such endomorphism.  Lemma~\ref{p3:lem:sharp-threshold-determinant-jump}
separately proves that exponentiation occurs only when the same additive
residue is represented as a multiplicative positive determinant character.
\end{proof}

\begin{remark}[What is and is not being identified]
The number \(\Vlogcap\) is not a Riemannian volume and is not obtained by
setting RG time equal to Reeb time.  It is the canonical critical contact
spectral residue of the principal unit-Chern Hopf atom.  The Reeb vector is Killing on
\(S^3\), while scale transformations are generated on the cone \(C(S^3)\)
by \(r\partial_r\).  Hopf geometry fixes the residue value; any physical
threshold location or protected rank event paired with that value must be
supplied independently.
\end{remark}

\begin{corollary}[Completed principal unit-Chern contact step]
\label{p3:cor:completed-primitive-rg-jump}
If an independently supplied protected rank-one event is paired with the
principal unit-Chern Hopf contact atom, then its matched topological logarithmic step is
\(\Vlogcap=4\pi^2\).  The contact geometry fixes this residue weight; the
existence, location, and ordering of protected events are separate data.
\end{corollary}

\begin{proof}
Theorems~\ref{p3:thm:rank-one-hopf-localization} and~\ref{p3:thm:Vlog}
compute the rank-one critical residue.  No Dirac-to-protected intertwiner is
used or inferred.
\end{proof}


\begin{remark}[No auxiliary protected finite branch]
\leavevmode\par\noindent
No auxiliary UV-subtracted finite/fuzzy comparison quantity built from the
same Hopf operator is identified with the smooth Hopf/contact quantity
\[
  \Vlogcap+\delta_{\rm pr}.
\]
No core derivation in Paper~III depends on that comparison branch.
\end{remark}

\paragraph{The Exponential Form}

\begin{theorem}[Neutral contact factor and scheme-labelled smooth completion]
\label{p3:thm:alpha_step}
The topological principal unit-Chern hierarchy factor is
\begin{equation}
\boxed{
\alpha_{\rm top}
=\exp(\Vlogcap)
=\exp\left((2\pi)^2\right)
=\exp(4\pi^2).
}
\end{equation}
For a formal rank-one smooth-completion parameter \(\delta_{\rm pr}\), the corresponding completed factor is
\begin{equation}
\boxed{
\alpha(\delta_{\rm pr})
=
\exp\!\big(\Vlogcap+\delta_{\rm pr}\big)
=
\exp\!\big(4\pi^2+\delta_{\rm pr}\big).
}
\end{equation}
This separates the topological threshold from the smooth slot.  The boxed
family is formal in \(\delta_{\rm pr}\).  Under the source prescription
\(({\rm REG}_0)\), the rank-one smooth calculation below gives
\(\delta_{\rm pr}=1/(6\sqrt2)\).  A nonzero \(\sigma_{\rm fin}\) adds to this
logarithmic completion.  No physical scale interpretation is used in either
calculation.
\end{theorem}

\begin{proof}
Theorem~\ref{p3:thm:Vlog} gives the canonical Heisenberg residue
\[
  \Vlogcap=(2\pi)^2=4\pi^2.
\]
Exponentiation yields \(\alpha_{\rm top}=\exp(4\pi^2)\).  By definition,
the finite completion variable adds to this logarithmic step, so
\[
  \alpha(\delta_{\rm pr})
  =\exp(4\pi^2+\delta_{\rm pr}).
\]
The normalized determinant-shell calculation gives
\(\delta_{\rm loc}=1/(6\sqrt2)\), and
Theorem~\ref{p3:thm:D1-primitive-smooth-uniqueness} identifies it with
\(\delta_{\rm pr}\) under \(({\rm REG})\).
\end{proof}


\section{Smooth Hopf and contact hierarchy correction}
\label{paper3c:sec:organized-contact}
\subsection{\texorpdfstring{The smooth Hopf/contact correction \(\delta\)}{The smooth Hopf/contact correction}}
\label{p3:sec:delta}

The quantity denoted \(\delta_{\rm loc}\) below is computed as the normalized
first-order momentum-shell readout of the reduced formal half-determinant anomaly.
The local scalar heat coefficient is retained as an independent consistency
check, not as the determinant correction itself.  The section then proves that
the \(({\rm REG})\) prescription gives
\[
  \delta_{\rm pr}=\delta_{\rm loc}=\frac{1}{6\sqrt2}
\]
inside the displayed determinant-shell scheme.

The supplementary \(\mathsf{Supp}_{2T}\) compatibility now includes the
exact equality \(\delta_{\rm loc}=\omega_{2O}\), the \(2O\) tangent-frame
reconstruction of the Levi metric factor, and
\(\delta_{\rm loc}^{-2}=\operatorname{rank}(E_6)h(E_6)=72\).  These are
recorded in Paper~III-G, Appendix~\ref{p3c:app:e6-compatibility}; they are not inputs to
the native heat-kernel or determinant calculation below and do not select a
regularization origin without the stated owner map.

\begin{remark}[The Levi-to-round factor is computed, not chosen]
\label{p3:rem:levi-round-factor}
On the horizontal distribution \(H=\ker\eta_0\) of the unit round \(S^3\),
with \(\eta_0=\tfrac12(d\psi+\cos\theta\,d\varphi)\) normalized so that
\(\int_{S^3}\eta_0\wedge d\eta_0=4\pi^2\), the round metric and the Levi form
satisfy
\[
        g(e_i,e_i)=\tfrac14,
        \qquad
        \lvert L(e_i,e_i)\rvert=\tfrac12,
        \qquad\text{hence}\qquad
        L=2\,g\ \text{ on } H .
\]
A Levi-orthonormal vector therefore has round length \(1/\sqrt2\), and the
first-order covector conversion from Levi to round normalization is
\(\sqrt{g/L}=1/\sqrt2\).  The factor is thus a ratio between two forms on the
same constituted contact distribution; the underlying \(2\) traces to
\(\int_{S^2}\omega=2\pi\) against a base area \(\pi\).  No convention enters.
\end{remark}

The derivation separates the global determinant data from the local
heat-kernel coefficient.  The determinant-shell value uses three inputs, all
internal to the present geometric package:

\begin{enumerate}
\item the exact fixed-charge trace, whose expanded constant gives
      \(\zeta_{\rm red}(0)=1/3\) and hence the declared half-determinant anomaly
      \(\Delta\log Z_{\det}=\log2/6\) under Levi-to-round rescaling;
\item the intrinsic quadratic-momentum interval
      \(|\Delta\log p^2|=\log2\); and
\item the exact first-order covector conversion between the Levi and round
      horizontal frames, which contributes \(1/\sqrt2\).
\end{enumerate}
They determine
\(\delta_{\rm loc}=(1/\sqrt2)(\Delta\log Z_{\det}/|\Delta\log p^2|)\).
The unexpanded global constant \(C_{\rm trace}=-1/6\) and the local
Laplace-type scalar-curvature coefficient \(1/6\) are independent consistency
checks; neither is substituted for \(\zeta_{\rm red}(0)\) in the determinant
scaling formula.  Identifying \(\delta_{\rm pr}\) with this normalized
determinant-shell value is the \((\mathrm{REG})\) prescription.

\begin{remark}[Geometric input imported from Paper~I-A]
The derivation below is specialized to the standard regular Hopf contact structure on the round CR three-sphere.  In particular, it uses the global sub-Riemannian geometry of $\SU(2)$, the standard pseudohermitian invariants, and the round harmonic package recorded in Paper~I-A.  No observational benchmark is imported.  The coefficient $1/6$ is the universal local scalar coefficient in the heat-kernel density of a scalar Laplace-type operator; the reduced Hopf scalar trace supplies a separate global check in the primitive oriented associated-line channel.
\end{remark}


\begin{remark}[Contact \(Q=4\) is an analytic scaling datum]
\label{p3:rem:contact-q4-not-spacetime-dimension}
The Paper-I-A Heisenberg tangent cone has homogeneous dimension \(Q=4\).  In the
present section this licenses the contact heat-kernel and horizontal Sobolev
scaling used by the principal unit-Chern Hopf comparison.  It is not used as an independent
proof of Lorentzian spacetime dimension, nor as an extra continuum-field input.
The Lorentzian scalar-cover statements remain the separate Born/Hamiltonian
closures proved elsewhere in the series.
\end{remark}

\begin{definition}[Associated-character atlas and primitive oriented comparison line]
\label{p3c:def:primitive-oriented-associated-line}
For the upstream Hopf principal bundle retain the complete associated-character
atlas
\[
 \mathfrak A_{\rm Hopf}^{\rm assoc}
 :=\{L^q:q\in\mathbb Z\},
 \qquad c_1(L^q)=q.
\]
The principal-bundle statement \(c_1(P)=+1\) fixes the frozen orientation of
the Hopf generator; it does not by itself select a physical associated charge.
For the present determinant comparison, independently constitute the analytic
role \(L_{\rm prim}^{+}\): the primitive associated scalar line compatible
with the frozen positive Hopf orientation.  The role is specified without
naming its answer: its signature is
\(q\in\mathbb Z\), primitive, and positive relative to the frozen orientation.
Hence
\[
 \boxed{L_{\rm prim}^{+}=L^{+1}.}
\]
Chronologically the full atlas remains unresolved upstream and the member
\(q=+1\) is \emph{derived later} only for this analytic comparison role.  No
matter-charge or gauged-\(U(1)\) assignment follows.
\end{definition}

\begin{remark}[No hidden Hopf-basic zero mode in the closure]
\label{p3:rem:no-hidden-basic-dolbeault-zero-mode}
Paper~I-A's calculation gives one basic scalar ground state and no basic
\((0,1)\) harmonic zero mode in the regular \emph{untwisted} \(q=0\)
Hopf-basic Dolbeault complex.  That result excludes an unaccounted transverse
zero-mode multiplet only in this untwisted auxiliary complex; it does not
apply automatically to the \(L^q\)-twisted sectors.  The comparison channel used below is instead the derived primitive oriented
associated-line sector \(L_{\rm prim}^{+}=L^{+1}\).  Its absence of a
zero mode follows independently from the explicit positive ladder in
Proposition~\ref{p3:prop:reduced-hopf-scalar-channel-uniqueness}.  Any further
zero mode, ghost, modulus, or base-complex deformation would be additional
spectral data outside the specified \(({\rm REG})\) calculation.
\end{remark}


\begin{proposition}[Primitive oriented associated-line Hopf scalar channel]
\label{p3:prop:reduced-hopf-scalar-channel-uniqueness}
Let \(L_{\rm prim}^{+}=L^{+1}\) be the analytic primitive oriented associated
line derived in Definition~\ref{p3c:def:primitive-oriented-associated-line},
and let
\(\mathcal V_{+1}\cong L^2(S^2,L_{\rm prim}^{+})\) be its scalar section
space.  On this space define the dimensionless horizontal operator
\[
       \Delta_h^{\rm red}:=\Rmom^2\Delta_H|_{\mathcal V_{+1}}.
\]
Within the following data this is the rank-one scalar comparison channel:
\begin{enumerate}[label=(\roman*),leftmargin=2.1em]
\item the standard regular Hopf contact structure on the round three-sphere, with the
Hopf/contact normalization imported from Papers~I-A--II;
\item the independently constituted primitive oriented associated-line role,
whose signature derives \(q=+1\) relative to the frozen positive orientation,
with no doubling by the oppositely oriented \(q=-1\) channel;
\item the scalar horizontal sub-Laplacian inherited from the round Hopf metric,
with no extra endomorphism potential, labelled projector, or sector-dependent profile;
\end{enumerate}
its spectrum is
\begin{equation}
\label{p3:eq:reduced-hopf-unique-ladder}
        \lambda_j=4j^2-2,
        \qquad
        m_j=2j,
        \qquad j=1,2,\ldots .
\end{equation}
In particular, the ladder has no zero mode.  Any different scalar ladder changes the
Chern charge, the metric/operator normalization, or the channel multiplicity and
therefore changes one of the displayed data.
\end{proposition}

\begin{proof}
The Paper-I-A fixed-charge spectral resolution gives
\[
 \Rmom^2\lambda_{k,q}=k(k+2)-q^2,
 \qquad |q|\le k,\quad k\equiv q\pmod2,
\]
with multiplicity \(k+1\) for each fixed pair \((k,q)\).  Definition~\ref{p3c:def:primitive-oriented-associated-line} retains the full
associated-character atlas upstream and independently constitutes the primitive
positive associated-line comparison role.  Its primitive positive signature
therefore gives \(q=+1\), so its scalar sections lie in that fixed comparison
sector.  The parity condition then makes
\(k=2j-1\), \(j\ge1\), and direct substitution gives
\[
 k(k+2)-1=(2j-1)(2j+1)-1=4j^2-2,
 \qquad k+1=2j.
\]
The lowest admissible \(k\) is one, so zero cannot occur.  Reversing the bundle
orientation replaces \(q=+1\) by the isospectral conjugate \(q=-1\), but
summing both would double the multiplicity and would constitute an additional
orientation channel.  Thus the stated primitive oriented associated-line channel and its
ladder follow from the later-derived \(q=+1\) comparison role and the fixed-charge Hopf resolution.
\end{proof}

\begin{theorem}[Primitive oriented associated-line Hopf scalar spectral trace]
\label{p3:prop:hopf-spectral-trace}
\leavevmode\par\noindent
In the primitive oriented associated-line scalar channel used for
contact-threshold normalization, the operator \(\Delta_h^{\rm red}\) has
eigenvalues
\[
  4j^2-2,
  \qquad
  j=1,2,\dots,
\]
with multiplicity \(2j\).  Its exact trace is
\begin{equation}
\Tr e^{-t\Delta_h^{\rm red}}
=
2e^{2t}\sum_{j=1}^{\infty}j e^{-4tj^2}.
\end{equation}
This is the later-derived primitive \(q=+1\) comparison-line trace used below, not the full horizontal
sub-Laplacian trace over all \(\SU(2)\) weights and not a sum over both
orientations.
\end{theorem}
\begin{proof}
Proposition~\ref{p3:prop:reduced-hopf-scalar-channel-uniqueness} identifies
the primitive oriented \(q=+1\) comparison line.  In that channel,
Paper~1B, Theorem~\ref{p1b:thm:dyadic-trace}~\cite{Kugelmann2026Paper1B},
gives exactly the displayed eigenvalues, multiplicities and heat trace
with the present unit-sphere normalization.
\end{proof}

\begin{theorem}[Dyadic Hopf null trace and square-root asymptotics]
\label{p3c:thm:dyadic-hopf-null-trace}
On the primitive oriented associated-line channel of
Proposition~\ref{p3:prop:reduced-hopf-scalar-channel-uniqueness}, define by
spectral calculus
\begin{equation}
 \mathcal J_H:=\frac12\bigl(\Delta_h^{\rm red}+2I\bigr)^{1/2},
 \qquad
 \mathcal D_H:=2^{-\mathcal J_H}.
\end{equation}
Then \(\mathcal J_H\) has eigenvalue \(j\) and \(\mathcal D_H\) has
eigenvalue \(2^{-j}\), both with multiplicity \(2j\), for
\(j=1,2,\ldots\).  In particular, \(\mathcal D_H\) is positive and trace
class, with
\begin{equation}
 \label{p3c:eq:dyadic-hopf-trace-class}
 \Tr\mathcal D_H=\sum_{j\ge1}2j\,2^{-j}=4.
\end{equation}
For \(x>0\), put
\(x^{\mathcal D_H}:=\exp((\log x)\mathcal D_H)\).  The difference
\(x^{\mathcal D_H}-I\) is trace class and its exact null trace is
\begin{equation}
 \label{p3c:eq:exact-dyadic-null-trace}
 \mathcal S_H(x)
 :=\Tr\bigl(x^{\mathcal D_H}-I\bigr)
 =\sum_{j=1}^{\infty}2j\bigl(x^{2^{-j}}-1\bigr).
\end{equation}
As \(x\to\infty\),
\begin{equation}
 \label{p3c:eq:dyadic-null-trace-asymptotic}
 \mathcal S_H(x)
 =2x^{1/2}+4x^{1/4}
   +O\!\bigl(x^{1/8}\log x\bigr)
 =2x^{1/2}+O(x^{1/4}).
\end{equation}
Thus, for example,
\[
 \mathcal S_H(\alpha_\sigma/7)
 =\frac{2}{\sqrt7}\sqrt{\alpha_\sigma}
   +O(\alpha_\sigma^{1/4})
 \qquad(\alpha_\sigma\to\infty).
\]
There is no constant \(C\) for which
\(\mathcal S_H(x)=C\sqrt{x}\) holds for all \(x>0\).  The square-root law
is an asymptotic spectral statement, not an exact identity.
\end{theorem}

\begin{proof}
With the primitive channel fixed by
Proposition~\ref{p3:prop:reduced-hopf-scalar-channel-uniqueness}, the
operators \(\mathcal J_H,\mathcal D_H\) coincide with those in
Paper~1B, Theorem~\ref{p1b:thm:dyadic-trace}~\cite{Kugelmann2026Paper1B}.
That theorem supplies the spectra, trace class property, exact trace and
asymptotic expansion in
Eqs.~\eqref{p3c:eq:dyadic-hopf-trace-class}--\eqref{p3c:eq:dyadic-null-trace-asymptotic}.
Substituting \(x=\alpha_\sigma/7\) gives the displayed hierarchy asymptotic.
Finally \(\mathcal S_H(1)=0\) would force \(C=0\) in an exact law
\(C\sqrt x\), whereas the same trace theorem gives
\(\mathcal S_H'(1)=\Tr\mathcal D_H=4\), a contradiction.
\end{proof}

\begin{theorem}[Reduced Hopf heat constant and zeta value]
\label{p3:thm:hopf-trace-constant}
The reduced scalar-channel small-\(t\) expansion is
\begin{equation}
 \Tr e^{-t\Delta_h^{\rm red}}
 =\frac{1}{4t}+\frac13+O(t).
\end{equation}
Equivalently, if the leading factor is kept unexpanded,
\[
 \Tr e^{-t\Delta_h^{\rm red}}
 =\frac{e^{2t}}{4t}-\frac16+O(t).
\]
Thus the renormalized constant in the latter convention is
\(C_{\rm trace}=-1/6\), while the spectral zeta value is
\begin{equation}
\label{p3:eq:zeta-red-zero}
 \boxed{\zeta_{\rm red}(0)=\frac13.}
\end{equation}
\end{theorem}

\begin{proof}
Paper~1B, Theorem~\ref{p1b:thm:dyadic-trace}~\cite{Kugelmann2026Paper1B},
gives the reduced heat expansion and \(\zeta_{\rm red}(0)=1/3\) for the
primitive channel identified above.  Since
\(e^{2t}/(4t)=1/(4t)+1/2+O(t)\), keeping that leading factor
unexpanded changes the displayed remainder constant to
\(1/3-1/2=-1/6\).  This is the stated \(C_{\rm trace}\) convention.
\end{proof}

\begin{lemma}[Local scalar heat coefficient in Levi-normalized horizontal units]
\label{p3:lem:local-sixth-heat}
Let \(L\to S^2\) be the degree-one Hermitian line bundle selected by the
primitive oriented associated line, let \(\nabla^L\) be its descended Hopf connection,
and let \(g_\gamma\) be the Levi-normalized transverse metric.  Define the
elliptic base operator
\[
 L_\gamma:=(\nabla^L)^*\nabla^L
\]
on \(L^2(S^2,L,d\mu_\gamma)\), with its self-adjoint Friedrichs realization
and form domain \(H^1(S^2,L)\) (equivalently operator domain \(H^2\) on the
closed smooth base).  No additional endomorphism potential is inserted.  Its
diagonal heat-kernel density has the local bundle-endomorphism form
\begin{equation}
K_{L_\gamma}(t;x,x)
\sim
(4\pi t)^{-1}
\left(I_L+\frac{t}{6}\operatorname{Scal}_\gamma(x)I_L+O(t^2)\right)
\end{equation}
in the fixed-charge scalar line-bundle channel.  Equivalently, the coefficient multiplying the normalized scalar curvature density is exactly
\begin{equation}
\label{p3:eq:local-sixth-lemma}
\frac16 .
\end{equation}
This is a two-dimensional elliptic base expansion.  It is not the heat
expansion of the full contact sub-Laplacian on \(S^3\), whose homogeneous
dimension is four and whose leading order is \(t^{-2}\).
\end{lemma}
\begin{proof}
This is the standard local Seeley--DeWitt coefficient for a minimally coupled
connection Laplacian on the closed two-dimensional base.  In local
Levi-normal coordinates, an operator
\(L=-(g^{ab}\nabla_a\nabla_b+E)\) has
\[
K_L(t;x,x)\sim (4\pi t)^{-q/2}\left(1+t\left(\frac{\operatorname{Scal}(x)}{6}+E(x)\right)+O(t^2)\right).
\]
Here \(E=0\) and the base dimension is \(q=2\).  Therefore the coefficient of
the normalized scalar curvature density is \(1/6\)
\cite{berline1992heat,ChamseddineConnes1997}.  The connection curvature first
enters later coefficients.  This local base calculation is distinct from both
the global fixed-charge trace constant of
Theorem~\ref{p3:thm:hopf-trace-constant} and the all-charge contact residue of
Theorem~\ref{p3:prop:contact-residue-normalization}.  Combining those distinct
objects is precisely the declared \(({\rm REG})\) prescription.
\end{proof}

\begin{proposition}[Levi-normalized scalar coefficient]
\label{p3:prop:delta-local-sixth}
\leavevmode\par\noindent
For the standard Hopf/CR three-sphere, the Levi-normalized local scalar
amplitude is
\begin{equation}
\label{p3:eq:delta-local-sixth}
c_{\mathrm{loc}}^{(\gamma)}=\frac16.
\end{equation}
This statement concerns the local heat-kernel density coefficient only; it is
not an identification with the global constant \(C_{\mathrm{trace}}\).
\end{proposition}

\begin{proof}
By definition, \(c_{\mathrm{loc}}^{(\gamma)}\) is the coefficient multiplying the normalized scalar curvature density in the Levi-normalized horizontal scalar channel.  Lemma~\ref{p3:lem:local-sixth-heat} computes that coefficient and gives \(c_{\mathrm{loc}}^{(\gamma)}=1/6\).
\end{proof}

\begin{lemma}[Paper-I-A dimensionless Levi metric versus the round Hopf metric]
\label{p3:lem:levi-vs-round}
Use the Paper~I-A dimensionless Hopf pair \((g_E,\eta_0)\), so that
\[
  \eta_0\wedge d\eta_0=2\,\vol_{g_E},
  \qquad \eta_0(X_H)=1,
  \qquad H=\ker\eta_0.
\]
Let \(J\) be the standard transverse complex structure and define
\[
  \gamma_0(X,Y):=d\eta_0(X,JY),\qquad X,Y\in H.
\]
Then
\begin{equation}
\label{p3:eq:levi-vs-round}
  \boxed{\gamma_0=2g_E|_H}.
\end{equation}
Thus, if \((e_1,e_2,X_H)\) is a positively oriented \(g_E\)-orthonormal
frame with \(e_2=Je_1\), the \(\gamma_0\)-orthonormal frame is
\begin{equation}
\label{p3:eq:levi-frame-conversion}
  f_1=e_1/\sqrt2,\qquad f_2=e_2/\sqrt2.
\end{equation}
The equivalent physical-radius identities are
\(g_{\Rmom}=\Rmom^2g_E\), \(\eta_{\Rmom}=\Rmom\eta_0\), and
\(d\eta_{\Rmom}(X,JY)=2g_{\Rmom}(X,Y)/\Rmom\); the two normalizations are not
mixed termwise.
\end{lemma}

\begin{proof}
This is exactly the Hopf scale dictionary fixed in Paper~I-A.  The displayed
frame conversion is the inverse square-root conversion for
\(\gamma_0=2g_E|_H\).
\end{proof}


\begin{lemma}[Levi-to-dimensionless-round determinant anomaly and momentum conversion]
\label{p3:lem:levi-round-linear-amplitude}
On the primitive oriented \(q=+1\) descended comparison-line channel, set explicitly
\[
        L_E:=\Delta_h^{\rm red}.
\]
Let \(L_\gamma\) be the same descended connection Laplacian represented in
the Levi-normalized transverse metric.  Since
\(\gamma_0=2g_E|_H\),
\[
 L_\gamma=\frac12L_E,
 \qquad
 p_\gamma^2=\frac12p_E^2,
 \qquad
 |p|_\gamma=\frac1{\sqrt2}|p|_E.
\]
With \(\zeta_{\rm red}(0)=1/3\), the positive formal half-determinant amplitude
satisfies
\begin{equation}
\label{p3:eq:integrated-det-anomaly}
 \Delta\ln Z_{\det}
 :=\ln Z_\gamma-\ln Z_E
 =\frac{\ln2}{6}.
\end{equation}
The corresponding logarithmic quadratic-momentum interval is
\begin{equation}
\label{p3:eq:quadratic-momentum-interval}
 |\Delta\ln p^2|=\ln2.
\end{equation}
\end{lemma}

\begin{proof}
For a positive operator \(L\),
\(\log\det_\zeta(cL)=\log\det_\zeta L+\zeta_L(0)\log c\).
Taking \(c=1/2\) and using
Theorem~\ref{p3:thm:hopf-trace-constant} gives
\[
 \log\det_\zeta L_\gamma-\log\det_\zeta L_E
 =-\frac13\ln2.
\]
By the declared formal amplitude \(Z_{1/2}:=(\det_\zeta L)^{-1/2}\),
\(\Delta\ln Z_{\det}=\ln2/6\).  The metric relation gives the displayed
quadratic- and first-order momentum scalings.
\end{proof}

\begin{proposition}[Determinant amplitude and real/complex measure dictionary]
\label{p3c:prop:v2211-determinant-amplitude-measure}
For the positive primitive complex line operator used here, define the formal
amplitude $Z_a=(\det_\zeta L)^{-a}$, $a>0$. Its Levi/round scale anomaly
and normalized shell readout are
\[
 \boxed{\Delta\log Z_a=\frac{a\log2}{3},\qquad
 \delta_a=\frac{a}{3\sqrt2}.}
\]
The present REG/compiler signature uses $a=1/2$ and therefore
$\delta_{1/2}=1/(6\sqrt2)$. This is a specified positive formal
half-determinant readout. It is not automatically the ordinary Gaussian
partition function of a section of the degree-one complex line $L^{+1}$.
For a finite positive Hermitian matrix, one complex Gaussian coordinate has
two real components and contributes determinant exponent $-1$, whereas one
real Gaussian coordinate contributes $-1/2$. A genuine complex-section
Gaussian consequently uses $a=1$ in the same complex spectral trace and has
twice the half-amplitude anomaly. The degree-one line has no real line form
whose complexification preserves its Chern class. Any later physical Gaussian
use must transport its actual real/complex measure, field content and
normalization into this dictionary. The analytic comparison role itself
constitutes no charged matter field.
\end{proposition}
\begin{proof}
The previously computed zeta value is $\zeta_L(0)=1/3$. The determinant
scaling law for $L_\gamma=L_E/2$ therefore gives
$\Delta\log Z_a=a\log2/3$; dividing by $\log2$ and applying the same
covector conversion gives $\delta_a$. For a real positive eigenvalue
$\ell$, the real integral is proportional to $\ell^{-1/2}$, while the
integral over its real and imaginary components is proportional to
$\ell^{-1}$. Products give the finite-dimensional determinant powers;
the printed zeta prescription preserves their declared multiplicities.
A complexification of a real rank-one bundle has Chern class of order at
most two; $H^2(S^2;\mathbb Z)$ is torsion-free, so it cannot have degree
one. Thus a real Gaussian interpretation cannot silently identify that
complex channel with one real coordinate. None of this changes the value
under the explicit half-determinant prescription.
\end{proof}

\begin{proposition}[Primitive determinant and complete density jets]
\label{p3c:prop:v2217-determinant-density-jets}
Hold the primitive momentum operator, domain, connection, measure and
regularization fixed while varying an independently supplied spacetime
curvature $R$. Its zeta determinant has zero first and second $R$ derivative.
If an additional incidence instead supplies $L(R)=(R/C)L(C)$, $R\ge C>0$,
and $Z=(\det_\zeta L)^{-a}$, $a>0$, its reference-centered total action is
$W=-\log[Z(R)/Z(C)]=(a/3)\log(R/C)$.
For an actual positive density factor $N(R)=\nu(R)V_4(R)$, the scalar
coefficient $f=W/N$ has the complete second jet
\[
 f''=\frac{W''}{N}-\frac{2W'N'}{N^2}
             +W\left(\frac{2(N')^2}{N^3}-\frac{N''}{N^2}\right).
\]
On the conditional homothetic chart $V_4=v_0/R^2$, $\nu$ constant, this gives
\[
 f=\frac{aR^2}{3\nu v_0}\log(R/C),\qquad
 f''=\frac{a}{3\nu v_0}[2\log(R/C)+3]>0.
\]
\end{proposition}
\begin{proof}
The stipulated fixed operator has constant spectrum and determinant. For
the other family, $\zeta_{L(R)}(s)=(R/C)^{-s}\zeta_{L(C)}(s)$.
Differentiation at zero, using $\zeta_{L(C)}(0)=1/3$, gives $W$.
Twice applying the quotient rule gives the general jet; substituting
$N=\nu v_0/R^2$ gives the last two formulas.
\end{proof}
The homothety is an explicitly supplied comparison incidence, not an
identification of the two curvature owners. Any moving localizer, band,
normalization or reference term contributes its own derivative. A total
logarithmic determinant therefore cannot be assigned a negative local
curvature susceptibility before this density and complete pullback test.
Nor does a constant-coefficient $1-C^2/R^2$ comparison reproduce the full
native factor $1+n_7(R)$ and its terminal jets.


\begin{proposition}[Normalized determinant-shell readout]
\label{p3:prop:delta-reduction}
The rank-one local smooth correction is obtained in two stages.  First,
normalize the determinant anomaly per unit logarithmic interval in \(p^2\).
Then express the result in first-order round-momentum units:
\begin{equation}
\label{p3:eq:delta-reduction}
 \boxed{
 \delta_{\rm loc}
 =\frac1{\sqrt2}
  \frac{\Delta\ln Z_{\det}}{|\Delta\ln p^2|}.}
\end{equation}
Thus \(\delta_{\rm loc}\) is not the integrated determinant correction itself;
it is its normalized first-order momentum/contact readout.
\end{proposition}

\begin{proof}
The reduced spectral action is a function of the quadratic variable \(p^2\), so
its intrinsic logarithmic shell coordinate is \(\ln p^2\).  Dividing the
integrated determinant anomaly by the exact interval gives the response density
per unit quadratic log scale.  The exported contact observable is first order,
so the Levi-to-round conversion contributes the covector factor \(1/\sqrt2\).
No additional scale or profile is introduced by the displayed \((\mathrm{REG})\) readout.  This smooth determinant statement is not inferred from the finite-Hilbert classification of \(\Dcl\).
\end{proof}

\begin{proposition}[Numerical value of the normalized momentum-shell correction]
\label{p3:prop:delta-sqrt2}
Using \(\Delta\ln Z_{\det}=\ln2/6\) and
\(|\Delta\ln p^2|=\ln2\),
\begin{equation}
\label{p3:eq:delta-from-levi-round}
 \delta_{\rm loc}
 =\frac1{\sqrt2}\frac{\ln2/6}{\ln2}
 =\frac{1}{6\sqrt2}.
\end{equation}
\end{proposition}

\begin{proof}
Substitute
Eqs.~\eqref{p3:eq:integrated-det-anomaly} and
\eqref{p3:eq:quadratic-momentum-interval} into
Proposition~\ref{p3:prop:delta-reduction}.
\end{proof}

\begin{remark}[Distinct trace and normalized-response sixths]
The two appearances of \(1/6\) have different origins.  The constant in
Theorem~\ref{p3:thm:hopf-trace-constant} is the global renormalized constant
of the reduced Hopf scalar-channel heat trace.  The normalized determinant
response instead follows from \(\zeta_{\rm red}(0)=1/3\), which gives the
integrated half-determinant anomaly \(\Delta\ln Z_{\det}=\ln2/6\); division by the
intrinsic quadratic-momentum interval \(\Delta\ln p^2=\ln2\) yields the
response density \(1/6\).  The global trace constant is therefore an
independent spectral check and is not identified with the normalized
momentum-shell coefficient.
\end{remark}

\begin{theorem}[Normalized determinant-shell value]
\label{p3:thm:delta_rigorous}
The rank-one smooth Hopf/contact correction is
\begin{equation}
 \boxed{\delta_{\rm loc}=\frac{1}{6\sqrt2}.}
\end{equation}
The integrated formal half-determinant anomaly is separately
\begin{equation}
 \boxed{\Delta\ln Z_{\det}=\frac{\ln2}{6}.}
\end{equation}
The two values are related by the normalized first-order momentum readout
\[
 \delta_{\rm loc}
 =\frac1{\sqrt2}
  \frac{\Delta\ln Z_{\det}}{|\Delta\ln p^2|}.
\]
\end{theorem}

\begin{proof}
Theorem~\ref{p3:thm:hopf-trace-constant} gives
\(\zeta_{\rm red}(0)=1/3\).  Lemma~\ref{p3:lem:levi-round-linear-amplitude}
therefore gives the determinant anomaly \(\ln2/6\) over the exact interval
\(|\Delta\ln p^2|=\ln2\).  Proposition~\ref{p3:prop:delta-reduction} then gives
\[
 \delta_{\rm loc}
 =\frac1{\sqrt2}\frac{\ln2/6}{\ln2}
 =\frac1{6\sqrt2}.
\]
\end{proof}
\begin{remark}[Exact octahedral shell match and owner criterion]
\label{p3:rem:2O-shell-match-owner-firewall}
Paper~I-B independently supplies the canonical normalizer-shell atom
\[
 \omega_{2O}
 =\frac{\operatorname{Re}(S_{2O})}{|S_{2O}|}
 =\frac1{6\sqrt2},
\]
so its scalar value agrees exactly with the analytically derived
\(\delta_{\rm loc}\).  Paper~I-B also reconstructs the same Levi metric from
the \(2O\) shell frame and identifies the defining and adjoint shell-walk
eigenvalues \(1/\sqrt2\) and \(1/3\).
Paper~III-G, Appendix~\ref{p3c:app:e6-compatibility} records the full finite/smooth
compatibility.

Equality of real values alone does not identify the finite shell owner with
the determinant regularization owner.  That companion appendix therefore constructs the
required operator-origin word before evaluation.  On the generated
octahedral normalizer/interface branch that word closes
\(\mathrm{REG}_0\); on branches without that anchor the owner firewall and
full affine family remain.
\end{remark}

\paragraph{Route-B closure: uniqueness of the \texorpdfstring{$(\mathrm{REG})$}{REG}-local rank-one smooth completion}
\label{p3:sssec:D1-local-delta-closure}

The previous theorem computes the normalized determinant-shell value.  We now
record its uniqueness after the \(({\rm REG})\) clauses are fixed.  Independent
finite counterterms, channel projectors, or alternative finite-part
prescriptions define other schemes.  This uniqueness claim is internal to
\((\mathrm{REG})\), not to the finite-Hilbert classification of \(\Dcl\).

\begin{definition}[Rank-one normalized determinant-shell germ]
\label{p3:def:first-order-contact-shell-germ}
The rank-one smooth correction is the local response density of one reduced
quadratic spectral shell,
\[
 \frac{\Delta\ln Z_{\det}}{|\Delta\ln p^2|},
\]
exported through the first-order horizontal momentum symbol.  Hence the
quadratic determinant normalization and the first-order covector conversion are
separate steps of the same readout.
\end{definition}

\begin{definition}[\((\mathrm{REG})\)-local rank-one smooth completion]
\label{p3:def:D1-local-primitive-smooth-completion}
Let \(Q\) be an arbitrary finite-dimensional real multiplicity/comparison
space.  A finite smooth threshold functional \(\mathcal S_{\rm sm}(Q)\) is a
\((\mathrm{REG})\)-local rank-one smooth completion when it satisfies:
\begin{enumerate}[label=(\roman*),leftmargin=2.1em]
\item orthogonal-presentation naturality: \(\mathcal S_{\rm sm}\) is invariant
under every orthogonal isomorphism of \(Q\), and no additional structure,
grade projector, endomorphism, channel label, or sector-dependent profile is
attached to \(Q\);
\item direct-sum additivity on independent real multiplicity spaces;
\item determinant-shell compatibility with the same reduced scalar operator,
so the rank-one integrated anomaly is \(\Delta\ln Z_{\det}=\ln2/6\);
\item normalization by the same intrinsic interval
\(|\Delta\ln p^2|=\ln2\), followed by the first-order Levi-to-round covector
factor \(1/\sqrt2\).
\end{enumerate}
The rank-one scalar \(\delta_{\rm pr}\) is the rank-one value.  Protected
quotients of \(W\) are a later specialization and are not required to define
this smooth scheme.
\end{definition}

\begin{theorem}[Uniqueness under the determinant-shell prescription \(({\rm REG})\)]
\label{p3:thm:D1-primitive-smooth-uniqueness}
Every smooth completion satisfying
Definition~\ref{p3:def:D1-local-primitive-smooth-completion} is the rank
valuation
\begin{equation}
\label{p3:eq:D1-smooth-completion-rank-valuation}
 \mathcal S_{\rm sm}(Q)
 =\dim_{\mathbb R}(Q)\,\delta_{\rm loc}.
\end{equation}
For a rank-one atom,
\begin{equation}
\label{p3:eq:delta-pr-equals-loc-D1}
 \boxed{\delta_{\rm pr}=\delta_{\rm loc}=\frac1{6\sqrt2}.}
\end{equation}
\end{theorem}

\begin{proof}
Orthogonal-presentation naturality and the absence of additional structure make
\(\mathcal S_{\rm sm}\) depend only on the real rank.  Direct-sum additivity therefore gives
\(\mathcal S_{\rm sm}(Q)=\dim_{\mathbb R}(Q)s_1\) for one rank-one constant
\(s_1\).  Determinant-shell compatibility and the fixed normalization give
\[
 s_1
 =\frac1{\sqrt2}\frac{\ln2/6}{\ln2}
 =\frac1{6\sqrt2}
 =\delta_{\rm loc}
\]
by Theorem~\ref{p3:thm:delta_rigorous}.  This proves both displayed formulas.
\end{proof}

\begin{corollary}[Alternative smooth completions are different schemes]
\label{p3:cor:non-D1-smooth-completions-are-insertions}
The value
\[
        \delta_{\rm pr}=\delta_{\rm loc}=\frac{1}{6\sqrt2}
\]
is unique as a rank-one smooth completion inside \(({\rm REG})\).
Any completion assigning a different rank-one value must abandon at least one
of the stated \(({\rm REG})\) compatibility, naturality, additivity, or
normalization clauses. Examples of additional or changed data include a finite
counterterm, determinant exponent or determinant/field measure, non-scalar
grade or channel projector, endomorphism potential, sector-dependent cutoff
profile, transverse zero-mode or boundary selector, and local finite-part
prescription. This is not an exhaustive classification of modifications.  Such a completion is
a different regularization or dynamical datum, not a contradiction of the
computed \(({\rm REG})\) value.
\end{corollary}

\begin{proof}
Theorem~\ref{p3:thm:D1-primitive-smooth-uniqueness} fixes the rank-one
constant, under orthogonal-presentation naturality, threshold additivity,
determinant-shell compatibility, and round-shell normalization, to
\[
 \frac{1}{6\sqrt2}.
\]
A different value must therefore violate at least one of these clauses.
Violating orthogonal-presentation naturality introduces extra presentation data, a labelled projector, potential,
channel, or profile.  Violating determinant-shell compatibility introduces a
different finite part, scale interval, or counterterm.  Violating the scalar
reduction introduces a transverse mode or boundary selector, while violating
round-shell normalization changes the constituted Hopf/contact normalization.
These are examples, not a necessary list of all modifications; changing the
determinant power or integration measure may also alter the rank-one finite
part. What the uniqueness theorem proves is the disjunction that at least
one defining clause must change, not an exhaustive taxonomy of how it changes.
Such changes may be studied in a larger spectral class but are not fixed by
the present calculation.
\end{proof}

\subparagraph{Framework export: rank-one smooth correction.}
Theorem~\ref{p3:thm:D1-primitive-smooth-uniqueness} exports the branch-internal identity
\[
        \ddef:=\delta_{\rm pr}=\delta_{\rm loc}=\frac1{6\sqrt2}.
\]
The finite denominator interpretation \(72=8\cdot9=|\Phi(E_6)|\) is then a Paper~III-G compatibility/origin check for the analytic value.

\begin{remark}[Scope of the closure]
\label{p3:rem:delta-closure-scope}
Theorem~\ref{p3:thm:D1-primitive-smooth-uniqueness} closes the identification
under \(({\rm REG})\).  Theories with other counterterms, projectors, profiles,
or finite-part prescriptions are different spectral data.  Later uses of the
defect scalar inherit the \(({\rm REG})\) tag.
\end{remark}



\section{Topological origin of the protected septet}
\label{paper3c:sec:organized-rank}
\subsection{Protected Clifford carrier and generated rank-seven endpoint}
\label{p3:sec:hierarchy-origin}

On a same-owner exceptional branch admitted by Theorem~\ref{p3c:thm:v236-exceptional-active-fork}, the active response
constructed below generates the basis-independent rank-seven non-scalar
complement from the protected Clifford operator sector.  By
Theorem~\ref{p3:thm:determinant-rank-character}, rank additivity
turns this rank into the factor \(\alpha_{\rm top}^7\); under
\(({\rm REG})\), the selected fixed-shell scheme gives the completed factor
\(\alpha(\delta_{\rm loc})^7\).  The exponent is not counted by arbitrary
basis vectors or by irreducible \(\SO(3)\) multiplets.  The endpoint rank
deficiency is seven.  For any active spectral projector its cumulative
deficiency \(q\) is integer-valued, but the realized values need not visit
every integer.  The active response derives the realized depths
\(q=0,4,5,7\).  Adopting \((\mathrm{FLAG})\) supplies the finer bookkeeping
list \(q=0,1,\ldots,7\), but no such flag is needed for the four active
supports.  The endpoint formula \(8\cdot7=56\) is
imported from Paper~I-B, Theorem~\ref{p1:thm:septet-octet-resolution-module}.

\begin{theorem}[Protected Clifford channel module on \((\mathrm{ND}_{\rm sp})\)]
\label{p3:thm:clifford-dim}
Assume \((\mathrm{ND}_{\rm sp})\).  Paper~III-B then supplies the represented
real protected Clifford algebra and regular module
\[
        A_{\rm prot}\cong A_{\rm ev}\cong\Cl(0,3),
        \qquad
        \mathcal S_{{\rm prot},\mathbb R}\cong\Cl(0,3)_{\rm reg}.
\]
The canonical grade decomposition is
\[
\Cl(0,3)=\Cl^{(0)}\oplus\Cl^{(1)}\oplus\Cl^{(2)}\oplus\Cl^{(3)},
\qquad
\dim_{\mathbb R}(\Cl^{(0)},\Cl^{(1)},\Cl^{(2)},\Cl^{(3)})=(1,3,3,1).
\]
This theorem is an import of the typed module bridge; it is not derived from
the numerical equality
\(\dim_{\mathbb R}\End_{\mathbb C}(S_+)=8\).
\end{theorem}

\begin{proof}
Under \((\mathrm{ND}_{\rm sp})\), Theorem~\ref{p3:thm:regular-module-two-eights}
of Paper~III-B identifies the protected diagonal-weight module with the real
regular \(\Cl(0,3)\)-module, and its preceding protected-algebra theorem
identifies the represented even operator algebra with \(\Cl(0,3)\).  The
standard basis
\[
1,\ e_1,e_2,e_3,\ e_{12},e_{13},e_{23},\ e_{123}
\]
has the displayed grade dimensions.  The lowest Dirac half-spinor window is
therefore only a separate analytic two-complex-dimensional window; its
endomorphism algebra is not used as the physical Clifford-algebra bridge.
\end{proof}



\begin{definition}[Finite mobile-Clifford role]
\label{p3c:def:v237-mobile-clifford-role}
Assume \((\mathrm{ND}_{\rm sp})\) and write
\[
 A:=A_{\rm prot}\cong\Cl(0,3),\qquad
 A^+:=A^{\rm even}\cong\mathbb H,\qquad
 G:=U_1(A^+)\cong\operatorname{Sp}(1).
\]
A \emph{finite mobile-Clifford role}, denoted \(\mathsf{MobCl}\), consists of
an occurring finite same-owner subgroup \(\Gamma<G\), one local tetrahedral
Clifford projector frame \(\mathcal F\subset\operatorname{End}(\operatorname{Im}\mathbb H)\),
and the conditions that \(\Gamma\) contains the orientation-preserving
binary-tetrahedral core associated with \(\mathcal F\) and that some element
of \(\Gamma\) moves \(\mathcal F\) to a distinct equivalent frame.  The role
contains no \(2I\), \(2O\), \(2T\), or \(C_6\) label as input.
\end{definition}

\begin{theorem}[Generated Clifford--exceptional resolution system]
\label{p3c:thm:v237-generated-clifford-exceptional-resolution}
On \((\mathrm{ND}_{\rm sp})+\mathsf{MobCl}\), Paper~I-B's
Theorem~\ref{p1:thm:v237-mobile-tetrahedral-completion} forces
\[
 \boxed{\Gamma\cong2I.}
\]
Its generated same-owner consequence packet consists of five tetrahedral cores
\(H_i\cong2T\), ten pairwise intersections
\(C_{ij}=H_i\cap H_j\cong C_6\), and five ambient normalizers
\[
 O_i=N_G(H_i)\cong2O.
\]
No core or axis member is selected.

Each \(C_{ij}\) canonically determines an unoriented quaternionic axis line
\(L_{ij}=\mathbb Ru_{ij}\subset\operatorname{Im}\mathbb H\) and hence the
sign-independent rank-one Clifford projector
\[
 P_{ij}:=u_{ij}\otimes u_{ij}.
\]
For each fixed core \(H_i\), its four incident projectors form a tetrahedral
tight frame,
\[
 \boxed{\sum_{j\ne i}P_{ij}=\frac43 I_{\operatorname{Im}\mathbb H}.}
\]
Globally the ten projectors satisfy
\[
 \boxed{\sum_{i<j}P_{ij}=\frac{10}{3}I_{\operatorname{Im}\mathbb H}.}
\]
Consequently the complete unmarked exceptional incidence generates the
Clifford metric on the imaginary-quaternion sector; a particular oriented
axis remains a member token only for a downstream claim that consumes its
orientation.
\end{theorem}

\begin{proof}
The group statements are Corollary~\ref{p1:cor:v237-five-core-ten-axis-packet}.
The local and global projector identities are now owned upstream by
Paper~I-B, Theorem~\ref{p1:thm:v285-ten-axis-quadrupole-exact-sequence}; this
paper imports them as generated resolution data on the same
\(\mathsf{MobCl}\) owner.
\end{proof}

\begin{corollary}[Trace-free quadratic resolution module of the ten-axis system]
\label{p3c:cor:v285-ten-axis-quadrupole-resolution}
On the branch of Theorem~\ref{p3c:thm:v237-generated-clifford-exceptional-resolution},
let \(D\) and \(\Phi\) be the incidence and quadrupole maps of Paper~I-B.
Then
\[
 \boxed{\mathbb R^{10}/D(\mathbb R^5)
 \cong\operatorname{Sym}^2_0(\operatorname{Im}\mathbb H).}
\]
This is the canonical five-dimensional trace-free quadratic resolution-data
module of the generated ten-axis system.  The local incidence coordinates are
the affine \(\mathcal H_0\oplus\mathcal H_2\) profiles of
Paper~I-C, Theorem~\ref{p1b:thm:v285-tetrahedral-affine-incidence}.  Globally
the exceptional siblings densely generate \(SU(2)\), while on a common
\(C_6\) axis their finite normalizers densely generate
\(\operatorname{Pin}(2)\), by Paper~I-B,
Theorems~\ref{p1:thm:2O-2I-generate-SU2} and
\ref{p1:thm:v285-shared-axis-pin2}.
\end{corollary}

\begin{remark}[Resolution-data firewall]
The quotient, affine incidence profile and continuous completions above are
mathematical descendants of the generated Clifford owner.  They do not by
themselves constitute a perturbation covariance, probability law, physical
time evolution, gauge-field occurrence, or preferred axis.
\end{remark}

\begin{corollary}[Generated global and local active-family ancestry]
\label{p3c:cor:v237-mobile-active-family-ancestry}
On the branch of Theorem~\ref{p3c:thm:v237-generated-clifford-exceptional-resolution},
each generated core \(H_i\) carries its four generated \(C_6\) axes and hence
one complete local active Clifford family of
Theorem~\ref{p3:thm:derived-active-support}.  The five local families form one
unmarked global incidence family with ten shared axes.  The rank/depth/endpoint
packet
\[
 8\to4\to3\to1,\qquad q=0,4,5,7,\qquad P_{\rm term}=P_0
\]
is therefore generated without independent \(2I\), \(2O\), \(2T\), \(C_6\),
body-diagonal, or unmarked active-orbit occurrences.  Selecting one oriented
member remains a G0 token only where an orientation-sensitive claim consumes it.
\end{corollary}

\begin{remark}[Constructor firewall]
\label{p3c:rem:v237-no-abstract-module-promotion}
The physical bridge used above is the generated same-owner map
\(C_{ij}\mapsto L_{ij}\mapsto P_{ij}\).  No identification is made merely
from abstract equalities such as a four-dimensional incidence module being
isomorphic, after recentering, to \(A^+\), or an eight-dimensional incidence
kernel having the same restricted representation type as \(\Cl(0,3)\).
\end{remark}

\begin{theorem}[Conditional entries and generated mobile completion of the active finite family]
\label{p3c:thm:v236-exceptional-active-fork}
Assume \((\mathrm{ND}_{\rm sp})\), so that
\(A_{\rm prot}\cong\Cl(0,3)\) is already constituted.  The finite active
family has three constitutionally distinct entries.
\begin{enumerate}[label=(\roman*),leftmargin=2.2em]
\item \emph{Local-only entry.}  If \(H\cong2T\subset U(A^{\rm even})\)
occurs, its four order-six subgroups generate the unmarked local tetrahedral
axis family.  For a generator \(c\) put \(u_C=(c-c^{-1})/\sqrt3\); changing
\(c\) to \(c^{-1}\) reverses \(u_C\) but leaves the line projector and subgroup
Reynolds average unchanged.
\item \emph{Interface entry.}  If \(O\cong2O\) occurs, then
\([O,O]=H\cong2T\) generates the local core and branch (i), while the
normalizer/interface parent remains active.
\item \emph{Global mobile entry.}  If the result-independent role
\(\mathsf{MobCl}\) of Definition~\ref{p3c:def:v237-mobile-clifford-role}
occurs, Theorem~\ref{p3c:thm:v237-generated-clifford-exceptional-resolution}
forces \(2I\) and generates the full five-core/five-normalizer/ten-axis
incidence family.  No exceptional group label or axis family is an independent
occurrence on this branch.
\end{enumerate}
Every generated local member has the common rank/depth/terminal packet
\[
 8\to4\to3\to1,\qquad q=0,4,5,7,\qquad P_{\rm term}=P_0.
\]
These common data may descend from an unresolved entry family; entry-specific
incidence data remain branch dependent.
\end{theorem}


\begin{theorem}[Generated equivariant orbit of active Clifford chains]
\label{p3:thm:derived-active-support}
Assume \((\mathrm{ND}_{\rm sp})\) and let \(H\cong2T\subset
U(A^{\rm even})\) be either an occurring local-only core from
Theorem~\ref{p3c:thm:v236-exceptional-active-fork}(i), the commutator core of
branch (ii), or one generated core of the mobile branch (iii).  Let
\(\mathcal O_{\rm diag}\) denote a coherent set of calculation representatives
for the four unoriented \(C_6\)-axis lines of this local core under the
canonical normalizer \(2O=N_{\operatorname{Sp}(1)}(H)\), represented by unit
bivectors
\[
 u\in A^{(2)},\qquad u^*=-u,\qquad u^2=-1,
 \qquad A:=A_{\rm prot}\cong\Cl(0,3).
\]
Replacing \(u\) by \(-u\) leaves the projector and subgroup Reynolds data below
unchanged.  No axis-line member is selected by the unmarked family; an
orientation-sensitive downstream claim must supply the corresponding member
token.
With
\[
 \langle a,b\rangle_{\rm Cl}:=\tau_8(L_{a^*b}),
 \qquad
 \Pi_8:=I_A,
 \qquad
 \Pi_4:=\frac{I_A+\gamma}{2},
\]
define, for every \(u\in\mathcal O_{\rm diag}\),
\[
 \Pi_u(a):=\langle u,a\rangle_{\rm Cl}u,
 \qquad
 \Pi_3(u):=\Pi_4-\Pi_u,
 \qquad
 E_3(u):=\operatorname{Ran}\Pi_3(u).
\]
Then \(\Pi_4A=A^{\rm even}\cong\mathbb H\) has rank four,
\(E_3(u)=\mathbb R\mathbf1\oplus u^\perp\) has rank three, and the
Hurwitz unit
\[
 c_u:=\cos\frac{\pi}{3}+u\sin\frac{\pi}{3},
 \qquad c_u^6=1,
\]
generates a cyclic group \(C_6(u):=\langle c_u\rangle\) whose conjugation
representation preserves \(E_3(u)\); its kernel is the central subgroup
\(\{1,c_u^3\}=\{\pm1\}\), so the induced action factors faithfully through
\(C_6(u)/\{\pm1\}\cong C_3\).  The carrier \(E_3(u)\) and the acting
group \(C_6(u)\) are distinct typed objects.  Put
\[
 \mathcal A_6(u):=\frac16\sum_{j=0}^{5}\operatorname{Ad}_{c_u^j},
 \qquad
 \Pi_1(u):=\Pi_3(u)\mathcal A_6(u)\Pi_3(u).
\]
For every orbit member,
\begin{equation}
\label{p3:eq:derived-active-projector-chain}
\begin{gathered}
 \Pi_1(u)\le\Pi_3(u)\le\Pi_4\le\Pi_8,\\
 \operatorname{rank}_{\mathbb R}
   (\Pi_8,\Pi_4,\Pi_3(u),\Pi_1(u))=(8,4,3,1),\\
 \boxed{\Pi_1(u)=P_0.}
\end{gathered}
\end{equation}
Thus the four carrier ranges are
\[
 A,\qquad A^{\rm even},\qquad E_3(u),\qquad\mathbb R\mathbf1,
\]
while the additional finite-support/action data are, respectively, the
protected Clifford structure, the Hurwitz units \(2T\), the action
\(C_6(u)\curvearrowright E_3(u)\), and the trivial action on the scalar line.

Set
\[
 Q_0(u):=\Pi_8-\Pi_4,\quad
 Q_4(u):=\Pi_4-\Pi_3(u),\quad
 Q_5(u):=\Pi_3(u)-P_0,\quad
 Q_7(u):=P_0
\]
and define
\begin{equation}
\label{p3:eq:active-depth-operator}
 \mathcal Q_{\rm act}(u)
 :=0Q_0(u)+4Q_4(u)+5Q_5(u)+7Q_7(u).
\end{equation}
Its upper spectral projectors are
\begin{equation}
\label{p3:eq:active-upper-spectral-projectors}
 \mathbf1_{[q,7]}(\mathcal Q_{\rm act}(u))
 =\Pi_{8-q}(u)
 \quad\text{for}\quad
 (q,8-q)=(0,8),(4,4),(5,3),(7,1),
\end{equation}
where \(\Pi_8(u):=\Pi_8\), \(\Pi_4(u):=\Pi_4\), and
\(\Pi_1(u):=P_0\).  Hence every member has the same invariant depth and clock
data
\begin{equation}
\label{p3:eq:derived-active-depth-clock}
 \boxed{q=0,4,5,7,
 \qquad K_{\rm sh}=8q=0,32,40,56.}
\end{equation}

The complete family
\[
  \mathscr C_{\rm act}
 :=\left(
 \mathcal O_{\rm diag};
 \Pi_8,\Pi_4,
 \{\Pi_3(u),P_0,C_6(u)\curvearrowright E_3(u),
   \mathcal Q_{\rm act}(u)\}_{u\in\mathcal O_{\rm diag}}
 \right)
\]
is equivariant under the unmarked conjugation action of the canonical
normalizer \(2O=N_{\operatorname{Sp}(1)}(2T)\) and is the
generated output.  A particular \(\Pi_3(u)\), cyclic embedding,
\(\mathcal Q_{\rm act}(u)\), or memberwise response requires a marked
diagonal token.  No mark is needed for the orbit family, ranks, spectra,
depths, or common endpoint.  No \((\mathrm{ACT}_{\rm Cl})\),
\((\mathrm{FLAG})\), particle assignment, or empirical mass value is used.
\end{theorem}

\begin{proof}
The grade involution is orthogonal for the normalized regular inner product.
Its positive eigenspace is \(A^{(0)}\oplus A^{(2)}\), hence has rank four and
is the quaternionic even subalgebra.  Paper~I-B identifies the units of the
occurring Hurwitz order with \(2T\).  For every unit bivector \(u\),
\(\Pi_u\) is a rank-one orthogonal subprojector of \(\Pi_4\); therefore
\(\Pi_3(u)\) is an orthogonal rank-three projector containing the scalar line.

Conjugation by \(c_u\) fixes \(1\) and \(u\) and rotates
\(u^\perp\subset\operatorname{Im}\mathbb H\) through \(2\pi/3\).
Consequently \(E_3(u)=\mathbb R1\oplus u^\perp\) is invariant, and the finite
average \(\mathcal A_6(u)\) is the orthogonal projector onto
\(\operatorname{span}_{\mathbb R}\{1,u\}\) in \(A^{\rm even}\).  Hence
\[
 \Pi_3(u)\mathcal A_6(u)\Pi_3(u)(a)
 =\langle a\rangle_0\mathbf1=P_0(a).
\]
This proves the projector identities and ranks.  The four differences
\(Q_0(u),\allowbreak Q_4(u),\allowbreak Q_5(u),\allowbreak Q_7(u)\) are mutually orthogonal and sum to the
identity, so finite functional calculus gives the displayed spectral
projectors and depths.

If \(g\in2O=N_{\operatorname{Sp}(1)}(2T)\) and \(u'=gug^{-1}\), invariance of the Clifford
inner product gives
\[
 \Pi_3(u')=\operatorname{Ad}_g\Pi_3(u)\operatorname{Ad}_{g^{-1}},
 \qquad
 \mathcal Q_{\rm act}(u')
 =\operatorname{Ad}_g\mathcal Q_{\rm act}(u)\operatorname{Ad}_{g^{-1}},
\]
and \(C_6(u')=gC_6(u)g^{-1}\).  The scalar projector \(P_0\) is fixed.
Thus the construction is a result-independent equivariant operation on the
complete upstream orbit: it neither chooses nor deletes a representative.
The orbit family, its action maps, and its finite functional-calculus closure
are exactly its nonredundant essential image.  The rank, spectrum, depth, and
endpoint are constant on the orbit and therefore descend to the unmarked
quotient.  The clock relation \(K_{\rm sh}=8q\) is the Paper-I-A
septet--octet resolution valuation.
\end{proof}

\begin{theorem}[Equivariant rank-two classification and the active depth-six obstruction]
\label{p3c:thm:active-depth-six-obstruction}
Fix \(u\in\mathcal O_{\rm diag}\) and let the effective
\(C_3=C_6(u)/\{\pm1\}\) action on
\(E_3(u)=\mathbb R\mathbf1\oplus u^\perp\) be the one generated in
Theorem~\ref{p3:thm:derived-active-support}.  The complete list of orthogonal
projectors \(P\le\Pi_3(u)\) commuting with this action is
\begin{equation}
 \label{p3c:eq:c3-equivariant-subprojectors}
 \boxed{0,\quad P_0,\quad
 P_{u^\perp}:=\Pi_3(u)-P_0,\quad\Pi_3(u).}
\end{equation}
The rank-two member \(P_{u^\perp}\) is the rotating plane itself and obeys
\(P_{u^\perp}P_0=0\).  Consequently
\begin{equation}
 \label{p3c:eq:no-nested-rank-two-carrier}
 \boxed{P_0\le P\le\Pi_3(u),\ [P,C_3]=0
 \quad\Longrightarrow\quad P\in\{P_0,\Pi_3(u)\}.}
\end{equation}
There is therefore no rank-two equivariant carrier in the nested active
upper-support chain.  Equivalently,
\begin{equation}
 \label{p3c:eq:active-depth-six-absent}
 \operatorname{spec}\mathcal Q_{\rm act}(u)=\{0,4,5,7\},
 \qquad
 \boxed{\mathbf1_{\{6\}}(\mathcal Q_{\rm act}(u))=0.}
\end{equation}
This is the active depth-six obstruction.  It does not deny the ambient
rank-two plane \(u^\perp\); it proves that this plane is not the missing
rank-two stage between the scalar endpoint and \(E_3(u)\).
\end{theorem}

\begin{proof}
Conjugation by \(c_u\) fixes \(\mathbb R\mathbf1\) and restricts to a
rotation through \(2\pi/3\) on the real plane \(u^\perp\).  That plane has no
nonzero proper real invariant line.  Moreover it has no fixed vector, so an
equivariant orthogonal projector has no off-diagonal block between the
trivial line and the rotating plane.  Its two diagonal blocks are therefore
independently either zero or the identity, giving precisely
\eqref{p3c:eq:c3-equivariant-subprojectors}.  Requiring the range to contain
\(P_0\) leaves only \(P_0\) and \(\Pi_3(u)\), proving
\eqref{p3c:eq:no-nested-rank-two-carrier}.  The spectral resolution
\eqref{p3:eq:active-depth-operator} has eigenvalues \(0,4,5,7\), so finite
functional calculus gives \eqref{p3c:eq:active-depth-six-absent}.
\end{proof}

\begin{corollary}[Generated co-realization of the active-chain orbit]
\label{p3c:cor:active-projector-corealization}
On the occurring active branch,
\[
 \boxed{
 \operatorname{CoReal}
 \bigl(\mathscr C_{\rm act},P_0\mid P_{\rm active}\bigr).}
\]
The complete equivariant family is the nonredundant finite-functional-calculus
image of the unmarked exceptional-core output and inherits the same role-level
occurrence witness.  Its invariant constituents \(\Pi_8,\Pi_4,P_0\), ranks,
spectra, and depths co-realize without a representative choice.  A formula
evaluated at one displayed \(u\) is a calculation representative; a physical
claim about that particular member requires an additional mark.  This
co-realization does not supply \((\mathrm{FLAG})\), does not choose an
external scale interpolation \((\mathrm{ACT}_{\rm Cl})\), and does not assert
a separate token occurrence for every later process parameter.
Relative to Paper~III-B, where the physical endpoint occurrence was still
unresolved, the equality \(P_{\rm term}=P_0\) is therefore a chronological
\texttt{DERIVED-LATER} result of the independently generated active orbit.
\end{corollary}

\begin{corollary}[Dynamical infrared projector and derived
\texorpdfstring{\((\mathrm{IR}_{O(3)})\)}{IR(O(3))} certification]
\label{p3:cor:derived-infrared-projector}
On the branch of Theorem~\ref{p3:thm:derived-active-support}, every orbit
member has the positive self-adjoint response
\begin{equation}
\label{p3:eq:active-infrared-semigroup}
 R_t(u):=\exp\!\bigl[-t(7I_A-\mathcal Q_{\rm act}(u))\bigr]
 =e^{-7t}Q_0(u)+e^{-3t}Q_4(u)+e^{-2t}Q_5(u)+Q_7(u).
\end{equation}
The family \(\{R_t(u)\}_{u\in\mathcal O_{\rm diag}}\) is equivariant; a
particular member requires a marked diagonal.  Every member preserves the unit
and satisfies the common norm limit
\begin{equation}
\label{p3:eq:active-infrared-limit}
 \boxed{\lim_{t\to\infty}R_t(u)
 =\mathbf1_{\{7\}}(\mathcal Q_{\rm act}(u))
 =\Pi_1(u)=P_0
 =\mathbb E_{\rm IR}.}
\end{equation}
The limit is unital, parity-even, and fully \(O(3)\)-equivariant; its range
is \(\mathbb R\mathbf1\) and its kernel has real dimension seven.  Hence
\((\mathrm{IR}_{O(3)})\) is a theorem-status certification of the terminal
active constituent.  It requires no independent G0 token beyond the occurring
active branch.
\end{corollary}

\begin{proof}
Equation~\eqref{p3:eq:active-infrared-semigroup} is the spectral expansion of
the exponential.  Since \(7I_A-\mathcal Q_{\rm act}(u)\) is nonnegative, the
semigroup is a positive self-adjoint contraction.  The unit lies in
\(Q_7(u)A=\mathbb R\mathbf1\), so \(R_t(u)\mathbf1=\mathbf1\).  All other
eigenvalues decay, proving the norm limit.  Theorem~\ref{p3:thm:derived-active-support}
identifies the limit with \(P_0\), while
Proposition~\ref{p3:prop:reeb-selection} identifies \(P_0\) with the full
orthogonal Haar projector and proves the covariance and kernel statements.
Equivariance of functional calculus gives
\[
 R_t(gug^{-1})
 =\operatorname{Ad}_gR_t(u)\operatorname{Ad}_{g^{-1}}.
\]
Because the exponential family and its common spectral limit are intrinsic
constituents of the frozen occurring orbit, GP2 co-realization propagates the
same G0 occurrence witness to the family and to the invariant limit without
selecting a member.
\end{proof}

\begin{remark}[Why the terminal rank distinguishes parity]
An \(SO(3)\)-only average on the full Clifford algebra would preserve both
\(1\) and the pseudoscalar \(e_{123}\), giving terminal rank two,
\(q=6\), and \(K_{\rm sh}=48\).  Every member of the independently
constructed grade-even active orbit instead has terminal rank one at \(q=7\).
Thus exclusion of the
pseudoscalar, and hence parity evenness of the infrared map, is a conclusion
of the active projector construction rather than an input used to select it.
\end{remark}

\begin{definition}[Active spectral-profile sort and residual enlargement]
\label{p3:def:active_corner}
Theorem~\ref{p3:thm:derived-active-support} fixes the equivariant orbit of four
discrete carrier ranks
and Corollary~\ref{p3:cor:derived-infrared-projector} fixes the infrared
endpoint.  Before evaluating a parameterized support relation, declare
\(\mathsf{ActSpecGraph}\) to be the sort of decreasing upper spectral
families obtained by real functional calculus from the generated
active-depth operator, with its positive depth orientation and fixed four
spectral projectors.  Theorem~\ref{p3c:thm:common-trace-defect-spectral-filtration}
constructs the unique member of this sort and its dimensionless pullback
along the generated trace-defect clock.

The optional tag \((\mathrm{ACT}_{\rm Cl})\) is henceforth reserved for an
enlarged profile
\[
 P_{\rm act}(\Lambda)\in\End_{\mathbb R}(A_{\rm prot}),
 \qquad
 \mathcal A_{\rm act}(\Lambda)
 :=P_{\rm act}(\Lambda)L(A_{\rm prot})P_{\rm act}(\Lambda).
\]
that is not the canonical spectral family, or for its adoption as an actual
external-scale process.  Such a profile may use a marked member's projectors
but may not change the orbit-invariant ranks, action types, or terminal value.
No external scale owner, calibration, marked diagonal, or actual transition
token is supplied by the dimensionless spectral construction.
\end{definition}

\begin{theorem}[Scalar remnant and seven-dimensional rank drop]
\label{p3:thm:hopf-decomposition}
On the occurring active branch of
Theorem~\ref{p3:thm:derived-active-support}, the protected Clifford block has
the canonical ultraviolet-to-infrared endpoint decomposition
\[
\Cl(0,3)=\RR\mathbf 1\oplus W,
\qquad
W:=\ker\mathbb E_{\mathrm{IR}},
\qquad
\dim_{\RR}W=7.
\]
Thus the protected non-scalar sector has total real rank seven.  The endpoint
and total rank drop are canonical.  If \((\mathrm{FLAG})\) is additionally
supplied, its seven one-dimensional quotients may be counted successively;
without that datum no preferred basis, order, or seven-event spectral flag is
claimed.  An \(SO(3)\)-only averaging operation would instead preserve the
pseudoscalar \(e_{123}\) and give \(2\oplus6\), but it is not the terminal
projector of the derived active orbit.
\end{theorem}

\begin{proof}
\leavevmode\par\noindent
Theorem~\ref{p3:thm:clifford-dim} identifies the ultraviolet protected corner
with an eight-dimensional real Clifford algebra.
Corollary~\ref{p3:cor:derived-infrared-projector} derives the infrared remnant
as the scalar channel \(\RR\mathbf 1\) and identifies the non-scalar
complement as
\[
  W=\ker\mathbb E_{\mathrm{IR}}.
\]
Lemma~\ref{p3:lem:o3-clifford-expectation} gives
\[
  \dim_{\RR}\Cl(0,3)=8,
  \qquad
  \dim_{\RR}\RR\mathbf 1=1,
\]
and hence
\[
\dim_{\RR}W
=
\dim_{\RR}\Cl(0,3)-\dim_{\RR}(\RR\mathbf 1)
=
8-1
=
7.
\]
For a different \(SO(3)\)-only averaging map, both \(1\) and \(e_{123}\) are
invariant, giving dimension \(8-2=6\) for the complement.  The exact Dirac staircase of
Theorem~\ref{p3:thm:exact-rg} fixes discrete Dirac thresholds but supplies no
intertwiner from those eigenspaces to one-dimensional subquotients of \(W\).
Thus seven successive losses arise only after \((\mathrm{FLAG})\) is adopted;
degenerate losses may otherwise occur at one threshold.
\end{proof}

\section{Determinant depth, active orbit filtration, and canonical continuation}
\label{p3c:sec:v196-depth-continuation}
\begin{definition}[Basis-independent cumulative rank depth]
\label{p3:def:cumulative-rank-depth}
For any orthogonal active projector \(P\) on \(A_{\rm prot}\), define
\begin{equation}
\label{p3:eq:cumulative-rank-depth}
 r_{\rm act}(P)
 :=\dim_{\RR}(PA_{\rm prot}),
 \qquad
 q(P):=8-r_{\rm act}(P).
\end{equation}
The integer \(q\) is the \emph{cumulative protected rank depth}.  For every
\(u\in\mathcal O_{\rm diag}\), the orbit member
\((\Pi_8,\Pi_4,\Pi_3(u),P_0)\) has \(q=(0,4,5,7)\).  Hence these values
descend to the unmarked orbit without selecting \(u\).  On a supplied
\((\mathrm{FLAG})\) rank-one refinement it takes the successive bookkeeping
values \(\{0,1,\ldots,7\}\).  It counts the
number of completed non-scalar quotient units with multiplicity.  It
does not record which vectors of \(W\) have been removed, and it remains defined
when several rank losses occur at one supplied threshold; in that case
some intermediate integers are skipped by the actual projector rank.
\end{definition}

\begin{theorem}[Protected determinant-depth clock]
\label{p3:thm:basis-free-sevenfold-shell-clock}
On the occurring active branch of
Theorem~\ref{p3:thm:derived-active-support}, the protected endpoint has rank
drop seven.  Define the dimensionless determinant-depth clock
\begin{equation}
\label{p3:eq:Ksh-rank-depth}
        \boxed{\Ksh:=8q,\qquad 0\le q\le7.}
\end{equation}
The active response canonically realizes \(q=0,4,5,7\), hence
\(\Ksh=0,32,40,56\).  The generated smooth resolution role below crosses the
integer depth list
\begin{equation}
\label{p3:eq:sevenfold-shell-ladder}
        \boxed{\Ksh\in\{8,16,24,32,40,48,56\}.}
\end{equation}
without selecting a basis of \(W\).  If \((\mathrm{FLAG})\) is additionally
supplied, the same list may also index an ordered rank-one projector
refinement; that stronger interpretation is not part of the scalar clock.
For any fixed-shell scheme coordinate \(\sigma_{\rm fin}\), the corresponding
neutral completed determinant readout is
\begin{equation}
\label{p3:eq:protected-action-density-shell-clock}
        H_{\Ksh}^{(\sigma)}
        :=\alpha_{\rm hier}(\sigma_{\rm fin})^{\Ksh/8}.
\end{equation}
At the endpoint,
\begin{equation}
\label{p3:eq:terminal-shell-action-relation}
        \boxed{H_{56}^{(\sigma)}
        =\alpha_{\rm hier}(\sigma_{\rm fin})^{7}
        =H_{56}^{(0)}e^{7\sigma_{\rm fin}}.}
\end{equation}
The source convention is \(H_{\Ksh}^{\rm REG}:=H_{\Ksh}^{(0)}\).  No
dimensionful reference magnitude or cosmological event name is fixed by this
theorem.  The generated neutral dual ratio and its exact orientation
classification are proved below.  A physical scale law is inherited only on
a generated same-owner branch; dimensionful normalization and any semantic
event identification remain typed downstream data.
\end{theorem}

\begin{proof}
Theorem~\ref{p3:thm:hopf-decomposition} gives the canonical endpoint rank
drop \(8-1=7\).  Paper~I-B's resolution-data theorem assigns dimension
\(8\dim Q\) to a rank-\(\dim Q\) scalar-free quotient, motivating the neutral
clock coordinate \(\Ksh=8q\).  The smooth semigroup of
Theorem~\ref{p3:thm:smooth-protected-resolution-semigroup} supplies the
ordered scalar coordinate; without \((\mathrm{FLAG})\), no intermediate
rank-one projector ordering is selected.
Theorem~\ref{p3:thm:determinant-rank-character} gives the neutral
rank-additive determinant character, so a rank-\(q\) completed non-scalar
quotient carries \(q\) copies of the same scheme-labelled factor.  This gives
\(H_{8q}^{(\sigma)}=\alpha_{\rm hier}(\sigma_{\rm fin})^q\) without invoking
a physical transfer rule or dimensionful reference scale.
\end{proof}

\begin{definition}[Visibility-flow calibration \((\mathrm{CAL}_{\rm vis})\)]
\label{p3c:def:visibility-flow-calibration}
The generated visibility role is naturally parameterized by the dimensionless
flow coordinate \(\tau\).  A physical time variable \(t\) is related to it
only after supplying a positive clock/rate calibration
\[
 \boxed{(\mathrm{CAL}_{\rm vis}):\qquad \tau=\kappa t,\quad \kappa>0.}
\]
This calibration does not alter the common-rate semigroup structure and does
not select a channel-dependent decay rate.
\end{definition}

\begin{theorem}[Generated basis-neutral resolution semigroup and rigidity]
\label{p3:thm:smooth-protected-resolution-semigroup}
At the beginning of this chronological stage, freeze the active-response
constructor schema that licenses formation of the retained complement
\(1-P_0\), continuous functional calculus in that retained projector, the
dimensionless additive parameter monoid, and the resulting process graph.
This is the exact model-specific constructor input used below; it is not a
consequence of mathematical canonicity or of GP1--GP3.

Before evaluating any flow member, declare the target sort
\(\mathsf{DynGraph}_{1\oplus7}^{+}\) to consist of dimensionless positive
contraction-semigroup roles on a bare unmarked split
\(\mathbb R\mathbf1\oplus W\), with distinguished stationary scalar line,
positive resolution orientation, and endpoint projector onto that line.  On
the occurring active orbit of Theorem~\ref{p3:thm:derived-active-support}, let
\(P_0\) be the representative-independent terminal projector of
Corollary~\ref{p3:cor:derived-infrared-projector}, put \(P_W=1-P_0\), and
regard \(W=P_WA_{\rm prot}\) as unmarked.  Its presentation automorphisms are
all orthogonal transformations of \(W\).

Define the forgetful/process constructor
\begin{equation}
\label{p3c:eq:visibility-derived-constructor}
 \mathfrak F_{\rm vis}:
 \mathscr C_{\rm act}\longrightarrow
 \mathsf{DynGraph}_{1\oplus7}^{+},
 \qquad
 \mathfrak F_{\rm vis}(\mathscr C_{\rm act})
 :=\mathcal V(P_0),
\end{equation}

\paragraph{Momentum-radius normalization of visibility time.}
Define
\[
 \boxed{\widehat\kappa_{\rm vis}:=\kappa\frac{R_{\rm mom}}{c_{\rm op}}.}
\]
The residual visibility-rate freedom is therefore dimensionless in the native
clock unit \(R_{\rm mom}/c_{\rm op}\).  This retyping does not derive the value
of \(\widehat\kappa_{\rm vis}\) or identify the visibility process with any
other relaxation channel.
where
\begin{equation}
\label{p3c:eq:generated-resolution-role}
 \mathcal V(P_0)
 :=\{V_\tau:=e^{-\tau P_W}:\tau\ge0\}
\end{equation}.
Let \(w_{\rm vis}\) be the displayed finite typed word
\[
 \mathscr C_{\rm act}\longmapsto P_0
 \longmapsto P_W=1-P_0
 \longmapsto
 \bigl(\mathcal V(P_0),\operatorname{Graph}\mathcal V\bigr).
\]
Every input is supplied by the complete active branch, each intermediate is
retained and reusable in its native sort, and the required output was declared
before evaluation.  Therefore
\[
 \operatorname{Orig}_{\Sigma_{\rm vis}}
 (w_{\rm vis};\mathscr C_{\rm act},
  \mathsf{DynGraph}_{1\oplus7}^{+})
\]
holds.
The constructor forgets orbit-member labels and the three marked-chain rates
but retains the common scalar/non-scalar split and the positive-decay
orientation already witnessed by every \(R_t(u)\).  Its graph with the
active-orbit family has an explicit GP1-admissible, GP2-exact realization, so
the FRC obstruction is absent.  Paper~II-A,
Theorem~\ref{p2a:thm:rooted-generated-image-corealization}, therefore retains
and co-realizes \(\mathcal V(P_0)\) in the declared process sort.  No
additional process-occurrence premise is required.  This co-realizes the
generated family and its operation law, not a separate contingent token for
every \(\tau\).

The assignment is also rigid.  If \((T_t)_{t\ge0}\) is any norm-continuous
real-linear semigroup on \(\mathbb R\mathbf1\oplus W\) that preserves the
unit, is natural under every unmarked automorphism, and has one decay witness
\(\|T_{t_*}|_W\|<1\), then there is a unique \(\kappa>0\) such that
\[
 T_t=P_0+e^{-\kappa t}P_W.
\]
Thus, after the separately supplied visibility-flow calibration
\((\mathrm{CAL}_{\rm vis})\), every such
candidate is the generated family \eqref{p3c:eq:generated-resolution-role}.
Its generator is uniquely
\begin{equation}
       K_W:=P_W,
       \qquad K_W|_W=1_W,
\end{equation}
and the visibility semigroup is
\begin{equation}
\label{p3:eq:smooth-visibility-semigroup}
 V_\tau:=P_0+e^{-\tau K_W}P_W,
 \qquad \tau\ge0.
\end{equation}
Self-adjointness, positivity, contraction, a nonnegative generator, and the
absence of channel-dependent rates are conclusions.  The only residual freedom is the positive calibration
\((\mathrm{CAL}_{\rm vis})\) of the flow parameter.  The common-rate
role is distinct from every marked multi-rate active response \(R_t(u)\) of
Corollary~\ref{p3:cor:derived-infrared-projector}.
It satisfies
\[
 V_0=1,
 \qquad
 V_{\tau_1+\tau_2}=V_{\tau_1}V_{\tau_2},
 \qquad
 \lim_{\tau\to\infty}V_\tau=P_0.
\]
For every finite \(\tau>0\), its spectrum is exactly
\[
 \operatorname{spec}(V_\tau)
 =\{1,e^{-\tau}\},
\]
with multiplicities \(1\) and \(7\).  Hence the semigroup has one
sevenfold-degenerate non-scalar spectral block and does not generate seven
spectral-projector rank changes.
Its rank-trace cumulative resolved depth is exactly
\begin{equation}
\label{p3:eq:smooth-cumulative-depth}
 q_{\rm sp}(\tau)
 :=\Tr_W\!\bigl(1-e^{-\tau K_W}\bigr)
 =7(1-e^{-\tau}).
\end{equation}
The rank-trace coordinate crosses the integers at
\begin{equation}
       \tau_j=-\log\!\left(1-\frac j7\right),
       \qquad j=1,\ldots,6,
\end{equation}
while \(q_{\rm sp}=7\) is approached only as \(\tau\to\infty\).  The smooth
bookkeeping coordinate is
\begin{equation}
\label{p3:eq:smooth-shell-clock}
 \Ksh(\tau):=8q_{\rm sp}(\tau),
\end{equation}
and its integer level sets are the scalar determinant-depth landmarks of
Theorem~\ref{p3:thm:basis-free-sevenfold-shell-clock}.  They require no
\((\mathrm{FLAG})\), but they are not projector-rank changes.  The exact Dirac
staircase of Theorem~\ref{p3:thm:exact-rg} is a separate sharp spectral count;
no equality between its eigenvalue thresholds and these protected
rank-trace crossings is asserted.
\end{theorem}

\begin{proof}
We verify the constructor-origin certificate and its nonobstructed exact
realization before proving rigidity.

First, the target process sort and its retained relations were declared before
the constructor was evaluated.  The map
\(P_0\mapsto P_W=1-P_0\mapsto e^{-\tau P_W}\) uses only subtraction,
continuous functional calculus, and the additive monoid
\((\mathbb R_{\ge0},+)\).  Since \(P_0\) is constant on
\(\mathcal O_{\rm diag}\), the map descends to the unmarked orbit quotient
and contains no representative choice, scale, or channel-dependent rate.
For every orthogonal isomorphism \(U\) preserving the scalar line,
\[
 Ue^{-\tau P_W}U^{-1}=e^{-\tau UP_WU^{-1}},
\]
so the constructor is host-natural under all unmarked automorphisms.

Second, every source response \(R_t(u)\) fixes \(\mathbb R\mathbf1\) and
strictly contracts every nonzero spectral vector in \(W\).  The
claim-relative G0 comparison “stationary scalar versus positively decaying
non-scalar” is therefore instantiated uniformly on the complete orbit.
Its image under \(\mathfrak F_{\rm vis}\) is the same typed comparison because
\(V_\tau\mathbf1=\mathbf1\) and
\(\|V_\tau|_W\|=e^{-\tau}<1\) for \(\tau>0\).  Conversely,
\[
 -\partial_\tau V_\tau|_{\tau=0}=P_W,
 \qquad
 \lim_{\tau\to\infty}V_\tau=P_0,
\]
so the output reflects the split, endpoint, and orientation asserted by the
target claim.  The forgotten relative rates \(7,3,2\) are not target content
and are not erased from the separately retained source role.

Third, the essential image is exactly the bare split, the dimensionless
additive parameter monoid, and the displayed semigroup graph.  No physical
time, calibration, flag, marked diagonal, state, or empirical token is added.
The semigroup law and operation envelope give a finite exact structured
presentation, so the output also satisfies the finitely presented
generated-role definition.

Fourth, the simultaneous source/target presentation is the graph of this
typed forgetful constructor.  It leaves every \(R_t(u)\) intact, declares
\(R_t(u)\) and \(V_\tau\) to be distinct roles, and inserts no equality
between their rates.  Their only shared relation is the already generated
endpoint \(P_0\) and positive resolution orientation.  Hence the combined
presentation is GP1-admissible and its insertion, omission, relation, typing,
and closure defect components all vanish.  This realization rules out
\(\operatorname{Obs}_{\rm FRC}(w_{\rm vis},\mathscr C_{\rm act})\).
Paper~II-A
Theorem~\ref{p2a:thm:rooted-generated-image-corealization} therefore retains
and co-realizes the exact required output from the same active-branch root.
The naturality, conservativity, and zero-defect calculations verify that
realization; none of them creates the frozen constructor schema.

For rigidity, the commutant of the standard \(O(W)\)-module \(W\) is
\(\mathbb R1_W\), and \(W^{O(W)}=0\).  Hence every natural endomorphism of
\(\mathbb R\mathbf1\oplus W\) is block scalar.  Unit preservation gives
\[
 T_t=P_0+b(t)P_W,
 \qquad b(0)=1.
\]
The semigroup law becomes \(b(t+s)=b(t)b(s)\).  Continuity at zero excludes
zeros, while \(b(t)=b(t/2)^2>0\).  Hence
\(-\log b(t)\) is a continuous additive function on
\(\mathbb R_{\ge0}\), so \(b(t)=e^{-\kappa t}\) for a real \(\kappa\).
The single decay witness gives \(\kappa>0\).  With \(\tau=\kappa t\),
\[
 T_t=P_0+e^{-\kappa t}P_W,
\]
so the calibrated candidate is exactly \(V_\tau=e^{-\tau P_W}\).  The
orthogonal-projector form proves positivity, self-adjointness, contractivity,
and the nonnegative generator \(K_W=P_W\).
The displayed formula also shows directly that \(V_\tau\) has eigenvalue
\(1\) on the scalar line and the single eigenvalue \(e^{-\tau}\) on all of
\(W\).  Taking the ordinary rank trace on \(W\) yields
\(q_{\rm sp}=7(1-e^{-\tau})\), whose derivative is
\(7e^{-\tau}>0\).  Solving \(q_{\rm sp}=j\) gives the displayed crossing
times, and the endpoint follows as \(\tau\to\infty\).
\end{proof}

\begin{corollary}[Protected visibility as a universal-continuation instance]
\label{p3c:cor:visibility-universal-continuation-instance}
For the occurring active branch, take the generated continuation signature of
Paper~III-A Definition~\ref{p3a:def:generated-exact-continuation-signature}
with
\[
\begin{aligned}
 \mathcal S_{\rm vis}&=(\mathbb R_{\ge0},+),&
 X_{\rm vis}&=\mathbb R\mathbf1\oplus W,\\
 \mathcal G_{\rm vis}&=O(W),&
 \mathsf{DynGraph}_{\rm vis}
   &=\mathsf{DynGraph}_{1\oplus7}^{+},
\end{aligned}
\]
and with the stationary scalar line, positive resolution orientation,
endpoint \(P_0\), and stationary-versus-decaying comparison as its frozen
diagram.  Its admissible generator cone is
\[
 \mathfrak g_{\rm vis}
 =\mathbb R_{\ge0}(-P_W).
\]
The transported active-branch decay comparison excludes the zero generator.
Consequently Paper~III-A
Theorem~\ref{p3a:thm:finite-generator-one-ray-rigidity} gives
\[
 T_t=P_0+e^{-\kappa t}P_W,
 \qquad \kappa>0,
\]
and the declared calibration \(\tau=\kappa t\) identifies the complete
moduli with \(V_\tau=e^{-\tau P_W}\).  The G0 witness-transport already proved
in Theorem~\ref{p3:thm:smooth-protected-resolution-semigroup} co-realizes this
law in its process sort.

This instantiation creates no cross-role equality with a process defined on a
different parent.  Any shared clock, rate, generator, or descendant scalar
defect requires a separately generated typed bridge and the fibre-product
test of Paper~II-A
Theorem~\ref{p2a:thm:generated-exact-solution-moduli}.
\end{corollary}

\begin{proof}
The commutant calculation in
Theorem~\ref{p3:thm:smooth-protected-resolution-semigroup} gives
\(\operatorname{End}_{O(W)}(W)=\mathbb R I_W\); unit preservation fixes the
scalar block and contraction orientation fixes the nonpositive generator
ray.  The same theorem supplies a strict decay witness and verifies all G0
transport clauses.  The universal-continuation and calibration conclusions
therefore apply directly.  The final statement is the relation-free/fibre-
product boundary of the cited Paper-II-A theorem.
\end{proof}

\begin{corollary}[Primitive residual at the final finite crossover]
\label{p3c:cor:primitive-final-octet-residual}
For the generated semigroup of
Theorem~\ref{p3:thm:smooth-protected-resolution-semigroup}, define the
unresolved trace and the unmarked positive unit-trace residual by
\[
 N_W(\tau):=\Tr_W(V_\tau|_W)=7e^{-\tau},
 \qquad
 D_*:=\frac1{7}P_W,
 \qquad \Tr_W(D_*)=1.
\]
At the $j$-th finite trace crossing $q_{\rm sp}(\tau_j)=j$,
\begin{equation}
\label{p3c:eq:finite-crossing-residual-ladder}
 P_WV_{\tau_j}
 =\frac{7-j}{7}P_W
 =(7-j)D_*,
 \qquad j=1,\ldots,6.
\end{equation}
Consequently
\begin{equation}
\label{p3c:eq:primitive-k48-residual}
 N_W(\tau)=1
 \quad\Longleftrightarrow\quad
 \tau=\ln7
 \quad\Longleftrightarrow\quad
 q_{\rm sp}=6
 \quad\Longleftrightarrow\quad
 \Ksh=48,
\end{equation}
and this is the unique finite nonzero integer crossing at which the unresolved
operator is $D_*$ itself.  Here
\begin{equation}
 \label{p3c:eq:primitive-residual-not-projector}
 D_*^2=\frac17D_*,\qquad
 \operatorname{rank}_{\mathbb R}D_*=7,\qquad
 \Tr_WD_*=1,
\end{equation}
so \(D_*\) is a normalized positive residual density, not a rank-one (or any
orthogonal) projector.  The only orthogonal projectors on \(W\) natural under
the full unmarked \(O(W)\) presentation symmetry are \(0\) and \(P_W\).
At the same smooth depth, Theorem~\ref{p3c:thm:active-depth-six-obstruction}
gives
\begin{equation}
 \label{p3c:eq:k48-smooth-active-separation}
 q_{\rm sp}=6,\quad \Ksh=48,
 \qquad
 \mathbf1_{\{6\}}(\mathcal Q_{\rm act}(u))=0.
\end{equation}
Thus the crossing is transverse because \(N_W'(\ln7)=-1\), but it is a
neutral smooth trace/resolution crossing with no native active carrier.  It
does not by itself constitute a particle, reciprocal-position current,
Hilbert source, transfer history, or cosmological event.
\end{corollary}

\begin{proof}
On $W$, $V_\tau=e^{-\tau}I_W$, so
$N_W=7e^{-\tau}$ and
$q_{\rm sp}=7-N_W$.  Substituting $q_{\rm sp}=j$ gives
Eq.~\eqref{p3c:eq:finite-crossing-residual-ladder}.  Since $N_W$ is
strictly decreasing, $N_W=1$ has the unique solution $e^{-\tau}=1/7$,
which is equivalent to the displayed depth and shell labels.  Among
$j=1,\ldots,6$, the equation $(7-j)D_*=D_*$ holds only for $j=6$.
Equation~\eqref{p3c:eq:primitive-residual-not-projector} follows from
\(P_W^2=P_W\).  If an orthogonal projector is natural under every
\(O(W)\)-automorphism, its restriction commutes with the standard
\(O(W)\)-representation and is therefore a scalar multiple of \(I_W\);
idempotence leaves only zero and one.  The final active statement is
\eqref{p3c:eq:active-depth-six-absent}.
\end{proof}

\section{Spectral filtration, neutral hierarchy, and finite thresholds}
\label{p3c:sec:v196-filtration}
\begin{theorem}[Common trace defect and canonical active spectral filtration]
\label{p3c:thm:common-trace-defect-spectral-filtration}
Let
\(A_{\rm prot}=\mathbb R\mathbf1\oplus W\), \(\dim_{\mathbb R}W=7\), be the
unmarked orthogonal split already generated above.  There is a unique
real-linear functional
\[
 \ell_W:\End_{\mathbb R}(W)_{\rm sa}\longrightarrow\mathbb R
\]
that is invariant under orthogonal conjugation and has
\(\ell_W(E)=1\) for every rank-one orthogonal projector \(E\).  It is
\(\ell_W=\Tr_W\).  Consequently, among defects affine-linear in the
unresolved part \(I_W-T|_W\), the unique normalized, orthogonally invariant,
spectrally additive defect on self-adjoint positive contractions satisfying
\(T\mathbf1=\mathbf1\) is
\begin{equation}
\label{p3c:eq:common-trace-defect}
 \boxed{\delta_W(T):=\Tr_W(I_W-T|_W).}
\end{equation}
It simultaneously gives
\begin{equation}
\label{p3c:eq:trace-defect-two-restrictions}
 \delta_W(P)=8-\operatorname{rank}_{\mathbb R}P=q(P),
 \qquad
 \delta_W(V_\tau)=7(1-e^{-\tau})=q_{\rm sp}(\tau)
\end{equation}
for every active support
\(P\in\{\Pi_8,\Pi_4,\Pi_3(u),P_0\}\) and every \(\tau\ge0\).

For each orbit member define the upper spectral family of the already
generated active-depth operator by
\begin{equation}
\label{p3c:eq:canonical-active-upper-filtration}
 \boxed{\mathsf P_u(s):=\mathbf1_{[s,7]}
       (\mathcal Q_{\rm act}(u)),\qquad 0\le s\le7.}
\end{equation}
It is decreasing, left-continuous at every interior spectral threshold,
equivariant under the unmarked orbit action, and is the unique upper spectral
resolution of \(\mathcal Q_{\rm act}(u)\) with positive depth orientation.
Explicitly,
\begin{equation}
\label{p3c:eq:canonical-active-upper-filtration-steps}
 \mathsf P_u(s)=
 \begin{cases}
  \Pi_8,&s=0,\\
  \Pi_4,&0<s\le4,\\
  \Pi_3(u),&4<s\le5,\\
  P_0,&5<s\le7.
 \end{cases}
\end{equation}
The canonical dimensionless pullback along the smooth defect clock is
\begin{equation}
\label{p3c:eq:canonical-active-flow-pullback}
 \boxed{P_{\rm spec}(\tau;u)
 :=\mathsf P_u\!\left(q_{\rm sp}(\tau)\right).}
\end{equation}
The hard and smooth defects agree exactly at the four spectral anchors
\begin{equation}
\label{p3c:eq:hard-smooth-defect-anchors}
\begin{array}{c|c|c}
 d & \tau & \mathsf P_u(d)\\ \hline
 0&0&\Pi_8\\
 4&\log(7/3)&\Pi_4\\
 5&\log(7/2)&\Pi_3(u)\\
 7&+\infty&P_0.
\end{array}
\end{equation}
Here \(\tau=+\infty\) denotes the endpoint limit.  Between spectral anchors
\(\delta_W(\mathsf P_u(s))\) is stepwise and need not equal \(s\).
Thus \eqref{p3c:eq:canonical-active-flow-pullback} is a spectral filtration,
not a continuously rank-varying interpolation or a chronology of physical
transitions.

The complete equivariant family
\(\{\mathsf P_u\}_{u\in\mathcal O_{\rm diag}}\), and not a marked member, is
the generated output in \(\mathsf{ActSpecGraph}\).  It uses no
\((\mathrm{ACT}_{\rm Cl})\), \((\mathrm{FLAG})\), external scale owner,
dimensionful calibration, or new occurrence root.
\end{theorem}

\begin{proof}
Let \(B=B^*\in\End_{\mathbb R}(W)\).  The real spectral theorem gives an
orthonormal eigenbasis and a rank-one resolution
\(B=\sum_{j=1}^{7}\lambda_jE_j\).  Orthogonal conjugation is transitive on
rank-one projectors, so the normalization and linearity give
\[
 \ell_W(B)=\sum_{j=1}^{7}\lambda_j\ell_W(E_j)
 =\sum_{j=1}^{7}\lambda_j=\Tr_W(B).
\]
This also proves existence and independence of the chosen eigenbasis.  If
\(T=T^*\), \(0\le T\le I\), and \(T\mathbf1=\mathbf1\), then self-adjointness
makes \(W\) invariant, so \eqref{p3c:eq:common-trace-defect} is defined.  The
same uniqueness calculation proves the asserted classification of affine
spectral defects.

Every active projector contains the scalar line.  Therefore
\(P|_W\) has rank \(\operatorname{rank}P-1\), and
\[
 \delta_W(P)=7-(\operatorname{rank}P-1)=8-\operatorname{rank}P.
\]
For the already generated visibility semigroup
\(V_\tau|_W=e^{-\tau}I_W\), Paper~1B,
Proposition~\ref{p1b:prop:trace-coordinate}~\cite{Kugelmann2026Paper1B},
gives \(\delta_W(V_\tau)=7(1-e^{-\tau})\).  This proves
\eqref{p3c:eq:trace-defect-two-restrictions}: the discrete and smooth depths
are restrictions of one operation on one generated owner.

The four mutually orthogonal spectral projectors
\(Q_0,Q_4,Q_5,Q_7\) and eigenvalues \(0,4,5,7\) of
\(\mathcal Q_{\rm act}(u)\) give
\eqref{p3c:eq:canonical-active-upper-filtration-steps} by finite spectral
calculus.  The spectral theorem proves uniqueness, monotonicity, and
left-continuity with the displayed inclusive convention.  Equivariance follows
from
\[
 \mathbf1_{[s,7]}(g\mathcal Q_{\rm act}(u)g^{-1})
 =g\mathbf1_{[s,7]}(\mathcal Q_{\rm act}(u))g^{-1}.
\]
Substituting \(q_{\rm sp}(\tau)\) gives the pullback, and solving
\(7(1-e^{-\tau})=4,5\) gives the two finite anchor times.  The source and
endpoint anchors are immediate.

Finally, ordinary trace and finite functional calculus are unmarked natural
operations already present in the frozen active composite.  Their graph is
exactly the displayed essential image, reflects the depth spectrum, and
transports the same active-branch occurrence witness.  Hence the invariant
descent/exact-transport clauses of Paper~II-A
Theorem~\ref{p2a:thm:geometric-defect-repair-admission} admit the complete
orbit-equivariant filtration without a new premise.  Evaluating at a
particular \(u\), assigning a physical scale, or asserting an actual
transition token lies outside that essential image.
\end{proof}

\begin{corollary}[No continuous rank-changing projector path]
\label{p3c:cor:rank-changing-projector-continuity-no-go}
Let \(I\) be a connected interval and let
\(P:I\to\End_{\mathbb R}(A_{\rm prot})\) be a norm-continuous family of
orthogonal projectors.  Then \(\operatorname{rank}P(t)\) is constant on
\(I\).  Because \(A_{\rm prot}\) is finite dimensional, the same conclusion
holds for strong- or weak-operator continuous projector families.  Therefore
no operator-continuous path can realize the rank chain \(8,4,3,1\); the
canonical family above is necessarily stepwise.  A continuous parameter does
not imply a continuous projector-valued evolution.
\end{corollary}

\begin{proof}
If orthogonal projectors \(P,Q\) satisfy \(\|P-Q\|<1\), then
\(Q|_{\operatorname{Ran}P}\) is injective: \(Qx=0\) would give
\(\|x\|=\|(P-Q)x\|<\|x\|\) for nonzero \(x\).  Hence
\(\operatorname{rank}P\le\operatorname{rank}Q\); interchanging \(P,Q\)
gives equality.  Continuity makes the rank locally constant, and connectedness
makes it constant.  All operator topologies coincide on the finite-dimensional
endomorphism space.
\end{proof}

\subsection{Rank-chain heat checksum and determinant-measure firewall}
\label{p3c:sec:v346-rank-heat-check}
\begin{corollary}[Axis-invariant active-support heat polynomial]
\label{p3c:cor:v346-active-heat-checksum}
On the generated unmarked active Hurwitz orbit, the nested orthogonal
support ranks are $8\to4\to3\to1$ at depth labels $q=0,4,5,7$.
The active-depth operator has spectral multiplicities
\[
 \operatorname{mult}_{\mathcal Q_{\rm act}}(0,4,5,7)=(4,1,2,1),
\]
so the positive finite heat polynomial is exactly
\[
 \boxed{\operatorname{Tr}_{\mathbb R}
 e^{-t(7I-\mathcal Q_{\rm act})}
 =1+2e^{-2t}+e^{-3t}+4e^{-7t},\quad t\ge0.}
\]
The formula is independent of the cyclic-axis representative.  The
number $7$ is the \emph{codimension} of the terminal scalar line, not
seven independently proved rank-changing physical transition events;
those need the separate $(\mathrm{FLAG})$ refinement.
\end{corollary}
\begin{proof}
The differences of consecutive nested projectors have ranks $4,1,2,1$,
and $7I-\mathcal Q_{\rm act}$ takes the respective eigenvalues
$7,3,2,0$.  Taking the real trace proves the formula.  The orbit changes
projectors by conjugacy and therefore preserves their multiplicities.
For completeness, the rank-three projector is
$\Pi_3(u)=\Pi_4-\Pi_u$ with image
$\mathbb R1\oplus u^\perp$, not
$\mathbb R1\oplus\mathbb Ru\oplus\mathbb R\omega$:
the grade-odd pseudoscalar $\omega$ has already been removed by
$\Pi_4$.  The cyclic average on $u^\perp$ then yields $P_0$.
\end{proof}

\begin{remark}[Formal hierarchy numbers versus a realized field measure]
\label{p3c:rem:v346-hierarchy-measure-scope}
The retained REG half-determinant convention $a=1/2$ in
Proposition~\ref{p3c:prop:v2211-determinant-amplitude-measure} gives
\[
\begin{gathered}
 \alpha_{1/2}=\exp(4\pi^2+1/(6\sqrt2))
 \simeq1.5719582791225\times10^{17},\\
 \boxed{\alpha_{1/2}^{-7}\simeq4.216119310922\times10^{-121}.}
\end{gathered}
\]
The equally mathematical complex-Gaussian exponent $a=1$ yields
$\delta_{1}=1/(3\sqrt2)$ and
$\alpha_1^{-7}\simeq1.84772878116\times10^{-121}$.
Only the first is the \emph{declared formal compiler} and neither is
identified as a physical measure without an actual represented action;
$\sigma_{\rm fin}$ is set to zero only on the declared REG$_0$ branch.
\end{remark}


\begin{corollary}[Closed heat loop and normalized boundary state]
\label{p3c:cor:heat-loop-versus-normalized-state}
Under the hypotheses of
Theorem~\ref{p3:thm:smooth-protected-resolution-semigroup}, the protected
process is an atomic heat-process instance,
\[
 (V_\tau)_{\tau\ge0}
 =(e^{-\tau P_W})_{\tau\ge0}\in\mathcal P_{\rm heat},
\]
and its unmarked extensive readout is the closed loop
\[
 \boxed{Z_W(\tau):=\operatorname{Tr}V_\tau
 =1+7e^{-\tau}\in\Dheat.}
\]
The normalized operator
\[
 \boxed{\rho_\tau:=\frac{V_\tau}{\operatorname{Tr}V_\tau}
 =\frac{P_0+e^{-\tau}P_W}{1+7e^{-\tau}},
 \qquad\operatorname{Tr}\rho_\tau=1,}
\]
is instead a state/boundary-conditioned object.  For every effect
\(0\le E\le1\), its normalized readout is
\[
 p_\tau(E)=\operatorname{Tr}(E\rho_\tau)\in[0,1].
\]
The loop \(Z_W\) and state \(\rho_\tau\) are not type-identical, even though
one is obtained from the other by scalar normalization.
\end{corollary}

\begin{proof}
Since \(P_W^2=P_W\), functional calculus gives
\[
 e^{-\tau P_W}=P_0+e^{-\tau}P_W=V_\tau.
\]
Paper~II-D Theorem~\ref{p2d:thm:heat-semigroup-rigidity} identifies this as
the unique heat semigroup with generator \(P_W\); the unmarked closure theorem
of Paper~II-A gives the trace.  The multiplicities \(1\) and \(7\) yield
\(Z_W=1+7e^{-\tau}\).  Division by this positive scalar gives a positive
unit-trace operator.  The finite state/effect theorem then gives
\(0\le\operatorname{Tr}(E\rho_\tau)\le1\).  Normalization introduces state
provenance, so the boundary mark cannot be erased.
\end{proof}

\begin{theorem}[Protected-octet resolved/residual trace reflection]
\label{p3c:thm:protected-octet-trace-reflection}
For the generated split
\(A_{\rm prot}=\mathbb R\mathbf1\oplus W\), \(\dim W=7\), define the
cumulative resolved trace and the full extensive surviving trace by
\[
 q_{\rm sp}(\tau)
 :=\Tr_W\!\left(I_W-V_\tau|_W\right),
 \qquad
 Z_W(\tau):=\Tr_{A_{\rm prot}}V_\tau.
\]
Then the same generated visibility operator gives the exact identity
\begin{equation}
\label{p3c:eq:protected-octet-trace-complement}
 \boxed{q_{\rm sp}(\tau)+Z_W(\tau)=8}
 \qquad(\tau\ge0).
\end{equation}
At every finite integer crossing \(q_{\rm sp}(\tau_j)=j\), and at the
terminal limiting level \(j=7\) with the endpoint convention
\(\tau_7:=+\infty\),
\begin{equation}
\label{p3c:eq:protected-octet-reflection-levels}
 \boxed{Z_W(\tau_j)=8-j.}
\end{equation}
Consequently the full-octet complement induces the order-reversing
involution
\begin{equation}
\label{p3c:eq:protected-octet-reflection-involution}
 \boxed{J(j):=8-j,\qquad J^2(j)=j,}
 \qquad j\in\{1,\ldots,7\},
\end{equation}
with pairs \((1,7),(2,6),(3,5)\) and fixed point \(4\).

The non-scalar unresolved trace is instead
\begin{equation}
\label{p3c:eq:w-only-versus-full-octet-residual}
 N_W(\tau)=Z_W(\tau)-1=7-q_{\rm sp}(\tau).
\end{equation}
Thus the shift from \(7-j\) to \(8-j\) is exactly the stationary scalar
line, not an additional channel.  In particular, \(N_W=1\) at \(j=6\) is
a unit unresolved trace of the common sevenfold block, not a spectral
projector-rank jump.

Equations~\eqref{p3c:eq:protected-octet-trace-complement}--
\eqref{p3c:eq:protected-octet-reflection-involution} are identities between
typed scalar readouts of the same process.  They do not identify
\(q_{\rm sp}\) with \(Z_W\), construct a reversal of the contractive
semigroup, equate \(\tau_j\) with \(\tau_{J(j)}\), or assign a physical event
name to either level.
\end{theorem}

\begin{proof}
The generated semigroup has
\(V_\tau=P_0+e^{-\tau}P_W\).  Therefore
\[
 q_{\rm sp}=7(1-e^{-\tau}),
 \qquad
 Z_W=1+7e^{-\tau},
\]
and their sum is eight.  Substitution of
\(e^{-\tau_j}=1-j/7\) gives \(Z_W(\tau_j)=8-j\); the same equation holds at
the terminal limit \(e^{-\tau}\to0\).  Direct substitution proves
\(J(J(j))=j\), and subtraction of the stationary scalar contribution gives
Eq.~\eqref{p3c:eq:w-only-versus-full-octet-residual}.

More generally, on a split \(\mathbb R\mathbf1\oplus W_m\), the identical
trace construction gives
\(q_m+Z_m=m+1\).  Hence the constant eight here is forced by
\(\dim W=7\), rather than selected as an independent numerical input.  Since
the proof uses only scalar traces, it supplies none of the excluded process,
time, or event-owner identifications.
\end{proof}

\begin{corollary}[Co-realization of the protected visibility family]
\label{p3c:cor:corealized-protected-visibility-family}
Let \(P\) be the G0-occurring active protected composite.  Freeze in its
constructor grammar the generated role \(\mathcal V(P_0)\) of
Theorem~\ref{p3:thm:smooth-protected-resolution-semigroup}, the unmarked scalar
closure
\[
 r_\tau:=\operatorname{Tr}V_\tau:\mathbf1\to\mathbf1,
 \qquad \operatorname{FSR}_{\Sigma}(r_\tau),
\]
and the result-independent normalization constructor
\[
 \mathcal N(A):=\frac{A}{\operatorname{Tr}A}
 \quad\text{on positive nonzero trace-class operators.}
\]
Then the following generated constituents are simultaneously realized in
their native sorts:
\[
 V_\tau,\qquad
 Z_W(\tau)=\Tr V_\tau,\qquad
 \rho_\tau=\frac{V_\tau}{Z_W(\tau)}.
\]
If the frozen presentation also contains the entropy postprocessor on the
finite state space, then
\[
 S(\rho_\tau)=-\Tr(\rho_\tau\log\rho_\tau)
\]
is co-realized as well.  The closed loop, normalized state, and entropy value
remain different typed constituents; no equality or addition of their
numerical values is implied.

In particular, the active protected occurrence already co-realizes this
complete visibility family.  Once a later response composite reuses it, no
separate occurrence assumption is needed for its normalized active state.  A
separate empirical effect \(E\), or a different process law not canonically
generated from the same split, still requires a new G0 root or its own
generated-role proof.

In the universal bridge normal form, the chain
\[
 \mathscr C_{\rm act}\longmapsto P_0\longmapsto P_W
 \longmapsto
 \bigl(\mathcal V(P_0),\operatorname{Graph}\mathcal V\bigr)
\]
is one rooted dynamics component.  On every separating evaluation role
declared closed under this family, its contravariant dual semigroup is
co-realized on the same domain \(\tau\ge0\) by
Theorem~\ref{p2a:thm:evaluation-dual-corealization}.  This is not a Born
exchange or a reverse process: the generated contraction monoid contains no
negative-\(\tau\) member, and its endpoint \(P_0\) is noninvertible.
\end{corollary}

\begin{proof}
Functional calculus generates \(V_\tau=e^{-\tau P_W}\).  The displayed
scalar FSR diagram closes it without a boundary, while \(\mathcal N\) divides
by the strictly positive generated scalar \(Z_W(\tau)\).  When frozen, the
entropy postprocessor supplies the final displayed scalar.  These are
result-independent licensed constructors in the presentation.  They form a
finite nonredundant constituent family: their
outer provenance is respectively operation, closed-loop scalar, normalized
boundary state, and state-decorated loop.  Paper~II-A
Theorem~\ref{p2:thm:generated-constituent-corealization} therefore realizes
all displayed constituents simultaneously while preserving their sorts.  Its
scope clause gives the final statement about unrelated effects and process
laws.  The last paragraph is the process-graph case of
Theorem~\ref{p2a:thm:universal-corealization-bridge-normal-form}; the
evaluation-dual conclusion uses exactly the closure hypothesis stated there.
Neither that theorem nor the present constructor supplies an inverse
contraction or a negative parameter.
\end{proof}

\begin{theorem}[Generated neutral determinant-depth hierarchy and exact orientation]
\label{p3c:thm:generated-dual-depth-hierarchy}
Fix one determinant-shell scheme coordinate \(\sigma_{\rm fin}\), and write
\[
 L_1^{(\sigma)}
 :=4\pi^2+\frac1{6\sqrt2}+\sigma_{\rm fin},
 \qquad
 \alpha_\sigma:=e^{L_1^{(\sigma)}}.
\]
On the occurring active branch, the generated semigroup and determinant
character canonically form the positive neutral depth character
\begin{equation}
\label{p3c:eq:clocked-depth-character}
 D_\sigma(\tau)
 :=\exp\!\bigl(L_1^{(\sigma)}q_{\rm sp}(\tau)\bigr)
 =\alpha_\sigma^{q_{\rm sp}(\tau)}.
\end{equation}
For every fixed weight \(w>0\), its algebraic dual supplies the normalized
neutral ratio
\begin{equation}
\label{p3c:eq:generated-dual-scale-law}
 \boxed{
 \chi_{\sigma,w}(\tau)
 :=D_\sigma(\tau)^{-w}
 =\alpha_\sigma^{-wq_{\rm sp}(\tau)}.}
\end{equation}
The character and its dual ratio are finitely presented generated-role
descendants of the occurring protected composite and require no independent
\((\mathrm{HT}_{\rm phys})\) or scale-by-scale occurrence premise.

Equivalently, define the \emph{depth-variable} character
\(\widehat\chi_{\sigma,w}(q):=\alpha_\sigma^{-wq}\) for
\(0\le q\le7\), so
\(\chi_{\sigma,w}(\tau)=\widehat\chi_{\sigma,w}(q_{\rm sp}(\tau))\).
The depth character is unique among continuous positive partial additive
characters with \(\widehat\chi(0)=1\) and one-step depth value
\(\widehat\chi(1)=\alpha_\sigma^{-w}\); the latter is not a
statement at flow time \(\tau=1\). For \(j=1,\ldots,6\), let
\(\tau_j=-\log(1-j/7)\), so \(q_{\rm sp}(\tau_j)=j\).
At the six finite hits and terminal limit,
\begin{equation}
\label{p3c:eq:generated-dual-scale-levels}
 \chi_{\sigma,w}(\tau_j)
 =\widehat\chi_{\sigma,w}(j)=\alpha_\sigma^{-wj},
 \quad j=1,\ldots,6,
 \qquad
 \lim_{\tau\to\infty}\chi_{\sigma,w}(\tau)
 =\widehat\chi_{\sigma,w}(7)=\alpha_\sigma^{-7w}.
\end{equation}
Its orientation is completely classified by the one-step logarithm:
\begin{equation}
\label{p3c:eq:dual-hierarchy-orientation-trichotomy}
 \begin{array}{ccl}
 L_1^{(\sigma)}>0 &\Longleftrightarrow&
   \chi_{\sigma,w}\text{ is strictly decreasing},\\
 L_1^{(\sigma)}=0 &\Longleftrightarrow&
   \chi_{\sigma,w}\equiv1,\\
 L_1^{(\sigma)}<0 &\Longleftrightarrow&
   \chi_{\sigma,w}\text{ is strictly increasing}.
 \end{array}
\end{equation}
Equivalently, decrease holds precisely when
\begin{equation}
\label{p3c:eq:positive-reg-domain}
 \boxed{\sigma_{\rm fin}>
 -4\pi^2-\frac1{6\sqrt2}.}
\end{equation}
The source member \(\mathrm{REG}_0\) lies in this open domain.  This theorem
generates a neutral determinant-owner ratio; it does not place an external
physical scale in that owner.
\end{theorem}

\begin{proof}
Equation~\eqref{p3c:eq:clocked-depth-character} uses only the generated scalar
\(q_{\rm sp}\), the frozen rank-one logarithm, multiplication, and
exponentiation.  Taking the algebraic dual sends the positive character
\(D_\sigma\) to \(D_\sigma^{-1}\); raising the positive character to a fixed
real weight is a licensed scalar constructor.  Hence the displayed neutral
family has a finite generated presentation, and Paper~II-A
Theorems~\ref{p2:thm:generated-role-corealization} and
\ref{p2:thm:generated-constituent-corealization} give co-realization from the
same witness.

For uniqueness, the logarithm of any continuous positive depth character is
additive.  Its one-step value is \(-w\log\alpha_\sigma\), so the continuous
Cauchy equation gives \(-wqL_1^{(\sigma)}\).  Exponentiation proves
\eqref{p3c:eq:generated-dual-scale-law}.  Since
\(q_{\rm sp}'(\tau)=7e^{-\tau}>0\),
\[
 \frac{d}{d\tau}\log\chi_{\sigma,w}(\tau)
 =-wL_1^{(\sigma)}q_{\rm sp}'(\tau),
\]
whose sign proves the trichotomy and the exact domain
\eqref{p3c:eq:positive-reg-domain}.  Substitution of the six hits and terminal
limit gives \eqref{p3c:eq:generated-dual-scale-levels}.  None of these
operations chooses a physical owner, dimensionful reference, or event
predicate.
\end{proof}

\begin{corollary}[Owner-typed physical dual-scale inheritance]
\label{p3c:cor:owner-typed-physical-dual-scale}
Let \(X_0>0\) be a reference magnitude in a constituted physical readout role.
Assume a generated, nonzero, host-natural ancestry morphism places its
one-dimensional owner in the weight-\(w\) representation of the determinant
dual owner, and assume the frozen grammar contains no independent deformation
of that owner.  Then Paper~II-A's generated-owner normalization-inheritance
theorem forces
\begin{equation}
\label{p3c:eq:owner-typed-physical-scale-law}
 \boxed{
 X_w(\tau)
 =X_0\chi_{\sigma,w}(\tau)
 =X_0\alpha_\sigma^{-wq_{\rm sp}(\tau)}.}
\end{equation}
No second hierarchy-transfer normalization or scale-by-scale occurrence
premise is admissible.  At the scalar levels,
\begin{equation}
\label{p3c:eq:owner-typed-physical-scale-levels}
 X_{w,j}=X_0\alpha_\sigma^{-wj},
 \quad j=1,\ldots,6,
 \qquad
 X_{w,7}=X_0\alpha_\sigma^{-7w}.
\end{equation}
If the same owner grammar generates positive square- and fourth-root readouts
from a master weight-one readout, functional calculus gives
\begin{equation}
\label{p3c:eq:root-weight-readouts}
 X_{1/2,j}=X_{1/2,0}\alpha_\sigma^{-j/2},
 \qquad
 X_{1/4,j}=X_{1/4,0}\alpha_\sigma^{-j/4}.
\end{equation}
These levels are decreasing exactly on the domain
\eqref{p3c:eq:positive-reg-domain}.  The reference magnitudes and units remain
calibrations, and a cosmological or particle name requires a downstream typed
predicate.  Without the ancestry morphism, only the neutral ratio theorem is
available: occurrence, equal rank, or a declared weight alone does not prove
physical determinant ownership.
\end{corollary}

\begin{proof}
The ancestry morphism is the missing source-to-owner port; because it is
nonzero between one-dimensional owner lines, it is an isomorphism onto the
target owner.  Thus the antecedent itself establishes the required support and
owner typing.  Paper~II-A's generated-owner normalization-inheritance theorem
then transports the dual character without an independent scalar endomorphism.
Multiplication by the reference section gives
\eqref{p3c:eq:owner-typed-physical-scale-law}; evaluation at integer depth
gives the level formula.  Positive functional calculus in the already typed
owner gives the root weights.  The converse scope statement is precisely the
declared-output and cross-role firewall: no constructor-free owner map can be
inferred from a coincident number.
\end{proof}

\begin{corollary}[Neutral and owner-typed octet scale reflection]
\label{p3c:cor:owner-typed-octet-scale-reflection}
Let \(J(j)=8-j\) be the involution of
Theorem~\ref{p3c:thm:protected-octet-trace-reflection}.  The generated neutral
dual character obeys
\begin{equation}
\label{p3c:eq:neutral-octet-scale-reflection}
 \boxed{
 \widehat\chi_{\sigma,w}(j)\,\widehat\chi_{\sigma,w}(J(j))
 =\widehat\chi_{\sigma,w}(4)^2.}
\end{equation}
Whenever the ancestry hypotheses of
Corollary~\ref{p3c:cor:owner-typed-physical-dual-scale} place a dimensionful
readout in that owner, its inherited levels satisfy
\begin{equation}
\label{p3c:eq:owner-octet-scale-reflection}
 \boxed{X_{w,j}X_{w,J(j)}=X_{w,4}^{\,2}.}
\end{equation}
For the shell coordinate \(K=8j\), the induced label involution is
\begin{equation}
\label{p3c:eq:shell-octet-reflection}
 \boxed{K\longmapsto64-K,}
\end{equation}
with pairs \((8,56),(16,48),(24,40)\) and fixed point \(32\).
This is a scale-product covariance inside an already typed owner.  It does
not generate a physical owner map between the two labels or a physical event
at either crossing.
\end{corollary}

\begin{proof}
Since \(j+J(j)=8\),
\[
 \widehat\chi_{\sigma,w}(j)\widehat\chi_{\sigma,w}(J(j))
 =\alpha_\sigma^{-8w}
 =\bigl(\alpha_\sigma^{-4w}\bigr)^2.
\]
Multiplication by the same inherited reference magnitude on both factors
gives Eq.~\eqref{p3c:eq:owner-octet-scale-reflection}.  Substituting
\(K=8j\) gives \(8J(j)=64-K\).  The scope statement follows from the owner
firewall in Corollary~\ref{p3c:cor:owner-typed-physical-dual-scale}.
\end{proof}


\begin{corollary}[Typed determinant-owner inheritance and rank firewall]
\label{p3c:cor:typed-determinant-owner-inheritance}
Fix a determinant-shell scheme member \(\mathrm{REG}_\sigma\) and its one-step
logarithm \(L_1^{(\sigma)}\), so
\(\alpha_\sigma=e^{L_1^{(\sigma)}}\).  Let \(Q\) be a finite-dimensional real
multiplicity object that has already been generated and typed in the same
neutral determinant-multiplicity owner as the rank-one shell character, with
its direct-sum character structure inherited without a generated deformation
of that owner.  Then
\[
 \boxed{
 H_\sigma(Q)
 :=\exp\!\bigl(\dim_{\mathbb R}Q\,L_1^{(\sigma)}\bigr)
 =\alpha_\sigma^{\dim_{\mathbb R}Q}.}
\]
There is no second independently selectable one-step factor \(\alpha_Q\) on
this same-owner branch.

Conversely, for an external or differently typed finite object \(Q'\), the
bare equality \(\dim_{\mathbb R}Q'=n\) does \emph{not} imply that the physical
protected depth satisfies \(q_{\rm sp}=n\).  Such an identification requires a
separately generated ancestry map placing \(Q'\) in the actual hierarchy
multiplicity/depth owner.
\end{corollary}

\begin{proof}
The determinant formula is Theorem~\ref{p3:thm:determinant-rank-character}
with the scheme-labelled rank-one logarithm.  Once the owner typing is already
established, Paper~II-A's generated-owner normalization inheritance theorem
forbids replacing the inherited one-step determinant normalization by an
independent factor.  The final statement is the owner-typing firewall: rank is
an invariant inside the determinant category, whereas \(q_{\rm sp}\) is the
coordinate of the specifically generated protected resolution role.  Equal
integers do not generate the missing cross-role map.
\end{proof}

\begin{remark}[Typing of the remaining III-C readouts]
The corner tag \((\mathrm{CORN})\) certifies Paper~III-B's generated ambient
rank-one constituent and its normalized real-channel value.  The \((\mathrm{REG})\) determinant-shell value is a generated determinant closure with an affine scheme coordinate on the unanchored domain; the occurring Hurwitz normalizer supplies the anchored origin.  The transfer label \((\mathrm{HT})\) certifies
the generated neutral dual-line postprocessing and, when separately proved,
owner-typed physical inheritance; it supplies neither a selector nor an
ancestry map.  None creates a third primitive outer
readout role, and none is selected by the generic \(\Dcl\) form theorem.
\end{remark}

\begin{corollary}[Protected scalar hitting hierarchy and scheme-labelled ledger]
\label{p3:cor:protected-spectral-resolution-landmarks}
Under the derived active branch and the generated process certification of
Theorem~\ref{p3:thm:smooth-protected-resolution-semigroup}, the
rank-trace coordinate crosses the first six scalar depth levels at finite
flow parameter and approaches the seventh only
as its asymptotic endpoint:
\begin{equation}
\label{p3:eq:protected-resolution-landmark-ledger}
 \begin{aligned}
  (\Ksh,H_{\Ksh}^{(\sigma)})
   &=(8j,\alpha_{\rm hier}(\sigma_{\rm fin})^{j}),
     &&j=1,\ldots,6,
       \quad \tau=\tau_j<\infty,\\
  (\Ksh,H_{\Ksh}^{(\sigma)})
   &\longrightarrow(56,\alpha_{\rm hier}(\sigma_{\rm fin})^{7}),
     &&\tau\longrightarrow\infty.
 \end{aligned}
\end{equation}
Together the six finite crossings and the limiting endpoint form a canonical
seven-entry neutral scalar hierarchy on the fixed protected module.  The
generated dual-line ratios are \(\widehat\chi_{\sigma,w}(j)=\alpha_\sigma^{-wj}\).
For a physical reference readout, the same formula holds only under the
owner-ancestry hypotheses of Corollary~
\ref{p3c:cor:owner-typed-physical-dual-scale}.  If
\((\mathrm{FLAG})\) is supplied, the list may additionally index successive
rank-one quotients; it is not thereby turned into seven finite
spectral-projector events.  The crossing and neutral dual-line statements use
no \((\mathrm{HT}_{\rm phys})\).  Strict decrease additionally uses
\(L_1^{(\sigma)}>0\); a dimensionful normalization, owner-ancestry map, and
any named cosmological or particle interpretation remain downstream typed
data.
\end{corollary}

\begin{proof}
Theorem~\ref{p3:thm:smooth-protected-resolution-semigroup} gives the unique
rank-trace crossing of each integer \(j=1,\ldots,6\) and the limiting endpoint
\(j=7\).  Substitution into
Theorem~\ref{p3:thm:basis-free-sevenfold-shell-clock} gives the displayed
neutral hierarchy readouts.  The semigroup's non-scalar spectral block
remains seven-dimensional throughout finite flow time, so no crossing is
promoted to a projector-rank event.  The dual ratios and their sign
classification are Theorem~\ref{p3c:thm:generated-dual-depth-hierarchy}; the
physical statement is exactly the separately typed inheritance corollary.
\end{proof}

\begin{corollary}[Parity-even rank-seven determinant and physical readouts]
\label{p3:cor:hierarchy-factor}
On the derived active branch, the protected non-scalar rank is seven, so the
neutral topological determinant character is
\begin{equation}
        \boxed{H_7^{\rm top}:=\alpha_{\rm top}^{7}=\exp(28\pi^2).}
\end{equation}
At finite-shell scheme coordinate \(\sigma_{\rm fin}\), the completed neutral
readout is
\begin{equation}
        \boxed{H_7^{(\sigma)}
        :=\alpha_{\rm hier}(\sigma_{\rm fin})^{7}
        =H_7^{(0)}e^{7\sigma_{\rm fin}}.}
\end{equation}
The source scheme is \(H_7^{\rm REG}:=H_7^{(0)}
=\exp(7(4\pi^2+1/(6\sqrt2)))\).  Their generated duals are neutral scale
ratios.  They act on a physical scale role only after the generated
owner-ancestry premise of Corollary~
\ref{p3c:cor:owner-typed-physical-dual-scale}; on that branch no additional
hierarchy-transfer normalization is allowed.  No dimensionful radius or
absolute scale is assigned without a downstream reference calibration.
\end{corollary}

\begin{proof}
Theorem~\ref{p3:thm:hopf-decomposition} gives \(\dim_{\mathbb R}W=7\).
Theorem~\ref{p3:thm:determinant-rank-character} assigns one common neutral
factor to each completed rank unit, giving the seventh power.  The scheme
formula follows from Definition~\ref{p3:def:alpha-two}; the neutral dual ratio
follows from Theorem~\ref{p3c:thm:generated-dual-depth-hierarchy}, and its
physical inheritance follows only from the owner-typed corollary.  No ordering
of seven projector losses is needed.
\end{proof}

\begin{remark}
The exponent seven is the endpoint rank drop from the protected Clifford
module to the infrared scalar remnant.  Outside the displayed occurring
active branch, the corresponding rank and generated scale statement are not
asserted.
\end{remark}

\paragraph{Physical momentum threshold and auxiliary finite-window family}
\label{p3:app:decoupling}

The generated momentum seed has level-one index \(m_{\rm seed}=1\).
Consequently its associated spinor comparison window which contains only the lowest
Dirac shell.  This indexing does not actualize a second physical spinor carrier.  A
multi-step Dirac staircase belongs only to the auxiliary mathematical family
\(\mathcal S_m\), whose exact formulas are proved in
Paper~III-G, Appendix~\ref{p3c:app:aux-dirac-windows}.

\begin{theorem}[Sharp-cutoff level-one associated Dirac threshold]
\label{p3:thm:exact-rg}
On the level-one associated spinor comparison window \(\mathcal S_1\), the retained round
\(S^3(\Rmom)\) Dirac shell is only
\[
        k=0,
        \qquad
        |\kappa_0|=\frac{3}{2\Rmom}.
\]
Under the separately supplied sharp spectral-frontier occurrence
\(\mathsf{SF}_\Lambda\), applied to the loop-derived finite spectral class
\(\Dcl\), every sharp-count coupling is therefore constant away from that
single associated-spinor Dirac threshold and has a distributional derivative supported only at
\(\Lcut=3/(2\Rmom)\).  No multi-step momentum staircase follows from the associated level-one comparison.
\end{theorem}

\begin{proof}
For the associated spinor family, the exact auxiliary shell bound is
\(0\le k\le\lfloor m-\tfrac12\rfloor\), proved in
Theorem~\ref{p3:thm:exact-fuzzy-s3-shell-trace}.  Setting the derived generated-seed level \(m_{\rm seed}=1\)
gives \(k_{\max}=0\).  Treating this level-one window as an
exhaustive full-physical-endpoint spectrum remains conditional on the
GP3 endpoint comparison; the present cutoff assertion is about the
stated associated window.  A finite sharp spectral count
is a finite sum of Heaviside functions of \(\Lcut-|\kappa_k|\); Paper~II-D
types this as the sharp spectral-resolution process rather than deriving the
projector from \(\Dcl\) alone.  With only
\(k=0\) present, its distributional derivative is supported only at the one
displayed eigenvalue magnitude.  The result is a finite-\(\Dcl\) statement
and does not use the smooth \((\mathrm{REG})\) determinant construction.
\end{proof}

\begin{remark}[Auxiliary multi-shell windows]
Generic \(\mathcal S_m\) windows with \(m>1\) may have several exact Dirac
thresholds and are useful comparison mathematics.  They are not additional
physical momentum carriers and do not constitute an \(m_p\)-flow.
\end{remark}


\section{Fixed-shell determinant and octahedral compilation}
\label{paper3c:sec:organized-determinant}
\subsection{Affine scheme geometry of the fixed-shell determinant family}
\label{p3c:sec:reg-affine-scheme-update}

Write the finite-shell scheme coordinate as $\sigma=\sigma_{\rm fin}$ and
let $T_c(\sigma)=\sigma+c$.  Then
\[
 \delta_{\sigma+c}=\delta_\sigma+c,
 \qquad
 L_1^{(\sigma+c)}=L_1^{(\sigma)}+c,
 \qquad
 \alpha_{\sigma+c}=e^c\alpha_\sigma.
\]

\begin{proposition}[No natural origin of the REG torsor]
\label{p3c:prop:no-natural-reg-origin}
The fixed-shell family is an affine torsor for $(\RR,+)$.  Consequently no
selection rule natural under all finite scheme translations can distinguish
$\sigma=0$ or any other member.
\end{proposition}

\begin{proof}
A translation-invariant selected point would satisfy $\sigma_*=\sigma_*+c$
for every $c\in\RR$, which is impossible.
\end{proof}

For determinant weight $w$, let the associated positive line transform as
$s_w(\sigma+c)=e^{-wc}s_w(\sigma)$.  Since $d\sigma$ is independent of the
choice of affine origin, the connection
\begin{equation}
 \label{p3c:eq:scheme-weight-connection}
 \nabla^{(w)}=d+w\,d\sigma
\end{equation}
is canonically defined on the scheme torsor and has horizontal local sections
proportional to $e^{-w\sigma}$.

\begin{theorem}[Exact scheme-continuation quotient]
\label{p3c:thm:exact-scheme-continuation-quotient}
Let \(\mathcal G_{\rm sch}\) be the translation groupoid of the positive
scheme domain \(I_+\): an arrow \(c:\sigma\to\sigma+c\) exists whenever both
endpoints lie in \(I_+\).  On a generated determinant-weight-\(w\) associated
line, define
\[
 P_c^{(w)}s=e^{-wc}s.
\]
Then
\[
 P_0^{(w)}=I,\qquad
 P_d^{(w)}P_c^{(w)}=P_{c+d}^{(w)},\qquad
 \nabla^{(w)}=d+w\,d\sigma,
\]
and this is the unique exact continuation compatible with the already fixed
weight-\(w\) associated representation.

For a scalar descendant with total associated weight
\(W=w_1+\cdots+w_k\), translation invariance is equivalent to
\[
 \boxed{W=0.}
\]
This is a statement about an invariant scalar coordinate on the scheme
quotient.  By contrast, an associated section may be horizontal at any
weight:
\[
 \nabla^{(W)}s=0
 \quad\Longrightarrow\quad
 s(\sigma)=s_0e^{-W\sigma}.
\]
Thus zero total weight is equivalent to scalar translation invariance, but
horizontality of an associated section is \emph{not} equivalent to zero
weight.  The complete \(\mathrm{REG}_\sigma\) family is one covariant object
over the torsor; a raw numerical value at a chosen \(\sigma\), an owner
comparison, or a boundary renormalization condition remains additional data.
In particular this theorem does not select \(\mathrm{REG}_0\).
\end{theorem}

\begin{proof}
The identity and composition formulas are the additive character law
\(e^{-w(c+d)}=e^{-wd}e^{-wc}\); differentiation gives the displayed flat
connection.  Tensor-product weights add, so a scalar coordinate transforms by
\(e^{-Wc}\), which is trivial for every translation exactly when \(W=0\).
On the other hand the first-order equation
\((d+W\,d\sigma)s=0\) integrates to \(s_0e^{-W\sigma}\) for arbitrary
\(W\).  For example \(s=e^{-\sigma}\) is a nonzero-weight horizontal section.
The no-origin proposition excludes promotion of an evaluation point and
leaves numerical boundary data outside the quotient theorem.
\end{proof}

\begin{proposition}[Full scheme--depth determinant connection and the
same-representation firewall]
\label{p3c:prop:full-scheme-depth-determinant-connection}
Let
\[
\begin{aligned}
 I_+&:=\left(-4\pi^2-\frac1{6\sqrt2},\infty\right),
 & B&:=I_+\times[0,7],\\
 \Phi(\sigma,q)&:=qL_1^{(\sigma)},
 & L_1^{(\sigma)}&=4\pi^2+\frac1{6\sqrt2}+\sigma .
\end{aligned}
\]
The fundamental dual determinant transport
\[
 \chi_1(\sigma,q):=e^{-\Phi(\sigma,q)}
\]
is horizontal for the canonical flat connection
\begin{equation}
 \label{p3c:eq:full-scheme-depth-determinant-connection}
 \nabla^{\rm sd}_1:=d+d\Phi
 =d+L_1^{(\sigma)}\,dq+q\,d\sigma .
\end{equation}
For every fixed $\lambda>0$, however,
\[
 \chi_\lambda(\sigma,q):=e^{-\lambda\Phi(\sigma,q)},
 \qquad
 \nabla^{\rm sd}_\lambda:=d+\lambda\,d\Phi
\]
is likewise flat.  On the contractible base $B$ it is even
connection-isomorphic to the fundamental line through multiplication by
\begin{equation}
 \label{p3c:eq:weighted-flat-gauge}
 g_\lambda(\sigma,q)
 :=e^{(\lambda-1)\Phi(\sigma,q)},
 \qquad
 \nabla^{\rm sd}_1(g_\lambda s)=g_\lambda\nabla^{\rm sd}_\lambda s .
\end{equation}
Thus flatness, topological line-bundle isomorphism, and normalization at
$q=0$ do not select $\lambda=1$.

If instead a comparison is induced before trivialization by the same
fundamental determinant-dual associated representation, with no
$q$-dependent gauge deformation or affine-origin choice, then necessarily
$\lambda=1$.  Indeed, under a scheme translation $\sigma\mapsto\sigma+c$,
the gauge in \eqref{p3c:eq:weighted-flat-gauge} changes by
\[
 \frac{g_\lambda(\sigma+c,q)}{g_\lambda(\sigma,q)}
 =e^{(\lambda-1)qc},
\]
which is independent of the depth $q$ for all $c$ only when $\lambda=1$.
Equivalently, the comparison must be induced by the identity intertwiner of
the fundamental associated representation rather than by a gauge equivalence
between two weighted flat lines.

This proposition classifies the required naturality but does not construct an
ancestry morphism from an independently constituted physical response owner.
That physical owner transition remains a separate downstream theorem object.
\end{proposition}

\begin{proof}
Since $\chi_\lambda=e^{-\lambda\Phi}$,
\(
 (d+\lambda d\Phi)\chi_\lambda=0
\), and $d(\lambda d\Phi)=0$, proving horizontality and flatness.  Moreover
$d\log g_\lambda=(\lambda-1)d\Phi$, so
\[
 (d+d\Phi)(g_\lambda s)
 =g_\lambda(d+\lambda d\Phi)s,
\]
which proves the connection isomorphism.  One has
$g_\lambda(\sigma,0)=1$ for every $\lambda$, so the upper anchor cannot
remove the counterfamily.  The displayed scheme-translation quotient proves
the affine-origin-free statement.  For $\lambda=1$, $g_1=1$ and the identity
intertwiner supplies the comparison.
\end{proof}

\begin{theorem}[Unique positive associated-line fourth root]
\label{p3c:thm:positive-associated-line-fourth-root}
Let \((\mathcal L^{\rm sd}_1,\nabla^{\rm sd}_1)\) denote the positive
fundamental determinant-dual line over \(B\) in
Proposition~\ref{p3c:prop:full-scheme-depth-determinant-connection}.  In the
category of oriented real lines with structure group \(\RR_{>0}\), there is
a unique fourth-root associated line with connection,
\[
 (\mathcal L^{\rm sd}_{1/4},\nabla^{\rm sd}_{1/4}),
 \qquad
 (\mathcal L^{\rm sd}_{1/4})^{\otimes4}\simeq
 \mathcal L^{\rm sd}_1,
 \qquad
 \nabla^{\rm sd}_{1/4}=d+\frac14\,d\Phi,
\]
whose fourth tensor power, including its connection, is the fundamental
line.  If \(H_{b_1b_0}>0\) is the horizontal transport of
\(\nabla^{\rm sd}_1\), then the root transport is uniquely
\[
 h_{b_1b_0}=H_{b_1b_0}^{1/4}
 =\exp\!\left[-\frac14\bigl(\Phi(b_1)-\Phi(b_0)\bigr)\right].
\]
\end{theorem}

\begin{proof}
The fourth-power map
\[
 p_4:\RR_{>0}\longrightarrow\RR_{>0},\qquad s\longmapsto s^4,
\]
is an isomorphism of Lie groups.  Therefore every positive transition
cocycle \(g_{ij}\) of \(\mathcal L^{\rm sd}_1\) has the unique positive
cocycle \(g_{ij}^{1/4}\); uniqueness of the positive root also preserves the
cocycle equation on triple overlaps.  The differential of \(p_4\) multiplies
connection forms by four, so the unique compatible root connection has form
one quarter of \(d\Phi\).  Equivalently, its horizontal section is
\(e^{-\Phi/4}\), whose fourth tensor power is \(e^{-\Phi}\).  A compatible
positive automorphism with fourth power one is pointwise the identity, which
proves uniqueness rather than uniqueness only up to a residual sign.
\end{proof}

\begin{corollary}[Unique positive associated-line half-root]
\label{p3c:cor:positive-associated-line-half-root}
In the notation of
Theorem~\ref{p3c:thm:positive-associated-line-fourth-root},
\[
 \mathcal L^{\rm sd}_{1/2}
 :=\bigl(\mathcal L^{\rm sd}_{1/4}\bigr)^{\otimes2}
\]
is the unique positive half-root of \(\mathcal L^{\rm sd}_1\), with
\[
 \bigl(\mathcal L^{\rm sd}_{1/2}\bigr)^{\otimes2}
 \simeq\mathcal L^{\rm sd}_1,
 \qquad
 \nabla^{\rm sd}_{1/2}=d+\frac12d\Phi,
 \qquad
 h^{(1/2)}_{b_1b_0}=H_{b_1b_0}^{1/2}.
\]
It remains internal to the same positive associated representation.
\end{corollary}

\begin{proof}
Tensoring the quarter-root cocycle and connection with themselves gives the
positive square-root cocycle and the connection one-form
\(\tfrac12d\Phi\).  Uniqueness follows because
\(s\mapsto s^2\) is an isomorphism of \(\mathbb R_{>0}\).  No cross-owner
identification is introduced.
\end{proof}

\begin{corollary}[Transport from determinant character to the operational root]
\label{p3c:cor:operational-generator-root-transport}
Suppose a typed occurrence supplies the fundamental positive character
\(H\) of \(\mathcal L^{\rm sd}_1\) as the determinant character of the
\(O(3)\)-equivariant rank-four causal generator bundle of Paper~III-A.  Then
there is no further generator-action choice: the unique identity-connected
action is
\[
                 H^{1/4}I_4,
\]
and its associated-line connection is
\[
                 \nabla=d+\frac14\,d\Phi .
\]
Its determinant is exactly \(H\), and its fourth tensor power is
\((\mathcal L^{\rm sd}_1,\nabla^{\rm sd}_1)\).
\end{corollary}

\begin{proof}
Paper~III-A proves that determinant is an isomorphism from the identity
component of the equivariant causal-cone automorphism group to
\(\mathbb R_{>0}\), with inverse \(H\mapsto H^{1/4}I_4\).  Theorem
\ref{p3c:thm:positive-associated-line-fourth-root} supplies the unique
positive associated fourth root and its one-quarter connection.  Their
uniqueness statements identify the two constructions once the typed
determinant ancestry in the hypothesis is present.
\end{proof}


\begin{corollary}[Logarithmic operational clock on a fixed determinant scheme]
\label{p3c:cor:v199-rgas-operational-clock}
On the typed occurrence of
Corollary~\ref{p3c:cor:operational-generator-root-transport}, fix a scheme
with \(A>0\), \(H(q)=e^{-Aq}\), and the hierarchy coordinate \(K=8q\).
For the energy-valued operational root \(\mu/\mu_0=H^{1/4}\),
\begin{equation}
\label{p3c:eq:v199-rgas-clock}
 \log(\mu/\mu_0)=-\frac{Aq}{4}=-\frac{AK}{32},
 \qquad
 \partial_{\log\mu}=-\frac4A\,\partial_q
                 =-\frac{32}{A}\,\partial_K .
\end{equation}
These derivatives follow this one-parameter family with other scheme
coordinates fixed.
\end{corollary}
\begin{proof}
The positive fourth-root theorem gives \(\mu/\mu_0=e^{-Aq/4}\).
Taking its logarithm and applying the chain rule proves the identities.
\end{proof}

\begin{theorem}[Reciprocal inversion transported through the positive fourth root]
\label{p3c:thm:born-fourth-root-inversion}
On the typed determinant occurrence of
Corollary~\ref{p3c:cor:operational-generator-root-transport}, a generated
Born involution propagates to every unmarked same-owner determinant descendant
by Paper~II-A, Theorem~\ref{p2a:thm:generated-born-inheritance}.  Suppose the
induced target action on the positive determinant character is identified as
\[
 H\longmapsto \frac{H_B^2}{H},
 \qquad H_B>0\ \text{fixed}.
\]
Let
\[
 \frac{\mu}{\mu_0}=H^{1/4},
 \qquad
 \frac{\mu_B}{\mu_0}=H_B^{1/4}.
\]
Then the unique positive fourth root transports the reciprocity uniquely as
\[
 \boxed{\mu\longmapsto\frac{\mu_B^2}{\mu}},
 \qquad
 \boxed{\ell:=\log(\mu/\mu_B)\longmapsto-\ell}.
\]
No extra root sign or scale-line selector survives.
\end{theorem}

\begin{proof}
Paper~II-A supplies the inherited target equivariance; the displayed
fractional-linear action on $H$ is the present target-action identification,
not a second reciprocity premise.  By
Theorem~\ref{p3c:thm:positive-associated-line-fourth-root}, positive fourth
root is a group isomorphism on the associated $\mathbb R_{>0}$ line.  Hence
\[
 \left(\frac{H_B^2}{H}\right)^{1/4}
 =\frac{H_B^{1/2}}{H^{1/4}}
 =\frac{(\mu_B/\mu_0)^2}{\mu/\mu_0}.
\]
Multiplying by $\mu_0$ gives the first formula, and taking the logarithm about
the fixed midpoint gives the second.
\end{proof}

\begin{corollary}[Weight-four response/clock relation]
\label{p3c:cor:weight-four-born-logrelation}
If a separately established same-owner response coordinate $\mathcal R$ and
the operational root satisfy
\[
 \frac{\mathcal R}{\mathcal R_B}
 =\left(\frac{\mu}{\mu_B}\right)^4,
\]
then their Born-odd logarithmic coordinates obey
\[
 \boxed{u:=\log(\mathcal R/\mathcal R_B)=4\ell}.
\]
This conclusion is conditional on the displayed same-owner relation and does
not identify an independently constituted response line with the determinant
line merely because both admit positive fourth roots.
\end{corollary}

The regulator scale \(\Lambda\), a quantum infrared cutoff \(k\), and this
operational \(\mu\) retain their owners.  On an established same-owner
matching \(k=c_\mu\mu\), with fixed \(c_\mu>0\), their logarithmic
derivatives coincide.  Connection flatness alone supplies neither that
matching nor a vanishing beta functional.  The native resolution family
\(V_\tau=P_0+e^{-\tau}P_W\), \(q=7(1-e^{-\tau})\), keeps its full UV
octet at \(\tau=0\); its geometric initial endpoint is retained when
Paper~VII-C constructs an admitted quantum RG comparison.


\begin{remark}[Same-representation firewall]
\label{p3c:rem:fourth-root-same-representation-firewall}
The theorem is internal to the positive associated representation generated
from \(\mathcal L^{\rm sd}_1\).  It neither identifies an independently
constituted response line nor the rank-four operational generator bundle with
this root.  Such a cross-role identification requires a typed physical parent
and cannot be inferred from the existence of the root or from flatness on
\(B\).
\end{remark}

Thus the structural exits are now explicit.  Either later geometry supplies a
normalized determinant trivialization strong enough to break the affine
freedom, or physical observables must be typed as horizontal/scheme-covariant
sections.  Proposition~\ref{p3c:prop:full-scheme-depth-determinant-connection}
adds that an external owner weight cannot be fixed by flatness or by a
connection isomorphism on the contractible base: the morphism must be induced
by the same associated representation before a scheme origin is chosen.  The
positive fourth root in
Theorem~\ref{p3c:thm:positive-associated-line-fourth-root} is canonical only
inside that same representation.  The source convention
$\mathrm{REG}_0$ is not promoted to a theorem by this section.




\subsection{Same-owner octahedral determinant compilation}
\begin{proposition}[Octahedral reconstruction of the Levi metric factor]
\label{p3:prop:2O-shell-levi-reconstruction}
Let \(F_O\) be the Paper-I-B octahedral shell-frame bilinear form.  On the same
round \(SU(2)\) tangent owner used for the Hopf structure,
\begin{equation}
\label{p3:eq:2O-shell-levi-reconstruction}
 \boxed{F_O|_H=\gamma_0=2g_E|_H.}
\end{equation}
Therefore the first-order conversion factor in the determinant readout is
\begin{equation}
\label{p3:eq:2O-shell-first-order-conversion}
 \boxed{\sqrt{\frac{g_E}{\gamma_0}}
 =\sqrt{\frac{g_E}{F_O|_H}}=\frac1{\sqrt2}.}
\end{equation}
This is a same-tangent-space metric identity; it does not select a
determinant finite part.
\end{proposition}

\begin{definition}[Embedded octahedral shell-to-determinant compiler]
\label{p3c:def:octahedral-shell-determinant-compiler}
Assume the generated Paper-I-B normalizer/interface role
\(\mathsf{Supp}_{2O}\) over its Hurwitz core.  Use the complete uniformly
weighted positive nearest shell \(S_{2O}\), not merely its scalar average.
Its equivariant second moment gives the shell-frame form \(F_O\), and
Proposition~\ref{p3:prop:2O-shell-levi-reconstruction} identifies
\[
 F_O|_H=\gamma_0 .
\]
On the primitive unit-Chern Hopf line with this horizontal metric, let
\(\nabla_O\) be the induced invariant connection and
\[
 L_O:=\nabla_O^*\nabla_O
\]
its primitive connection Laplacian.  The frozen compiler
\[
 \mathfrak C_O:
 (S_{2O},F_O,\text{Hopf line})
 \longmapsto
 (\mathcal L_{O,\det},\nabla_O,L_O,\det\nolimits_\zeta L_O)
\]
uses only equivariant shell averaging, horizontal restriction, formation of
the induced line connection and Laplacian, functional calculus, zeta
determinant, and the normalized first-order readout.  All operations are
fixed before the resulting scalar is evaluated.
\end{definition}

\begin{theorem}[Canonical octahedral owner intertwiner]
\label{p3c:thm:canonical-octahedral-owner-intertwiner}
The compiler of
Definition~\ref{p3c:def:octahedral-shell-determinant-compiler} is natural
under the unmarked \(SU(2)\) host action and generates the same primitive
oriented rank-one owner on its finite and determinant presentations.  Its
distinguished determinant-germ section is
\begin{equation}
\label{p3c:eq:octahedral-compiled-determinant-section}
 \delta_O
 =\frac12\,\lambda_{\operatorname{Ad}}(T_{2O})
              \lambda_{\rm def}(T_{2O})
 =\frac12\left(\frac13\right)\left(\frac1{\sqrt2}\right)
 =\frac1{6\sqrt2}
 =\omega_{2O}.
\end{equation}
Consequently there is a unique positive owner map
\[
 \mathfrak c_O:\mathcal L_{O,\rm fin}
 \longrightarrow\mathcal L_{O,\det},
 \qquad
 \mathfrak c_O(\omega_{2O})=\delta_O ,
\]
with no free multiplier.
\end{theorem}

\begin{proof}
The complete shell is invariant as a weighted set, so its second moment and
the horizontal form \(F_O|_H\) are transported equivariantly.  Formation of
the induced Hopf connection, its adjoint Laplacian, functional calculus and
zeta determinant commutes with unitary host transport.  Thus the finite shell
and determinant germ are two outputs of one frozen operator word rather than
independently owned scalars.

The compiler's declared half-determinant amplitude contributes \(1/2\); this is not an inference about the measure of a complex charged field.  The primitive reduced zeta
coefficient is
\(\zeta_{L_O}(0)=1/3=\lambda_{\operatorname{Ad}}(T_{2O})\), while the
quadratic-to-first-order metric conversion is
\(\sqrt{g_E/\gamma_0}=1/\sqrt2=\lambda_{\rm def}(T_{2O})\).
This proves
\eqref{p3c:eq:octahedral-compiled-determinant-section}.  Because the two
nonzero distinguished sections arise in one oriented line word, the positive
intertwiner carrying one to the other is unique and admits no independent
normalization.
\end{proof}

\begin{corollary}[Octahedral determinant quarter-root and bivector half-character]
\label{p3c:cor:v2184-octahedral-quarterroot-bivector-half-character}
On the generated octahedral defining block of
Theorem~\ref{p1:thm:2O-nearest-shell-tangent-frame}, the scalar defining
average is \(T_{2O}^{\rm def}=2^{-1/2}I_2\), hence its complex determinant is
\begin{equation}
\label{p3c:eq:v2184-octahedral-H-half}
 \boxed{H_{2O}:=\det_{\mathbb C}T_{2O}^{\rm def}=\frac12.}
\end{equation}
Paper~III-A, Theorem~\ref{p3a:thm:causal-cone-determinant-root-isomorphism},
therefore gives the canonical rank-four determinant-balanced representative
\begin{equation}
\label{p3c:eq:v2184-octahedral-quarterroot}
 \boxed{\rho_4(H_{2O})=H_{2O}^{1/4}I_4=2^{-1/4}I_4.}
\end{equation}
On the generated bivector representation this acts by exterior square,
\begin{equation}
\label{p3c:eq:v2184-octahedral-bivector-half-character}
 \boxed{\Lambda^2\rho_4(H_{2O})=H_{2O}^{1/2}I_6=2^{-1/2}I_6.}
\end{equation}
This is an associated-representation normalization on the already generated
rank-four owner and its bivector descendant.  It does not by itself assert an
absolute curvature-amplitude bound, a black-hole transition occurrence, or a
new action coefficient.
\end{corollary}

\begin{proof}
The first statement is the determinant of the scalar operator
\(2^{-1/2}I_2\).  Apply the rank-four determinant inverse of Paper~III-A to
\(H_{2O}=1/2\).  Finally,
\((a u)\wedge(a v)=a^2(u\wedge v)\) with \(a=2^{-1/4}\), proving the
exterior-square factor \(2^{-1/2}\).  No range statement enters the proof.
\end{proof}

\begin{theorem}[Octahedral shell-owner closure of the affine origin]
\label{p3:prop:conditional-2O-owner-REG0-closure}
Let
\[
 \delta_\sigma=\frac1{6\sqrt2}+\sigma_{\rm fin}
\]
be the fixed-shell regularization family.  On the generated \(2O\)
normalizer/interface branch,
Theorem~\ref{p3c:thm:canonical-octahedral-owner-intertwiner} identifies the
canonical Paper-I-B shell atom with the completed primitive oriented rank-one
determinant germ in the common owner, so that
\begin{equation}
\label{p3:eq:conditional-2O-owner-bridge}
 \delta_\sigma=\omega_{2O}.
\end{equation}
Then
\begin{equation}
\label{p3:eq:conditional-REG0-closure}
 \boxed{\sigma_{\rm fin}=0,\qquad\mathrm{REG}=\mathrm{REG}_0.}
\end{equation}
\end{theorem}

\begin{proof}
Paper~I-B gives \(\omega_{2O}=1/(6\sqrt2)\).  Substitution in
\eqref{p3:eq:conditional-2O-owner-bridge} forces
\(\sigma_{\rm fin}=0\).
\end{proof}

\begin{remark}[Anchored and unanchored REG branches]
The \((\mathrm{REG})\) family is an affine torsor, and scheme-translation
symmetry has no distinguished origin.  A scheme-neutral finite-shell number
therefore cannot select \(\mathrm{REG}_0\) by numerical equality alone.
Theorem~\ref{p3c:thm:canonical-octahedral-owner-intertwiner} supplies the
stronger same-operator origin precisely when the generated octahedral
interface is present, and
Theorem~\ref{p3:prop:conditional-2O-owner-REG0-closure} then selects
\(\mathrm{REG}_0\).  A branch without that interface anchor retains the full
affine torsor and the no-natural-origin proposition remains unchanged.
\end{remark}



\subsubsection{The occurring Hurwitz normalizer fixes the REG origin}

\begin{theorem}[Occurring Hurwitz REG origin]
\label{p3c:thm:occurring-hurwitz-reg-origin}
On the occurring Hurwitz/Hopf branch, use the same embedded
\(H=2T\subset SU(2)\) in the normalizer and the canonical octahedral
shell determinant compiler.  Then
\[
 N_{SU(2)}(H)=2O,\qquad
 \lambda_{\mathrm{def}}=1/\sqrt2,\quad
 \lambda_{\mathrm{ad}}=1/3,\quad
 \delta_{\mathcal O}=\omega_{2O}=1/(6\sqrt2).
\]
The same-owner equality
\(\delta_\sigma=\omega_{2O}\) fixes \(\sigma_{\mathrm{fin}}=0\).
Thus the anchored branch uses \(\mathrm{REG}_0\), with
\[
 \alpha_0=\exp(4\pi^2+1/(6\sqrt2)),\qquad
 E_K/E_J=\alpha_0^{-(K-J)/32}.
\]
\end{theorem}
\begin{proof}
Apply Paper I-B's exceptional-normalizer theorem to the same
occurring embedded Hurwitz group.  Normalization is a constructor
on this input, so it does not require an independently occurring
second group.  The generated normalizer is the \(2O\) input to the
already defined six-point positive uniform nearest-shell compiler.
Its defining and adjoint averages are the displayed scalars;
the compiler contracts them to \(1/(6\sqrt2)\).
The canonical octahedral owner intertwiner and the conditional
REG-origin proposition therefore apply with their generated-owner
antecedent discharged.  In their unchanged normalization the
affine finite-part displacement is zero exactly when the two
defects agree.  Substitution into the hierarchy character gives
the two final formulas.
\end{proof}

The no-natural-origin proposition still holds on the unanchored
affine scheme torsor.  A scheme translation by \(c\) multiplies a
weight-\(w\) readout by \(e^{-wc}\); only total-weight-zero
combinations descend as invariant scalar coordinates without the anchor.
Nonzero-weight associated sections may still be horizontal for their
associated connections, as the exact scheme-continuation theorem proves.  On the
anchored branch the stabilizer of its distinguished origin under
translation is \(\{0\}\).  This is a scheme-coordinate statement,
not a physical-time Noether law.  Neither anchoring nor horizontal
descent fixes dimensional calibration, particle assignment, the
running/matching map or a physical pole.

\subsection{First anisotropy, physical support, and the REG-origin incidence test}
\label{p3c:sec:v1188-reg-incidence}

The following consequences use the existing Paper-I-B exceptional-frame
classification and Paper-II-A finite-operation comparison.  They do not
change the value or the claimed scope of
Theorems~\ref{p3:thm:delta_rigorous},
\ref{p3:thm:D1-primitive-smooth-uniqueness}, or
\ref{p3c:thm:occurring-hurwitz-reg-origin}.
In particular, the rank-one $(\mathrm{REG})$ value
$\delta_{\rm loc}=1/(6\sqrt2)$ is already derived; the question here is
whether its fixed-shell origin is identified on an \emph{actual} physical
operator owner.

\begin{proposition}[Exact octahedral stabilizer of the first local Born-even quartic]
\label{p3c:prop:v1188-quartic-octahedral-stabilizer}
On $\operatorname{Im}\mathbb H\cong\mathbb R^3$, let
$r^2=x^2+y^2+z^2$ and
\[
  Q_4(x,y,z)=x^4+y^4+z^4-\tfrac35 r^4.
\]
Then $Q_4$ is harmonic and Born-even. Its full
orientation-preserving rotational stabilizer is the octahedral group
$O\subset SO(3)$; its inverse image in $\operatorname{Spin}(3)=SU(2)$
is $2O$.  This identifies the symmetry of the \emph{quartic response},
not the support group of an operation producing that response.
\end{proposition}
\begin{proof}
In three Euclidean dimensions,
$\Delta(x^4+y^4+z^4)=12r^2$ and
$\Delta r^4=20r^2$, so $\Delta Q_4=0$.
On the unit sphere,
$Q_4=\sum_i x_i^4-3/5$ has its six global maxima precisely at
$\{\pm e_1,\pm e_2,\pm e_3\}$.  Every symmetry of $Q_4$ therefore
permutes those six axes, hence belongs to $O$; conversely every
orientation-preserving signed permutation of the axes preserves $Q_4$.
The full stabilizer is $O$, whose spin preimage is $2O$.
The binary tetrahedral subgroup $2T\subset2O$ also preserves $Q_4$,
so invariance of a quartic alone does not establish that an actually
performed finite comparison has operation support $2O$.
\end{proof}

\begin{proposition}[Five-frame cancellation and the Petersen sextic]
\label{p3c:prop:v1188-fiveframe-quartic-cancellation}
Inside a \emph{mathematically fixed} $2I$ completion, let the five
conjugate tetrahedral cores and local octahedral frames be indexed by
$A_5/A_4$.  For coset representatives $R_i$, set
$Q_i(v)=Q_4(R_i^{-1}v)$, $i=1,\ldots,5$.  Then
\[
  \boxed{\sum_{i=1}^5Q_i(v)=0},\qquad
  \boxed{B_8(v):=\sum_{i=1}^5Q_i(v)^2
      =\tfrac{16}{105}|v|^8+\tfrac9{10}|v|^2H_6^P(v)},
\]
where $H_6^P$ is the aligned and normalized ten-axis Petersen
sextic of Paper~III-B,
Theorem~\ref{p3b:thm:v319-petersen-moment-six}.
The $Q_i$ span the real standard $A_5$ module $\mathbf4$.
Consequently, an $A_5$-invariant actual source and state cannot have
any nonzero \emph{linear} expectation in this $\mathbf4$ module, even
though the invariant positive response $B_8$ may be nonconstant.
\end{proposition}
\begin{proof}
The sum is an $A_5$-invariant harmonic quartic, and the
icosahedral rotation group has no nonzero invariant harmonic of
degree four.  Its vanishing also shows that the five equal-norm
translates span the nontrivial four-dimensional standard module.
The nonnegative degree-eight polynomial $B_8$ is $A_5$-invariant;
its only admissible harmonic components at degrees $0,2,4,6,8$
are the scalar and the unique $\ell=6$ icosahedral harmonic.  Thus
$B_8=a|v|^8+b|v|^2H_6^P(v)$.  The angular mean is
$a=5\langle Q_4^2\rangle=5(16/525)=16/105$.
At a unit icosahedral vertex $v_*$, every $Q_i(v_*)=0$ and
$H_6^P(v_*)=-32/189$, in the same source normalization. Hence
$0=16/105-(32/189)b$, giving $b=9/10$.
Finally an invariant linear expectation vector in an irreducible
nontrivial $A_5$ module is zero, by averaging over $A_5$.
This conclusion concerns an invariant source/state; a marked source
or a sourced symmetry-breaking sector is not excluded.
\end{proof}

\begin{corollary}[Raw Levi-frame uniqueness is not a physical occurrence certificate]
\label{p3c:cor:v1188-raw-levi-frame-conditional}
For the \emph{unnormalized}, positive shell-frame metrics already
classified in Paper~I-B,
Theorem~\ref{p1b:thm:v11-three-shell-frames},
\[
 F_{2T}=\tfrac83g_E,\qquad F_{2O}=2g_E,\qquad F_{2I}=4g_E,
 \qquad \gamma_0=2g_E|_H,
\]
any nonnegative convex mixture satisfying
$\sum_\Gamma w_\Gamma F_\Gamma|_H=\gamma_0$ has
$w_{2O}=1$.  This selects $2O$ \emph{within that raw-frame
comparison role}.  It is not implied by selecting an arbitrary
normalized second-moment readout; individually trace-normalized
metrics all equal $2g_E$.
\end{corollary}
\begin{proof}
The raw comparison gives
$\tfrac83w_{2T}+2w_{2O}+4w_{2I}=2$ together with
$\sum w_\Gamma=1$.  Thus
$\tfrac23w_{2T}+2w_{2I}=0$.  Nonnegativity forces
$w_{2T}=w_{2I}=0$, and $w_{2O}=1$.
Paper~I-B computes the trace-normalized equality, showing why the
raw/normalized distinction is essential.
\end{proof}

\begin{proposition}[Born heat trace versus integrated determinant shell]
\label{p3c:prop:v1188-heat-determinant-distinction}
Let $L$ be the positive primitive charge-one horizontal Hopf
Laplacian of the existing $(\mathrm{REG})$ calculation, and write
\[
 H(t)=\operatorname{Tr}e^{-tL}=\frac{A}{t}+B+O(t),
 \qquad A=\tfrac14,\quad B=\tfrac13.
\]
For $a>0$, $0<h<1$, define the \emph{different} comparison
functionals
\[
 T_\varepsilon=a\bigl(H(h\varepsilon)-H(\varepsilon)\bigr),
 \qquad
 I_\varepsilon=a\int_\varepsilon^\infty
         \frac{H(ht)-H(t)}{t}\,dt .
\]
Subtracting their respective common leading power term
$aA(h^{-1}-1)/\varepsilon$, their finite parts obey
\[
 \boxed{\operatorname{FP}T_\varepsilon=0},\qquad
 \boxed{\operatorname{FP}I_\varepsilon=aB\log(1/h)}.
\]
For the normalized Born coefficient $a=1/2$ and $h=1/2$,
these are $0$ and $\log2/6$, respectively.  Hence the
half-weighted \emph{unintegrated} heat-trace action of
Definition~\ref{p3:def:physical_spectral_action} does not, by
its $1/2$ prefactor alone, establish physical performance of the
proper-time-integrated determinant readout.
\end{proposition}
\begin{proof}
The $B$ coefficients cancel in $T_\varepsilon$, leaving the
stated power term plus $O(\varepsilon)$.  Positivity of $L$ makes
the integrated comparison well-defined at large times; changing
variables in its two terms gives exactly
\[
  I_\varepsilon=a\int_{h\varepsilon}^\varepsilon\frac{H(s)}s\,ds.
\]
Inserting the heat expansion produces the same power term and
the finite term $aB\log(1/h)$, with remainder $O(\varepsilon)$.
The actual action may admit additional derived output functionals,
but the displayed bare trace does not imply this integration.
\end{proof}

\begin{corollary}[Sufficient conditional discharge of the REG-origin incidence]
\label{p3c:cor:v1188-conditional-reg-incidence}
Suppose the \emph{actual constituted physical action} generates an
occurring finite identity-rooted multiplication-closed comparison
$Z\subset SU(2)$ exact on the entire
$\mathcal P_2=E_0\oplus E_1\oplus E_2$ band, with the role and
subgroup-restriction fidelity prescribed in Paper~II-A,
Theorems~\ref{p2a:thm:v233-quadratic-support-discharge} and
\ref{p2a:thm:v236-qop-exceptional-frontier}.  Suppose further
that the resulting normalizer shell is consumed by the \emph{same
actual primitive Hopf determinant owner} through
Definition~\ref{p3c:def:octahedral-shell-determinant-compiler},
with the finite-part identification of
Theorem~\ref{p3c:thm:occurring-hurwitz-reg-origin}.
Then the GP3-least unmarked comparison type is $[2T]$, its
generated normalizer type is $[2O]$, and that performed
same-owner comparison has $\sigma_{\rm fin}=0$, i.e.
$\mathrm{REG}_0$.
\end{corollary}
\begin{proof}
The hypotheses discharge the \emph{occurrence}, band-image and
subgroup-restriction antecedents of the cited Paper-II-A
$\mathsf{QOp}_2$ theorem, which selects $[2T]$.  The already
proved exceptional normalizer identifies $N_{SU(2)}(2T)=2O$;
normalization requires no independent occurrence of $2O$.
The additional physical same-owner hypothesis permits direct
application of
Theorem~\ref{p3c:thm:occurring-hurwitz-reg-origin},
which fixes $\sigma_{\rm fin}=0$.  No part of this proof establishes
the extra physical-incidence hypotheses themselves.
\end{proof}

\begin{remark}[Open physical interface; no global REG promotion]
\label{p3c:rem:v1188-incidence-status}
Neither the existence of a mathematical $2T$ sampler, the
$SO(3)$ stabilizer of $Q_4$, the $A_5$-invariant scalar $B_8$,
nor the raw-frame convex uniqueness proves that the native action
performs the required finite $\mathsf{QOp}_2$ role or sends its
normalizer shell into its own completed primitive determinant.
The existing $\mathrm{REG}_0$ result remains proved on its
\emph{occurring Hurwitz/Hopf same-owner branch}; outside that
branch the fixed-shell affine scheme torsor is retained.
A future physical discharge must exhibit a generated point-operation
and a source-preserving map to the performed determinant readout.
The numerical identity $\delta_{\rm loc}=\omega_{2O}$ and the
protected rank chain require no revision.
\end{remark}




\subsection{Fixed-charge determinant evaluation and physical kinetic completion}
\label{p3c:sec:v2216-fixed-charge-completion}

The primitive horizontal comparison has an exact absolute determinant as
well as its already derived scale anomaly.  These are different readouts.
Moreover, restoring the vertical and temporal kinetics changes the operator
whose anomaly is being evaluated.  The following results make both
distinctions explicit without selecting a matter field, measure, physical
clock or additional action.

\begin{theorem}[Fixed-charge horizontal heat coefficients]
\label{p3c:thm:v2216-all-charge-zeta}
Let $L_q$ be the dimensionless pure connection Laplacian on the integer
Hopf line of charge $q$ over the round base of radius $1/2$.  Put
$m=|q|$.  Its spectrum and complex multiplicities are
\[
 \lambda_{r,q}=4r^2+4(m+1)r+2m,\qquad
 d_{r,q}=2r+m+1,\qquad r=0,1,\ldots .
\]
For every fixed $q$,
\begin{equation}
 \Tr e^{-tL_q}=\frac1{4t}+\frac13+
       \left(\frac4{15}-\frac{q^2}{6}\right)t+O(t^2).
 \label{p3c:eq:v2216-charge-heat}
\end{equation}
Consequently $\zeta_{L_q}(0)=1/3$ for $q\ne0$, whereas exclusion of
the single constant zero mode at $q=0$ gives
$\zeta_{L_0,\mathrm{red}}(0)=-2/3$.
\end{theorem}
\begin{proof}
In the inherited fixed-charge resolution set $k=m+2r$.  Subtracting the
vertical eigenvalue $q^2$ from $k(k+2)$ gives the displayed spectrum,
with multiplicity $k+1$.  With $b=(m+1)/2$ the exact heat trace is
\[
 2e^{(m^2+1)t}\sum_{r=0}^{\infty}(r+b)e^{-4t(r+b)^2}.
\]
Euler--Maclaurin summation gives for the last sum, including its factor
$2$, the expansion $1/(4t)+2\zeta(-1,b)-8t\zeta(-3,b)+O(t^2)$.
Substitution of $\zeta(-1,b)=-B_2(b)/2$ and
$\zeta(-3,b)=-B_4(b)/4$ proves
\eqref{p3c:eq:v2216-charge-heat}.  Equivalently its constant is
$(24\pi)^{-1}\int_{S^2}\operatorname{Scal}=1/3$.
For $q\ne0$ the lowest eigenvalue is $2|q|>0$; for $q=0$ the
kernel is the scalar constant.  The Mellin heat formula gives
$\zeta(0)=a_1-\dim\ker L_q$, proving the assertions.  The expansion
is at fixed charge and is not asserted uniformly as $|q|\to\infty$.
\end{proof}

\begin{corollary}[Charge-independent anomaly with charge-dependent spectra]
\label{p3c:cor:v2216-all-charge-anomaly}
For every $q\ne0$, the Levi-normalized operator is $L_{q,\gamma}=L_q/2$.
For the declared formal amplitude $Z_b(L)=(\det_\zeta L)^{-b}$,
$b>0$, one has
\[
 \log Z_b(L_q/2)-\log Z_b(L_q)=\frac{b\log2}{3},\qquad
 \delta_{q,b}=\frac{b}{3\sqrt2}.
\]
Thus the formal half-amplitude gives the established
$1/(6\sqrt2)$ for every nonzero charge.  The ordinary complex Gaussian,
if actually licensed in this complex spectral trace, instead has $b=1$.
Neither anomaly equality nor the primitive comparison selects charged
matter or identifies distinct line bundles.
\end{corollary}
\begin{proof}
Use $\log\det_\zeta(cL)=\log\det_\zeta L+\zeta_L(0)\log c$,
then the existing first-order Levi-to-round factor $1/\sqrt2$.
The charge-sensitive coefficient in \eqref{p3c:eq:v2216-charge-heat}
also shows directly that anomaly equality is not equality of spectra.
For a dimensionful radius $a$, introduce the actual reference inverse
length $\mu$: the dimensionless determinant is that of
$L_q/(\mu a)^2$.  Its logarithm differs by
$-2\zeta_{L_q}(0)\log(\mu a)$, which cancels in the same-radius
Levi/round comparison.
\end{proof}

\begin{proposition}[Convergent primitive determinant and certified tail]
\label{p3c:prop:v2216-convergent-determinant}
For $L_H$ with spectrum $4j^2-2$ and multiplicity $2j$, $j\ge1$,
\begin{align}
 \log\det_\zeta L_H
 &=\frac{\log4}{3}-4\zeta_{\mathbb R}'(-1)-\gamma_{\mathrm E}
   -2\sum_{r=2}^{\infty}\frac{2^{-r}}{r}
                              \zeta_{\mathbb R}(2r-1),
 \label{p3c:eq:v2216-primitive-determinant}\\
 \left|\text{tail after }r=N\right|
 &\le\frac{2\zeta_{\mathbb R}(3)}{N+1}\,2^{-N},\qquad N\ge2.
 \nonumber
\end{align}
In particular $\log\det_\zeta L_H=0.1063922291873472810\ldots$.
This absolute analytic readout is distinct from the normalized anomaly
$1/(6\sqrt2)$ and supplies no independent physical subtraction rule.
\end{proposition}
\begin{proof}
Initially for $\operatorname{Re}s>1$, the binomial expansion gives
\[
 \zeta_{L_H}(s)=2\,4^{-s}\sum_{r=0}^{\infty}
   \frac{(s)_r}{r!}\,2^{-r}\zeta_{\mathbb R}(2s+2r-1).
\]
The $r\ge2$ sum and its derivative converge uniformly near zero.
The $r=1$ term is
$4^{-s}s\zeta_{\mathbb R}(1+2s)
=1/2+s(\gamma_{\mathrm E}-\log4/2)+O(s^2)$.
Using $\zeta_{\mathbb R}(-1)=-1/12$ and
$\partial_s[(s)_r/r!]_{s=0}=1/r$ proves
\eqref{p3c:eq:v2216-primitive-determinant} by
$\log\det_\zeta L_H=-\zeta_{L_H}'(0)$.
In the tail, bound $\zeta_{\mathbb R}(2r-1)$ by
$\zeta_{\mathbb R}(3)$ and $1/r$ by $1/(N+1)$ and sum the geometric
series.  This is an evaluation in the canonical zeta prescription;
adding a finite action/readout counterterm remains a change of data.
\end{proof}

\begin{theorem}[Round kinetic completion of the primitive Hopf channel]
\label{p3c:thm:v2216-round-hopf-completion}
On a supplied round physical slice $S^3_a$, use the native integer
generator $Q=-2i\partial_\psi$.  Then
\[
 \Delta_{S^3_a}=\Delta_H+a^{-2}Q^2.
\]
For either $q=+1$ or $q=-1$, $k=2j-1$, and
\begin{equation}
 \lambda_{H,j}=\frac{4j^2-2}{a^2},\qquad
 \lambda_{\mathrm{vert},j}=\frac1{a^2},\qquad
 \lambda_{\mathrm{sp},j}=\frac{4j^2-1}{a^2}.
 \label{p3c:eq:v2216-physical-hopf-spectrum}
\end{equation}
Each single charge has $2j$ complex modes.  A real scalar on the paired
$q=\pm1$ sector has $4j$ real coordinates, as does a complex scalar
on one charge sector.  Counting the two charges as independent complex
fields instead doubles this determinant multiplicity.
\end{theorem}
\begin{proof}
The period-$2\pi$ fibre parameter is $\alpha=\psi/2$; its generator
$\partial_\alpha$ has length $a$, so the positive vertical operator
is $-a^{-2}\partial_\alpha^2=a^{-2}Q^2$.  The orthogonal splitting
and inherited scalar spectrum $k(k+2)$ prove
\eqref{p3c:eq:v2216-physical-hopf-spectrum}.  Fixed $(k,q)$ has
multiplicity $k+1=2j$.  Complex conjugation exchanges the two charges,
so one complex amplitude supplies two real coordinates in the paired
real sector.  This proves the count.  No common constant rescaling
identifies $4j^2-2$ with $4j^2-1$ at every $j$.
\end{proof}

\begin{proposition}[Owner-specific completed zeta values]
\label{p3c:prop:v2216-completed-zeta}
In the complex fixed-charge trace with multiplicity $2j$, let
$L_{\mathrm{sp}}=L_H+I$ and $L_{\mathrm{conf}}=L_H+2I$.  Then
\[
 \zeta_{L_H}(0)=\frac13,\qquad
 \zeta_{L_{\mathrm{sp}}}(0)=\frac1{12},\qquad
 \zeta_{L_{\mathrm{conf}}}(0)=-\frac16.
\]
On a static Euclidean product with periodic time of length $\beta>0$,
the positive minimal fixed-charge operator
$P=-\partial_\tau^2+a^{-2}L_{\mathrm{sp}}+m^2$, $m^2\ge0$,
has $\zeta_P(0)=0$.  These statements concern different complete
operators and do not change the primitive horizontal compiler.
\end{proposition}
\begin{proof}
The horizontal heat trace is $1/(4t)+1/3+O(t)$.
Multiplication by $e^{-ct}$ changes its constant to $1/3-c/4$.
All three operators are positive, so this constant equals $\zeta(0)$.
For the periodic clock, Poisson summation gives
$\Tr e^{t\partial_\tau^2}=\beta/\sqrt{4\pi t}$ up to exponentially
small terms.  The fixed-charge spatial operator is a Laplace-type
operator on the closed two-dimensional Hopf base and has heat powers
$t^{-1+n}$, $n\ge0$.  The product has powers $t^{-3/2+n}$ and no
constant; the spatial gap excludes a kernel.  The Mellin formula proves
$\zeta_P(0)=0$.  This projected problem is on $S^1\times S^2$;
summing all charges reconstructs the four-dimensional parent and has
different heat coefficients.  Other time boundaries or field content
require their own complete operator and measure.
\end{proof}

The temporal completion is an action pullback, proved in Paper~III-F,
Theorem~\ref{p3f:thm:v2216-physical-hopf-mode}.  Its canonical scalar
frequency is $4j^2+a^2m^2+(\xi-1/6)a^2R_g$.  The native odd scalar
has $\xi=1/6$ and its inherited potential Hessian; this does not
constitute an occupied charge-one field or a positive Gaussian by itself.
The fixed horizontal comparison, the unital scalar heat process and an
actual odd-field source therefore keep their own bundles, action domains
and normalizations.  The Hodge scalar zero line is uncharged and cannot
be identified with the gapped primitive charged ladder.


\section{Clifford transitions, rooted boundary, and response typing}
\label{paper3c:sec:organized-response}
\subsection{Clifford transitions and the balanced hierarchy response}
\label{p3c:sec:v179-clifford-transition-response}

The active support theorem supplies the inputs of the following constructor:
the even quaternionic carrier, its scalar line, and the transported Hopf axis.
All operations below take place in the already represented finite operator
envelope.  GP1 supplies its finite domain; rooted generated-image
co-realization (Paper II-A,
Theorem~\ref{p2a:thm:rooted-generated-image-corealization}) retains the
complete source and output family.  GP3 leastness is applied to the specified
completion problem, with its inputs fixed.

\begin{theorem}[Multiplication source and Hopf-projector tangent]
\label{p3c:thm:v179-clifford-projector-tangent}
Use the real trace metric on \(A=\Cl(0,3)\) and the active projectors of
Theorem~\ref{p3:thm:derived-active-support}.  In
\(A^{\rm even}\simeq\mathbb H\), let \(u\) be a unit imaginary Hopf axis,
\(V=u^\perp\cap\operatorname{Im}\mathbb H\),
\(Q_4=|u\rangle\langle u|\), \(Q_5=P_V\), and
\(Q_7=P_0=|1\rangle\langle1|\).  Extend these by zero on the odd complement
and put \(Q_0=I-\Pi_4\).  For \(v\in V\), define
\[
 T_u(v)=Q_5L_vQ_4,\qquad S_u(v)=T_u(v)+T_u(v)^\dagger,
 \qquad J_v=\tfrac12(L_v-R_v).
\]
Then
\begin{equation}
\label{p3c:eq:v179-transition-and-tangent}
 T_u(v)=|vu\rangle\langle u|,\qquad
 T_u(v)^\dagger T_u(w)=\langle v,w\rangle Q_4,\qquad
 S_u(v)=[J_v,Q_4].
\end{equation}
The source \(T_u\) is an isomorphism onto
\(\operatorname{Hom}_{\mathbb R}(\mathbb Ru,V)\).  The operator \(S_u(v)\)
is the unique self-adjoint completion retaining that prescribed off-diagonal
block and having zero other blocks.  With the protected whole-octet trace
\(\tau_8=\operatorname{Tr}_{\mathbb R}/8\),
\begin{equation}
\label{p3c:eq:v179-transition-metric}
 \tau_8(S_u(v)^\dagger S_u(w))=\tfrac14\langle v,w\rangle .
\end{equation}
Conjugation transports \(u,v\), the rooted Hurwitz datum, and these operators
together.  This is covariance of that family.
\end{theorem}
\begin{proof}
The active projectors and quaternionic realization above identify these
maps with Paper~1B, Proposition~\ref{p1b:prop:block-edges}~\cite{Kugelmann2026Paper1B}.
It gives Eqs.~\eqref{p3c:eq:v179-transition-and-tangent}
and~\eqref{p3c:eq:v179-transition-metric} in the same coefficient metric.
Right multiplication by \(u\) maps \(V\) isometrically onto \(V\), so
\(T_u\) has the asserted full two-dimensional image.  Self-adjointness
determines the reverse block from the prescribed block, and the zero-block
condition determines the remaining blocks.  Simultaneous conjugation
transports the active projectors, multiplication and adjoint together,
proving covariance of the rooted family.
\end{proof}

\begin{theorem}[Positive balanced response and its dynamical conversion]
\label{p3c:thm:v179-balanced-hierarchy-response}
On a positive quarter-character branch of
Theorem~\ref{p3c:thm:positive-associated-line-fourth-root}, choose its
positive reference \(E_0\), put \(E_K=E_0\alpha_{\rm hier}^{-K/32}\), and
represent that character on the active depths by
\[
 \mathcal Q=0Q_0+4Q_4+5Q_5+7Q_7,\qquad
 H_{\rm car}=E_0\alpha_{\rm hier}^{-\mathcal Q/4}
 =E_0Q_0+E_{32}Q_4+E_{40}Q_5+E_{56}Q_7 .
\]
Here \(H_{\rm car}\) is a carrier grading; a physical energy interpretation
uses the same-owner character ancestry and its unit map.  The positive
superoperator \(X\mapsto H_{\rm car}XH_{\rm car}\) has the unique positive
square root
\[
 \mathscr E_H(X)=H_{\rm car}^{1/2}XH_{\rm car}^{1/2}.
\]
It satisfies
\begin{equation}
\label{p3c:eq:v179-e36-response}
 \mathscr E_H(S_u(v))=E_{36}S_u(v),\qquad
 E_{36}=\sqrt{E_{32}E_{40}},\qquad
 |\mathscr E_H\circ S_u|=\tfrac12E_{36}I_V .
\end{equation}
The last magnitude uses exactly the source metric in
\eqref{p3c:eq:v179-transition-metric}.  The actual orbit derivative is
instead the explicitly related operator
\begin{equation}
\label{p3c:eq:v179-carrier-orbit-conversion}
\begin{aligned}
 [J_v,H_{\rm car}]
 &=(E_{32}-E_{40})S_u(v)\\
 &=(\alpha_{\rm hier}^{1/8}-\alpha_{\rm hier}^{-1/8})
       \mathscr E_H(S_u(v)).
\end{aligned}
\end{equation}
\end{theorem}
\begin{proof}
Apply Paper~1B, Proposition~\ref{p1b:prop:block-edges}~\cite{Kugelmann2026Paper1B},
to the positive grading \(H_{\rm car}\).  The positive quarter-character
formula of Proposition~\ref{p1b:prop:trace-coordinate}, with the
same-owner ancestry established above, gives
\(h_i=E_0\alpha_{\rm hier}^{-q_i/4}\).
The \((4,5)\) and \((5,4)\) blocks therefore have weight
\(\sqrt{E_{32}E_{40}}=E_{36}\); the source metric
\eqref{p3c:eq:v179-transition-metric} gives the factor \(1/2\) in
\eqref{p3c:eq:v179-e36-response}.  The commutator formula in the same
block proposition is \([J_v,H_{\rm car}]=(E_{32}-E_{40})S_u(v)\).
Dividing by \(\sqrt{E_{32}E_{40}}\) gives the second line of
\eqref{p3c:eq:v179-carrier-orbit-conversion}.
\end{proof}


\begin{corollary}[Canonical Petersen contraction into the protected $E_{36}$ source]
\label{p3c:cor:v371-petersen-e36-comparison}
Consume Paper~I-B, Proposition~\ref{p1:prop:v371-petersen-quadrupole-flux},
and the normalized source of
Theorem~\ref{p3c:thm:v179-clifford-projector-tangent}.
For a fixed projective $C_6$ Hopf axis $u$, let
$P_u=I_3-uu^{\mathsf T}$ on $\operatorname{Im}\mathbb H$ and let
$K:\mathbb R^{15}_{\rm oriented}\to\mathbb R^3$ be Paper~I-B's
quaternionic flux map. The geometrically defined \emph{comparison}
operator
\[
 \mathcal J_u=\mathscr E_H\circ S_u\circ P_u\circ c\circ\Phi
 =-{2\over3}E_{36}\,S_u\circ P_u\circ K
\]
uses only the existing off-diagonal active $Q_4$--$Q_5$ block.
Its source-metric Gram satisfies
\[
 \mathcal J_u^\dagger\mathcal J_u
 ={E_{36}^2\over9}K^\dagger P_uK,
 \qquad\operatorname{Spec}(\mathcal J_u^\dagger\mathcal J_u)
 =\{(5E_{36}^2/9)^{(2)},0^{(13)}\},
\]
\[
 \sum_{e=1}^{15}\|\mathcal J_u(e)\|_{\tau_8}^2
 ={10\over9}E_{36}^2.
\]
Thus fifteen geometric input labels give only two independent
fixed-axis tangent amplitudes, \emph{not} fifteen particle states.
For a mathematical direct sum of the ten comparison fibres, the
locally generated trimer source has Gram
$E_{36}^2(3I-B^\dagger B/2)/9$ and kernel the uniform edge line.
Neither comparison is a same-action physical interaction vertex;
$E_{36}$ remains a protected \emph{operator-transition} weight, not
a decay width or a newly inserted carrier state.
\end{corollary}
\begin{proof}
The tensor contraction in Paper~I-B is
$c\Phi=-2K/3$. The positive balanced response multiplies the
$Q_4$--$Q_5$ source by $E_{36}$, while its actual normalized trace
metric gives $\|S_u(v)\|_{\tau_8}^2=\|v\|^2/4$.
The nonzero eigenvalues of $K^\dagger P_uK$ equal those of
$P_uKK^\dagger P_u=5P_u$, so the displayed Gram has rank two and
the indicated spectrum and trace. The global trimer identity is
Paper~I-B, Theorem~\ref{p1:thm:v371-petersen-tangent-exact-sequence}.
\end{proof}

\begin{corollary}[Complete even-operator weights and retained action types]
\label{p3c:cor:v179-complete-balanced-weights}
On the even carrier, the multiplication words
\(L_aP_0L_b^\dagger=|a\rangle\langle b|\) span
\(\operatorname{End}_{\mathbb R}(A^{\rm even})\).  The balanced weights of
its blocks are
\[
\begin{array}{c|ccc}
\text{output / input}&\mathbb Ru&V&\mathbb R1\\ \hline
\mathbb Ru&E_{32}&E_{36}&E_{44}\\
V&E_{36}&E_{40}&E_{48}\\
\mathbb R1&E_{44}&E_{48}&E_{56}
\end{array}
\]
with real multiplicities \(1,4,4,2,4,1\) at
\(K=32,36,40,44,48,56\), respectively.
The balanced depth \(\tfrac12(\mathcal QX+X\mathcal Q)\) has value \(9/2\)
on \(T_u(V)\), whereas \([\mathcal Q,T_u(v)]=T_u(v)\).
Thus \(E_{36}\) is an operator-transition weight; the active projector
spectrum remains \(q=0,4,5,7\).
\end{corollary}
\begin{proof}
The displayed multiplication word is a rank-one operator, giving the full
endomorphism envelope.  Paper~1B,
Proposition~\ref{p1b:prop:block-edges}~\cite{Kugelmann2026Paper1B},
assigns each block the weight \(\sqrt{h_ih_j}\).
Substitute the active depths \((4,5,7)\), dimensions \((1,2,1)\) and
\(h_i=E_0\alpha_{\rm hier}^{-q_i/4}\): the half-sum depths give the
table, and the products of dimensions give the listed multiplicities.
On \(T_u(V)\) the half-sum and difference of the depths are
\((4+5)/2=9/2\) and \(5-4=1\), respectively.
\end{proof}

The normalized geometric source is \(2S_u\).  If a licensed quadratic
action uses this same source with \(Z=\tfrac14I_V\) and
\(K=\tfrac14E_{36}^2I_V\), its normalized operator is \(E_{36}^2I_V\).
Physical use must transport both forms by
Theorem~\ref{p3f:thm:v179-normalized-scalar-hessian}; the raw half-factor
alone is not a fermion mass.  Nor does host covariance constitute a
continuous symmetry of a fixed finite cell.  Noether applies to the complete
action with its retained marks
(Theorem~\ref{p3a:thm:v149-native-noether-frontier}).

\subsection{Hopf--Clifford scale character of the protected hierarchy}
\label{p3c:sec:v222-hopf-clifford-character}

The positive determinant/rank character established in this paper has a
natural \emph{unitary character companion}.  The companion is a function on
logarithmic scale space, not a Reeb standing wave, a shell selection rule,
or a new action density.  This distinction is essential: the earlier
unsupported seven-rung Reeb source-phase selector is not reinstated.

\begin{definition}[Primitive and completed logarithmic scale circles]
\label{p3c:def:v222-log-scale-circles}
On the same-owner protected rank branch with upper normalization $U>0$,
let
\[
  L_{\rm top}:=\Vlogcap=4\pi^2,\qquad
  \alpha_{\rm top}:=e^{L_{\rm top}}.
\]
If the $(\mathrm{REG})$ fixed-shell completion is admitted, put
\[
  L_\sigma:=L_{\rm top}+\delta_{\rm loc},\qquad
  \delta_{\rm loc}=\frac1{6\sqrt2},\qquad
  \alpha_\sigma=e^{L_\sigma}.
\]
For either choice $L\in\{L_{\rm top},L_\sigma\}$ define
\[
  q_L(\scal):=\frac{\log(U/\scal)}{L},\qquad
  \mathsf T_L:=\mathbb R/L\mathbb Z.
\]
The quotient $\mathsf T_L$ means that scales differing by the
multiplicative factor $e^L$ have the same phase class.  It is a
constructed logarithmic-scale circle; no identification with the
underlying Hopf fiber or with its connection holonomy is assumed.
\end{definition}

\begin{theorem}[Fundamental unitary companion of the Hopf rank character]
\label{p3c:thm:v222-rank-character-unitary-companion}
Under the hypotheses of Definition~\ref{p3c:def:v222-log-scale-circles},
the continuous unitary characters of $\mathsf T_L$ with $\chi(U)=1$ are
\begin{equation}
\label{p3c:eq:v222-unitary-character-family}
 \chi_{L,m}(\scal)
  :=\exp\!\left(\frac{2\pi i m}{L}\log\frac{\scal}{U}\right)
  =e^{-2\pi i m q_L(\scal)},\qquad m\in\mathbb Z.
\end{equation}
The primitive faithful characters are $m=\pm1$, exchanged by complex
conjugation.  With the orientation $m=1$, the associated real scale
character is
\[
 F_L(\scal):=\operatorname{Re}\chi_{L,1}(\scal)
  =\cos\!\left(\frac{2\pi}{L}\log\frac{\scal}{U}\right).
\]
Its integer phase returns are at $\scal_j=Ue^{-jL}$, and its signed
winding along the \emph{continuous interpolating} depth $0\le q_L\le7$
is
\begin{equation}
\label{p3c:eq:v222-seven-winding}
 \nu_L=\frac1{2\pi i}\int_{0}^{7}
    \chi_{L,1}^{-1}\frac{d\chi_{L,1}}{dq_L}\,dq_L=-7.
\end{equation}
Consequently the primitive contact residue supplies the phase period,
and the protected Clifford rank defect $\dim_{\mathbb R}W=7$
supplies its finite mathematical depth.  Neither quantity is
fitted to the gravitational action.
\end{theorem}

\begin{proof}
A continuous character $\mathbb R/L\mathbb Z\to U(1)$ pulls back to
a continuous homomorphism $t\mapsto e^{i\omega t}$ on $\mathbb R$.
Periodicity requires $\omega L=2\pi m$ for some $m\in\mathbb Z$.
Normalizing the phase at $U$ gives
Eq.~\eqref{p3c:eq:v222-unitary-character-family}; faithfulness forces
$|m|=1$.  At $\scal_j=Ue^{-jL}$ the exponent is $-2\pi i mj$.
For $m=1$, $d\chi/\chi=-2\pi i\,dq_L$, giving
Eq.~\eqref{p3c:eq:v222-seven-winding}.
The finite rank-seven endpoint is imported from
Theorem~\ref{p3:thm:clifford-dim} and the protected-rank result,
not inferred from a numerical cosmological endpoint.
\end{proof}

\begin{remark}[Physical-rank visibility and non-dynamical status]
\label{p3c:rem:v222-phase-scope}
The continuous path $q_L\in[0,7]$ is a convenient \emph{interpolation}
of the determinant rank valuation; without $(\mathrm{FLAG})$ the
established active support depths remain $q=0,4,5,7$, not seven
independently proven shell transitions.  Integer phase returns are
available scale classes, not dynamically selected vacua.  The
construction fixes the primitive character of a \emph{chosen}
logarithmic-scale quotient, but does not derive a physical transport
operator that realizes it as the Hopf connection holonomy, nor make
$F_L$ an $f(R)$ coefficient.  The real cosine alone also loses the
orientation and integer winding information carried by $\chi$.
\end{remark}

\begin{remark}[Quarter-character transport and its exact owner]
\label{p3c:rem:v242-quarter-character-owner}
On a common logarithmic period \(L\), the positive carrier grading
\(H_{\rm car}=E_0e^{-L\mathcal Q/4}\) admits the constructed unitary
companion
\[
 U_H=\exp\!\left[\frac{2\pi i}{L}\log(H_{\rm car}/E_0)\right]
       =e^{-i\pi\mathcal Q/2}.
\]
On the \(Q_4\oplus Q_5\) corridor it is \(Q_4-iQ_5\); hence
\(U_HT(v)U_H^\dagger=-iT(v)\) and the reverse edge has phase \(+i\).
The angle \((2\pi/4)(5-4)=\pi/2\) is exact in this construction.
Under \(J_{\rm HG}\), the carrier operator is
\(Q_{\rm fam,0}-iQ_{\rm fam,\perp}\), not the diagonal rooted
joint-quadratic phase \(R_G\). Its conjugation is the identity on
\(Q_5\operatorname{End}Q_5\), including the operator trimer.
The existing joint Hopf phase has instead the unique signed real-vector
lift of Corollary~\ref{p1:cor:v242-joint-phase-vector-lift}.
These distinguish exact geometric constructions from a proof that a
physical source consumes either unitary port. The positive fourth-root
theorem does not by itself supply that unitary coupling, and the cosine
alone does not fix its orientation or its physical leg assignment.
\end{remark}


\subsection{IR family pullback of the native response and operator grading}
\label{p3c:sec:v2208-ir-family-response}
On the occurring active branch of \ref{p3:thm:derived-active-support},
retain one already rooted member. Let $J_{\rm IR}$ be the geometric
isometry of \ref{p1b:thm:v2208-ir-family-bridge}, and write
$Q_c=|c\rangle\langle c|$, $Q_\perp=I_G-Q_c$.
All results below concern this owner-level pullback. They neither project
the independent matter family factor nor specify a physical mass compiler.

\begin{proposition}[Exact response in the IR corridor]
\label{p3c:prop:v2208-preterminal-response}
For the native semigroup of Paper III-C,
\[
 R_\theta=e^{-7\theta}Q_0+e^{-3\theta}Q_4
              +e^{-2\theta}Q_5+Q_7,
\]
the geometric pullback is
\[
 \boxed{J_{\rm IR}^\dagger R_\theta J_{\rm IR}
        =Q_c+e^{-2\theta}Q_\perp.}
\]
For $r=e^{-2\theta}>0$, its inverse is $Q_c+r^{-1}Q_\perp$.
The literal endpoint is $Q_c$, whose Moore--Penrose inverse is also $Q_c$.
Scalar functional calculus of this one response stays in the algebra
$\mathbb C Q_c\oplus\mathbb C Q_\perp$.
\end{proposition}
\begin{proof}
The bridge image has no $Q_0$ or $Q_4$ component, and the earlier projector
identities evaluate the two remaining summands. Orthogonality gives
$(aQ_c+bQ_\perp)(dQ_c+eQ_\perp)=adQ_c+beQ_\perp$, proving the inverse
and functional-calculus claims. The endpoint is a self-adjoint projector
and hence its own Moore--Penrose inverse.
\end{proof}
Inverse before the endpoint limit diverges on the plane; taking the
Moore--Penrose inverse after the limit removes that plane. Operation order
therefore matters. This comparison does not restore the obsolete response
word or identify $\theta$ with cosmic time or a K16 preparation.
A family operator obtained solely by scalar functional calculus of this
single response has one singlet value and a repeated plane value; that
recipe alone cannot yield three distinct eigenvalues. This is not a
no-go for a separately generated role packet or a larger owner word.

By \ref{p3c:thm:v179-balanced-hierarchy-response}, the native balanced carrier grading is
\[
 H_{\rm car}=E_0Q_0+E_{32}Q_4+E_{40}Q_5+E_{56}Q_7,
 \qquad\mathcal E_H(A)=H_{\rm car}^{1/2}AH_{\rm car}^{1/2}.
\]
On $E_3(u)$ it gives
\[
 \begin{array}{c|c|c}
 \text{operator corner}&\text{dimension}&\text{balanced weight}\\ \hline
 Q_7\operatorname{End}Q_7&1&E_{56}\\
 Q_7\operatorname{End}Q_5\oplus Q_5\operatorname{End}Q_7
       &4&\sqrt{E_{56}E_{40}}=E_{48}\\
 Q_5\operatorname{End}Q_5&4&E_{40}
 \end{array}
\]
Conjugating by $J_{\rm IR}$ carries these corners to the family
singlet, mixed and plane operator corners. Their dimensions follow from
$1^2,2\cdot1\cdot2,2^2$; their weights follow by multiplying the two
square-root grades. The same calculation on the charged $Q_4\oplus Q_5$
corridor gives the already recorded labels $32,36,40$.
The mixed label $48$ here is an operator midpoint, not an additional active
K48 carrier or a predicted neutrino mass. Carrier grading and particle
mass coefficients remain differently typed quantities.


The coefficient-factor consequence is separately typed in
\ref{p5:thm:v2208-terminal-tensor-extension}. In particular the geometric
endpoint element $J_{\rm IR}^\dagger P_0J_{\rm IR}=Q_c$ is not the
terminal map on the independent family coefficient factor.


\section{Transverse normalizer bridge and conditional boundary promotion}
\label{p3c:sec:v166-transverse-normalizer-bridge}

\begin{definition}[Result-independent transverse normalizer-bridge target]
\label{p3c:def:v166-transverse-normalizer-bridge-target}
Fix an orbit member \(u\in\mathcal O_{\rm diag}\), put
\[
 N_u:=N_{2O}(C_6(u)),
\]
and use the conjugation representation on
\(A^{\rm even}=\mathbb R\mathbf1\oplus\operatorname{Im}\mathbb H\).
A transverse normalizer-bridge candidate is an \(N_u\)-stable real subspace
\(L\le A^{\rm even}\) such that
\begin{enumerate}[label=\textup{(\roman*)}]
 \item \(\mathbb R\mathbf1\subset L\);
 \item \(C_6(u)\) acts trivially on \(L\); and
 \item the induced action of \(N_u/C_6(u)\) on
 \(L/\mathbb R\mathbf1\) is faithful.
\end{enumerate}
The target role is the complete unmarked orbit family of the least such
subspaces, if a unique least family exists.  This definition mentions no
depth, shell coordinate, energy, particle, boundary operator, RAR law, or
desired rank.
\end{definition}

\begin{theorem}[Unique binary-dihedral transverse bridge and derived
depth-six placement]
\label{p3c:thm:v166-unique-transverse-k48-bridge}
On the occurring active branch, define the even Reynolds projector
\begin{equation}
 \label{p3c:eq:v166-transverse-reynolds-projector}
 B_\perp(u):=\Pi_4\mathcal A_6(u)\Pi_4.
\end{equation}
Then
\begin{equation}
 \label{p3c:eq:v166-transverse-bridge-range}
 \boxed{
 B_\perp(u)=P_0+\Pi_u,
 \qquad
 \operatorname{Ran}B_\perp(u)
 =\mathbb R\mathbf1\oplus\mathbb Ru,
 \qquad
 \operatorname{rank}_{\mathbb R}B_\perp(u)=2.}
\end{equation}
Moreover Paper~I-B's diagonal-normalizer theorem gives
\begin{equation}
 \label{p3c:eq:v166-binary-dihedral-quotient}
 N_u\cong2D_3,
 \qquad C_6(u)\triangleleft N_u,
 \qquad N_u/C_6(u)\cong C_2,
\end{equation}
and the restricted real representation is
\begin{equation}
 \label{p3c:eq:v166-even-normalizer-decomposition}
 A^{\rm even}\!\downarrow_{N_u}
 \cong
 \mathbf1\oplus\operatorname{sgn}\oplus\rho_2,
 \qquad
 \operatorname{Ran}B_\perp(u)
 \cong\mathbf1\oplus\operatorname{sgn}
 \cong\mathbb R[C_2],
\end{equation}
where \(\rho_2\) is the real two-dimensional reflection representation on
\(u^\perp\).  Hence \(\operatorname{Ran}B_\perp(u)\) is the unique candidate
of Definition~\ref{p3c:def:v166-transverse-normalizer-bridge-target}, so the
complete orbit family is the GP1-finite, GP2-exact and host-faithful,
GP3-unique-least realization of that target.

Finally the common trace defect gives
\begin{equation}
 \label{p3c:eq:v166-transverse-bridge-depth}
 \boxed{
 \delta_W(B_\perp(u))=8-2=6,
 \qquad K_{\rm sh}=8\delta_W=48.}
\end{equation}
Only after this calculation may the structurally placed projector be denoted
\(B_{48}(u):=B_\perp(u)\).
\end{theorem}

\begin{proof}
The proof of Theorem~\ref{p3:thm:derived-active-support} already shows that
conjugation by \(C_6(u)\) fixes precisely
\(\mathbb R\mathbf1\oplus\mathbb Ru\) inside \(A^{\rm even}\).  Therefore its
normalized Reynolds operator is the orthogonal projector
\(P_0+\Pi_u\), proving \eqref{p3c:eq:v166-transverse-bridge-range}.

Paper~I-B, Theorem~\ref{p1:thm:binary-dihedral-diagonal-lift}, proves
\eqref{p3c:eq:v166-binary-dihedral-quotient}.  The cyclic subgroup fixes the
oriented axis \(u\) and rotates \(u^\perp\) through \(2\pi/3\).  Every element
in the nontrivial quotient coset reverses the oriented diagonal, hence sends
\(u\mapsto-u\), while fixing the scalar.  Its action on \(u^\perp\), together
with the order-three rotation, is the standard real reflection
representation of \(D_3\cong S_3\).  This proves
\eqref{p3c:eq:v166-even-normalizer-decomposition}.

If \(L\) satisfies Definition~\ref{p3c:def:v166-transverse-normalizer-bridge-target},
condition~(ii) places \(L\) inside the \(C_6(u)\)-fixed space
\(\mathbb R\mathbf1\oplus\mathbb Ru\).  Conditions~(i) and~(iii) force the
scalar and quotient-odd lines both to occur.  Thus
\(L=\mathbb R\mathbf1\oplus\mathbb Ru\), proving existence, uniqueness, and
leastness without using the output rank or clock.  Reynolds averaging,
normalizer restriction, finite image, and complete-orbit retention are
result-independent natural finite operations on the already occurring
branch.  The generated-role theorem of Paper~II-A therefore retains their
nonredundant image in the declared structural-subspace target.  The defect
formula is Theorem~\ref{p3c:thm:common-trace-defect-spectral-filtration}; since
\(B_\perp\) contains the scalar line and has rank two, it gives six and hence
\(K_{\rm sh}=48\).
\end{proof}

\begin{theorem}[Transverse diamond, active-gap noncollapse, and exact
half-step readouts]
\label{p3c:thm:v166-transverse-diamond-halfsteps}
Put \(P_{40}(u):=\Pi_3(u)\).  In the orthogonal projector lattice of
\(A^{\rm even}\),
\begin{equation}
 \label{p3c:eq:v166-transverse-projector-diamond}
 \boxed{
 P_{40}(u)B_{48}(u)=P_0,
 \qquad
 P_{40}(u)\wedge B_{48}(u)=P_0,
 \qquad
 P_{40}(u)\vee B_{48}(u)=\Pi_4.}
\end{equation}
In particular \(B_{48}(u)\not\le P_{40}(u)\), whereas
\begin{equation}
 \label{p3c:eq:v166-active-gap-preserved}
 \boxed{
 \mathbf1_{\{6\}}(\mathcal Q_{\rm act}(u))=0.}
\end{equation}
Thus the rank-two transverse bridge exists without filling the nested active
carrier gap.

The canonical equal-weight algebraic barycentres of the displayed endpoint
pairs are
\begin{align}
 \label{p3c:eq:v166-k44-barycentre}
 M_{44}(u)
 &:=\frac12\bigl(P_{40}(u)+B_{48}(u)\bigr)
   =\frac12(\Pi_4+P_0),
 &\delta_W(M_{44})&=\frac{11}{2},\\
 \label{p3c:eq:v166-k52-barycentre}
 M_{52}(u)
 &:=\frac12\bigl(B_{48}(u)+P_0\bigr)
   =P_0+\frac12\Pi_u,
 &\delta_W(M_{52})&=\frac{13}{2}.
\end{align}
Consequently their exact trace placements are \(K=44\) and \(K=52\).
The first is independent of \(u\); the second is retained only as its complete
equivariant orbit family.  Both are positive contractions and neither is an
orthogonal projector.

Independently, the already generated common-rate semigroup has unique points
at the same three trace coordinates:
\begin{equation}
 \label{p3c:eq:v166-generated-halfstep-semigroup-points}
 \boxed{\begin{gathered}
 V_{44}:=V_{\log(14/3)}=P_0+\frac{3}{14}P_W,\\
 V_{48}:=V_{\log7}=P_0+\frac17P_W,\\
 V_{52}:=V_{\log14}=P_0+\frac1{14}P_W.
 \end{gathered}}
\end{equation}
Equations~\eqref{p3c:eq:v166-k44-barycentre}--
\eqref{p3c:eq:v166-generated-halfstep-semigroup-points} are exact
contraction/trace identities.  They do not identify the barycentres with the
semigroup operators, insert active carriers, or assert physical events or a
transition history.
\end{theorem}

\begin{proof}
The orthogonal decomposition
\[
 A^{\rm even}
 =\mathbb R\mathbf1\oplus\mathbb Ru\oplus u^\perp
\]
gives
\(P_{40}=P_0+P_{u^\perp}\) and
\(B_{48}=P_0+\Pi_u\).  Their product and range intersection are \(P_0\),
while their range sum is all of \(A^{\rm even}\), proving the lattice
identities.  Since the bridge contains the omitted axis rather than lying in
\(E_3(u)=\mathbb R\mathbf1\oplus u^\perp\), it does not satisfy the nesting
hypothesis of the active chain.  Equation~\eqref{p3c:eq:v166-active-gap-preserved}
is therefore exactly the independent active-depth-six obstruction, not a
contradiction.

Adding the two orthogonal-block formulas gives the displayed barycentres.
On \(W\), \(M_{44}\) has three eigenvalues \(1/2\) and four zeros, whereas
\(M_{52}\) has one eigenvalue \(1/2\) and six zeros.  Hence
\[
 \delta_W(M_{44})=7-\frac32=\frac{11}{2},
 \qquad
 \delta_W(M_{52})=7-\frac12=\frac{13}{2}.
\]
Their eigenvalue \(1/2\) proves non-idempotence.  Finally
\(q_{\rm sp}=7(1-e^{-\tau})=K/8\) gives respectively
\(e^{-\tau}=3/14,1/7,1/14\) at \(K=44,48,52\), proving
\eqref{p3c:eq:v166-generated-halfstep-semigroup-points}.  No equality between
operators with merely equal trace defect follows or is used.
\end{proof}

\begin{corollary}[Exact scope of the depth-six bridge]
\label{p3c:cor:v166-transverse-bridge-scope}
The generated conclusion at depth six is the structural orbit family
\(\{B_{48}(u)\}_{u\in\mathcal O_{\rm diag}}\) in the transverse-normalizer
bridge sort.  It is not a native one-particle carrier, pole, state,
occupation, stress tensor, action atom, boundary jet, RAR kernel, or
cosmological transition.  Promotion to any such target requires a separately
proved owner-preserving map into that target sort.  In particular the active
carrier-gap theorem and every zero-production consequence of
\(P_{48}^{\rm act}=0\) remain unchanged.
\end{corollary}


\subsection{owner-typed boundary promotion of the depth-six bridge}
\label{p3c:sec:v158-b48-boundary-promotion}

\begin{definition}[Owner-preserving boundary promotion interface]
\label{p3c:def:v158-b48-boundary-promotion}
Let
\[
 \mathbb B_{48}
 :=\{\operatorname{Ran}B_{48}(u):u\in\mathcal O_{\rm diag}\}
\]
be the complete structural transverse-normalizer bridge family of
Theorem~\ref{p3c:thm:v166-unique-transverse-k48-bridge}.  A physical boundary
promotion of this family consists of an actually constituted boundary split
of a same-action stationary problem,
\[
 T_{\bar q}\mathcal C_{\rm stat}=J_\partial\oplus J_I,
\]
together with a generated, owner-preserving family of linear maps
\begin{equation}
 \label{p3c:eq:v158-b48-boundary-map}
 \boxed{
 \beta_{48,u}:\operatorname{Ran}B_{48}(u)\longrightarrow J_\partial}
\end{equation}
that is equivariant under the retained unmarked orbit action and whose
constructor ancestry connects the structural bridge root to the boundary
source/trace owner.  The definition does not assume that such a family
exists, is injective, or is unique.
\end{definition}

\begin{theorem}[Current III-C data do not generate a boundary owner for
\(B_{48}\)]
\label{p3c:thm:v158-b48-boundary-owner-nogo}
The generated data of this paper determine the domain family
\(\mathbb B_{48}\), its binary-dihedral quotient action, its rank two, and its
trace-defect placement \(K_{\rm sh}=48\).  They do not determine a physical
boundary tangent \(J_\partial\), a boundary source/trace constructor, or a
word with output sort
\[
 \operatorname{Hom}(\operatorname{Ran}B_{48}(u),J_\partial).
\]
Consequently G0, GP2 co-realization, equality of dimensions or trace defects,
and the representation identity
\(\operatorname{Ran}B_{48}\cong\mathbb R[C_2]\) do not by themselves produce
\eqref{p3c:eq:v158-b48-boundary-map}.  In particular
\begin{equation}
 \label{p3c:eq:v158-b48-no-direct-promotion}
 \boxed{
 B_{48}\not\Rightarrow
 J_\partial,\quad
 T_\partial,\quad
 F_{\rm contact},\quad
 a_{\rm RAR}.}
\end{equation}
Any positive promotion requires the separately generated owner-preserving
interface of Definition~\ref{p3c:def:v158-b48-boundary-promotion}.
\end{theorem}

\begin{proof}
Theorem~\ref{p3c:thm:v166-unique-transverse-k48-bridge} is a theorem in the
structural-subspace target sort.  Its constructor word uses finite Reynolds
averaging, normalizer restriction, and orbit retention; none has a boundary
source/trace output.  Paper~II-A's constructor-origin and same-owner rules can
retain and constrain the output of a licensed constructor, but they cannot
create a missing cross-sort constructor merely from numerical or
representation-theoretic equality.  Since the present signature contains no
constituted \(J_\partial\) owner and no map word into it, the displayed
boundary Hom-sort has no generated inhabitant.  This is exactly the owner
firewall already stated qualitatively in
Corollary~\ref{p3c:cor:v166-transverse-bridge-scope}, now expressed in the
current same-owner theorem language.
\end{proof}

\begin{theorem}[Same-action factorization once the boundary map exists]
\label{p3c:thm:v158-b48-schur-factorization}
Assume Definition~\ref{p3c:def:v158-b48-boundary-promotion} is independently
satisfied on a branch for which Paper~III-F's source-conditioned boundary
Schur theorem applies.  Let
\[
 K_{\partial,y}
 =H_{\partial\partial}(y)
  -H_{\partial I}(y)H_{II}(y)^{-1}H_{I\partial}(y)
\]
be that same-action boundary Hessian.  Then the quadratic response seen by
the promoted depth-six bridge is not a second fitted kernel; it is uniquely
the pullback quadratic form
\begin{equation}
 \label{p3c:eq:v158-b48-schur-pullback}
 \boxed{
 K_{48,y}(u)
 :=\beta_{48,u}^{*}K_{\partial,y}\beta_{48,u}.}
\end{equation}
If the promoted map and the boundary split are generated descendants of the
same occurring root, GP2 co-realizes the structural bridge, its boundary
image, and \eqref{p3c:eq:v158-b48-schur-pullback}; no second occurrence is
required.  If \(\beta_{48,u}\) itself depends on the source, derivatives of
the response obey the full chain rule and are not equal merely to
\(\beta_{48,u}^{*}(\nabla_yK_\partial)\beta_{48,u}\).
\end{theorem}

\begin{proof}
Once \(\beta_{48,u}\) and the physical boundary owner are given, the
quadratic boundary action evaluated on
\(u_\partial=\beta_{48,u}b\) is
\[
 \frac12\langle \beta_{48,u}b,
 K_{\partial,y}\beta_{48,u}b\rangle
 =\frac12\langle b,K_{48,y}(u)b\rangle,
\]
which proves \eqref{p3c:eq:v158-b48-schur-pullback}.  Paper~III-F proves that
\(K_{\partial,y}\) is already the unique stationary-interior descendant of
the same action.  Therefore an independently fitted \(K_{48}\) would insert a
second response owner.  The source-dependent statement is the ordinary
chain rule and prevents a hidden assumption that the promotion map is rigid
under source variation.
\end{proof}

\begin{theorem}[Terminal order of operations for a promoted \(B_{48}\) response]
\label{p3c:thm:v158-b48-terminal-order}
Suppose in addition that an active-side terminal family satisfies the
rank-changing hypotheses of Paper~III-F's eliminate-before-quotient theorem
and that \(\beta_{48,u,\varepsilon}\) has a declared finite terminal limit
\(\beta_{48,u,T}\).  Then the terminal depth-six boundary response is
\begin{equation}
 \label{p3c:eq:v158-b48-terminal-response}
 \boxed{
 K_{48,T}(u)
 =\beta_{48,u,T}^{*}
 \left[\lim_{\varepsilon\downarrow0}
 \left(A_\varepsilon
       -B_\varepsilon^{*}C_\varepsilon^{-1}B_\varepsilon\right)\right]
 \beta_{48,u,T}.}
\end{equation}
Replacing the bracket by the quotient-first block
\(A_T=\lim_{\varepsilon\downarrow0}A_\varepsilon\) is justified only when the
mixed Schur memory obeys
\[
 \lim_{\varepsilon\downarrow0}
 B_\varepsilon^{*}C_\varepsilon^{-1}B_\varepsilon=0
\]
or an independently proved identity gives the same result.  Thus a terminal
scalaron coordinate may disappear while a finite boundary response on the
promoted \(B_{48}\) image survives.
\end{theorem}

\begin{corollary}[Ward and critical-law firewalls on the promoted bridge]
\label{p3c:cor:v158-b48-ward-critical-firewall}
Under the same hypotheses, Paper~III-F's Ward--Schur theorem gives
\[
 K_{\partial,y}R_{A,\partial}=0
\]
for every actual boundary component of a generated symmetry direction.  Thus
if \(b\in\operatorname{Ran}B_{48}(u)\) happens to satisfy
\(\beta_{48,u}b=R_{A,\partial}\), then
\[
 K_{\partial,y}\beta_{48,u}b=0.
\]
Noether therefore constrains the promoted response after the owner map exists,
but it does not imply that all of \(\beta_{48,u}(\operatorname{Ran}B_{48})\)
is a Ward-null space.  Nor do
Theorems~\ref{p3c:thm:v158-b48-schur-factorization} and
\ref{p3c:thm:v158-b48-terminal-order} prove a particular all-order critical
law such as \(\mu^2(y)=y\), a square-root RAR lift, or a boundary-stress
normalization.  Those remain calculations of the actual same-action source,
solution, promotion map, and higher variations.
\end{corollary}

\begin{remark}[Scope of the normalized boundary response]
\label{p3c:rem:v158-b48-boundary-scope}
Version~1.58 closes an order/provenance question, not the physical RAR owner.
The positive result is the forced factorization
\[
 \boxed{
 B_{48}
 \xrightarrow{\ \beta_{48}\ }
 J_\partial
 \xrightarrow{\ \text{same-action Schur}\ }
 K_\partial
 \xrightarrow{\ \text{pullback}\ }
 K_{48}}
\]
conditional on the first arrow.  The first arrow itself, a terminal crossing
history, a physical boundary stress, and the all-order critical response law
remain open.
\end{remark}


\subsection{Conditional normalization after a licensed physical $B_{48}$ boundary port}
\label{p3c:sec:v184-g25-conditional-normalization}
The abstract $C_2$-equivariant map family is
\[
 \beta_{a,b}(x1+yu)=ax\,d+by\,n.
\]
Unit preservation fixes $a=1$, leaving $\beta_b:=\beta_{1,b}$.  Suppose,
\emph{in addition}, that a licensed owner-preserving physical port
$\beta_b:B_{48}\to J_\partial$ has already been constructed and that the same
owner supplies nonzero scalar functionals
\[
 \lambda_{48}:\mathbb Ru_-\to\mathbb R,
 \qquad
 \lambda_\partial:\mathbb Rn_-\to\mathbb R
\]
with the physical source/flux compatibility
\[
 \lambda_\partial\circ\beta_b=\lambda_{48}.
\]
Then
\[
 \boxed{b=\frac{\lambda_{48}(u_-)}{\lambda_\partial(n_-)}}.
\]
If both functionals are normalized to one with the same generated orientation,
this conditional formula gives $b=1$.

This calculation classifies the residual normalization \emph{after} the
physical cross-sort port and its common owner exist.  It does not weaken
Theorem~\ref{p3c:thm:v158-b48-boundary-owner-nogo}: the current Paper-III-C
data still do not generate a word with output sort
$\operatorname{Hom}(B_{48},J_\partial)$ or the required source/trace owner.
The G25 physical target requires more than the local geometric/PCAD
port. The source construction supplies the saturated Clifford-normal parent and its
\emph{local} PCAD germ through
Theorem~\ref{p3c:thm:v187-b48-boundary-parent}; this does not independently
construct a stationary physical same-action boundary action, its source
trace \(J_\partial\), or the owner-preserving physical
\(\beta_{48}:B_{48}\to J_\partial\). Thus the earlier no-go remains
a firewall against promoting the local Hom data to that physical sort.
\emph{If} those physical antecedents are supplied, the displayed same-owner
normalization is a valid consistency check and
Theorem~\ref{p3c:thm:v158-b48-schur-factorization} produces the
conditional \(K_{48}\) Hessian family. The local PCAD result alone does
not discharge this physical-response premise.



\subsection{The affine normal phase retains its local connection contact}
\label{p3c:sec:v229-normal-phase}
The conormal comparison uses the explicit frozen phase
$\theta(R)=\arctan(R/R_s)$ of
Proposition~\ref{p1b:prop:v229-affine-conormal}. If curvature varies,
\[
 U_R^\dagger\nabla U_R=\tfrac i2\Gamma_Wd\theta,\qquad
 d\theta=\frac{R_s}{R_s^2+R^2}\,dR.
\]
These follow by differentiating the same unitary frame. Any use in a
local gravitational action retains this connection and the conformal
metric contacts already supplied by Paper~III-F. A frozen spectral
identity alone does not identify this normal phase with a separately
calibrated hierarchy clock or visibility phase. The complete cap
application retains its actual curvature-to-source map and the
absolute boundary weight of Eq.~\eqref{p3e:eq:v229-amplitude-target}.




\section{Carrier, response, and terminal visibility interfaces}
\label{p3c:sec:v184-carrier-response-terminal-midpoint}

\subsection{Hard carriers and balanced response grades}
On the protected carrier skeleton the realized carrier depths are
\[
  \boxed{q_{\rm car}\in\{0,4,5,7\}},
\]
while balanced pair-response grades occur at
\[
  \boxed{K\in\{16,20,28,36,44,48\}}.
\]
For a carrier pair $i<j$ with depths $q_i,q_j$, the balanced response has
\[
 q_{ij}=\frac{q_i+q_j}{2},\qquad K_{ij}=4(q_i+q_j),\qquad
 E_{ij}=\sqrt{E_iE_j}.
\]
The physically consumed adjacent-carrier midpoints include
\[
 E_{16}=\sqrt{E_0E_{32}},\qquad
 E_{36}=\sqrt{E_{32}E_{40}},\qquad
 E_{48}=\sqrt{E_{40}E_{56}}.
\]
These are response grades; they are not additional particle carriers.

\subsection{Terminal $q=5\to7$ corridor and $K=48$ realization bifurcation}
With the terminal axis fixed, write
\[
 A^{\rm even}=S_7\oplus N_4\oplus T_5,
 \qquad P_{40}=P_S+P_T,\qquad P_{56}=P_S.
\]
The real $C_3$ action on $T_5$ is irreducible, so there is no invariant real
line between $S_7$ and $S_7\oplus T_5$.  Hence
\[
 \boxed{P_{48}^{\rm act}=0.}
\]
The smooth relative contraction on the active terminal carrier is
\[
 U_\sigma=P_S+e^{-\sigma}P_T.
\]
At the balanced midpoint $q=6$, equivalently $K=48$,
\[
 \boxed{U_{48}=P_S+\frac12P_T,\qquad U_{48}^2\ne U_{48}.}
\]
Thus $K=48$ is the midpoint of the full $q=5\to7$ depletion and the onset of
the final primitive $q=6\to7$ subtransfer; it is not the beginning of the
entire terminal depletion.

The sharp transverse static object remains
\[
 \boxed{B_{48}=P_S+P_N,\qquad B_{48}^2=B_{48}.}
\]
For the primitive residual $D_{48}=P_W/7$ one has the normalized compressions
\[
 \frac{P_TD_{48}P_T}{\operatorname{Tr}(P_TD_{48}P_T)}=\frac12P_T,
 \qquad
 \frac{P_ND_{48}P_N}{\operatorname{Tr}(P_ND_{48}P_N)}=P_N.
\]
These are covariance/compression identities, not quantum-state conditioning.

\subsection{Static/dynamic typing firewall}
The balanced depth-six operator-response module is
\[
 \mathcal R_{48}^{\rm op}
 =\operatorname{Hom}(T_5,S_7)\oplus\operatorname{Hom}(S_7,T_5).
\]
It is distinct from $B_{48}=S_7\oplus N_4$.  On an unbroken $C_3$ branch,
\[
 \boxed{\operatorname{Hom}_{C_3}(T_5,N_4)=0,}
\]
so no unmarked first-order $E_{48}\to B_{48}$ map or generic linear
$T_5\leftrightarrow B_{48}$ Hessian mixing is available without an additional
transforming background/order parameter.


\subsection{Carrier, response, and action support are distinct}
\label{p3c:sec:v187-response-support-taxonomy}

\begin{definition}[Three support levels]
\label{p3c:def:v187-three-support-levels}
For a hierarchy grade distinguish
\emph{carrier support} (an actual state/eigenspace),
\emph{response support} (a nonzero allowed Hom/operator block), and
\emph{action support} (an actually consumed physical action/source block).
Neither implication
\[
\begin{gathered}
 \mathsf{CarrierSupport}=0\Rightarrow\mathsf{ResponseSupport}=0,\\
 \mathsf{ResponseSupport}\ne0\Rightarrow\mathsf{ActionSupport}\ne0
\end{gathered}
\]
is valid without an additional theorem.
\end{definition}

\begin{proposition}[Terminal $E_{48}$ response without an unmarked invariant source]
\label{p3c:prop:v187-e48-terminal-symmetry-firewall}
The $q=5\leftrightarrow7$ response block has the balanced weight
$E_{48}=\sqrt{E_{40}E_{56}}$ even though there is no $K=48$ carrier.  On the
unmarked terminal $C_3$ branch, $T_5$ is the nontrivial real two-dimensional
rotation module while $S_7$ is trivial; hence
\[
 \boxed{\operatorname{Hom}_{C_3}(T_5,S_7)=0.}
\]
Thus a nonzero physical $E_{48}$ action/source coupling requires an additional
transforming source/background or another licensed same-action incidence.
\end{proposition}

\begin{proposition}[Invariant $E_{44}$ endpoint edge]
\label{p3c:prop:v187-e44-endpoint-edge}
On the rooted quaternionic endpoint lines
$Q_4=|u\rangle\langle u|$ and $Q_7=|1\rangle\langle1|$, left multiplication
by $u$ gives
\[
 \boxed{Q_7L_uQ_4=-|1\rangle\langle u|}.
\]
Its self-adjoint completion is a nonzero $q4\leftrightarrow q7$ response edge
with balanced weight
\[
 \boxed{E_{44}=\sqrt{E_{32}E_{56}}},\qquad
 \frac{E_{44}^2}{E_{32}}=E_{56}.
\]
This is a geometric operator-envelope result.  Its use by the physical finite
fermion action remains an action-incidence question.
\end{proposition}

\begin{remark}[No adjacency selector]
\label{p3c:rem:v187-no-adjacency-selector}
The complete response table contains $E_{36}$, $E_{44}$ and $E_{48}$.
Adjacency in hierarchy depth does not select physical consumption.  Action
ownership, representation typing and source incidence decide which response
blocks enter a physical claim.
\end{remark}


\subsection{Rank-seven visibility/local-jet type firewall}
\label{p3c:sec:v188-n7-visibility-firewall}
\begin{corollary}[Unique normalized unresolved rank-seven trace]
\label{p3c:cor:v188-normalized-unresolved-rank-seven-trace}
On the protected decomposition $\mathbb R\mathbf1\oplus W_7$ with
$V_\tau|_{W_7}=e^{-\tau}I_{W_7}$, the normalized unresolved trace is
\[
 \boxed{\nu_7(\tau):=\frac1{7}\operatorname{Tr}_{W_7}V_\tau=e^{-\tau}}
 \qquad\text{and}\qquad
 q_{\rm sp}=7(1-\nu_7).
\]
Through the generated hierarchy character $R(q)=U\alpha_\sigma^{-q}$ and
$C_7=U\alpha_\sigma^{-7}$ this is equivalently
\[
 \boxed{\nu_7(R)=\frac{\log(R/C_7)}{\log(U/C_7)}}
 \quad(C_7\le R\le U).
\]
In particular $\nu_7(C_6)=1/7$ and $\nu_7(C_7)=0$.
\end{corollary}
\begin{remark}[Visibility is not the local curvature Hessian]
\label{p3c:rem:v188-visibility-not-curvature-hessian}
The scalar $\nu_7$ is a normalized unresolved visibility trace.  It is not
by itself a theorem identifying the local metric susceptibility $f_{RR}$.
That local second jet is a configuration-tangent object and is derived in
Papers~III-D--III-F by holding the generated scheme/depth coefficient fixed
during the local variation.  Thus the provisional identification
$f_{RR}\propto\nu_7$ is not part of the active branch.
\end{remark}




\subsection{successor-cell homology and reciprocal orientation}
\label{p3c:sec:v190-successor-cell-homology}

\begin{theorem}[Same-role cell translation versus octet reflection]
\label{p3c:thm:v190-cell-translation-reflection}
On the protected hierarchy let the early cell be
$K=8+8s$ and the late cell be $K=48+8s$, $0\le s\le1$.
The translation
\[
 \boxed{T_{40}:K\mapsto K+40}
\]
preserves the local completion parameter and maps
\[
 8\mapsto48,\qquad16\mapsto56.
\]
By contrast the global octet reflection
\[
 \boxed{J_8:K\mapsto64-K}
\]
acts by
\[
 \boxed{s\mapsto1-s}
\]
and maps $8\leftrightarrow56$, $16\leftrightarrow48$.
Thus $T_{40}$ compares the same local response role, whereas $J_8$ exchanges
reciprocal/complementary roles.
\end{theorem}

\begin{corollary}[Exact normalized-energy homology]
\label{p3c:cor:v190-normalized-energy-homology}
For
\[
 E_K=\sqrt U\,\alpha^{-K/32},
\]
one has for every $0\le s\le1$
\[
 \boxed{
 \frac{E_{8+8s}}{E_8}
 =\frac{E_{48+8s}}{E_{48}}
 =\alpha^{-s/4}.}
\]
In particular
\[
 \frac{E_{16}}{E_8}=\frac{E_{56}}{E_{48}}=\alpha^{-1/4}.
\]
This is a locally rebased cell-role equivalence; it does not assert raw global
identities such as $E_8=E_{48}$.
\end{corollary}

\subsection{K48 terminal quadratic-response typing}
\label{p3c:sec:v192-k48-quadratic-typing}

\begin{proposition}[Complete invariant K48 terminal packet]
\label{p3c:prop:v192-k48-complete-invariant-return}
On the unbroken terminal $C_3$ branch there is no selected nonzero linear
$T_5\to S_7$ member, while the equivariant Clifford family
\[
 A_v=Q_5L_vQ_7,
 \qquad
 S_{57}(v)=A_v+A_v^\dagger
\]
remains generated.  Its inherited edge Gram form is
\[
 \tau_8\!\left(S_{57}(v)S_{57}(w)\right)
 =\frac14\langle v,w\rangle,
 \qquad
 \boxed{G_{48}^{\rm edge}=\frac14 I_{T_5}}.
\]
The normalized unmarked diffuse terminal tensor is
\[
 \boxed{\Sigma_{48}=\frac12P_T,
 \qquad \operatorname{Tr}_{T_5}\Sigma_{48}=1.}
\]
Thus the symmetry-safe terminal datum is the complete invariant quadratic
packet $(G_{48}^{\rm edge},\Sigma_{48})$, not a preferred transverse member.
\end{proposition}

\begin{corollary}[Quarter contraction is an edge normalization, not an H response law]
\label{p3c:cor:v192-k48-quarter-typing}
The unmarked scalar contraction of the generated edge packet is
\[
 \boxed{\chi_{48}^{\rm edge}
 =\operatorname{Tr}_{T_5}
 (\Sigma_{48}G_{48}^{\rm edge})=\frac14.}
\]
This identity is a structural quadratic-response normalization.  It is not,
without an additional same-action boundary theorem, the normalized homogeneous
Schur self-energy $\Sigma_H/K_H$, a finite $H$-stiffness decrement, a Cayley
return coordinate, or an abundance multiplier.
\end{corollary}

\begin{proposition}[All-orders uninserted spectral symmetry of the isolated K48 edge]
\label{p3c:prop:v1101-terminal-edge-all-orders}
On the isolated generated $S_7\oplus T_5$ self-adjoint Clifford block,
choose an orthonormal basis and write
\[
 D(v)=\begin{pmatrix}a&t v^T\\ t v&bI_2\end{pmatrix},\qquad
 (a,b,t)=(E_{56},E_{40},E_{48}),\quad t^2=ab,
 \qquad v\in T_5.
\]
For $P=\operatorname{diag}(-1,1,1)$ one has
$P D(v)P=D(-v)$; for every $O\in O(2)$ the orthogonal matrix
$\operatorname{diag}(1,O)$ conjugates $D(v)$ into $D(Ov)$. Moreover
\[
 \boxed{\det(\lambda I-D(v))=(\lambda-b)
 \bigl[(\lambda-a)(\lambda-b)-t^2|v|^2\bigr].}
\]
Its eigenvalues are
\[
 \boxed{b,\qquad \frac{a+b}{2}\pm
 \sqrt{\frac{(a-b)^2}{4}+t^2|v|^2}.}
\]
Every uninserted finite spectral trace $\operatorname{tr}g(D(v))$
therefore depends on $v$ only through $|v|^2$; wherever the trace is
smooth, all odd homogeneous Taylor degrees vanish. For the normalized
octet trace $\tau_8=\operatorname{Tr}_8/8$ and sufficiently smooth $g$,
\[
 \boxed{D_v^2\,\tau_8g(D(v))\big|_{v=0}[w_1,w_2]
 =\frac{t^2}{4}\frac{g'(a)-g'(b)}{a-b}
 \langle w_1,w_2\rangle,}
\]
with continuous $a=b$ limit $(t^2/4)g''(a)\langle w_1,w_2\rangle$.
Spectator blocks independent of $v$ do not alter this Hessian.
This accidental $O(2)$-evenness of the \emph{isolated, uninserted}
block is stronger than the terminal $C_3$ permission rules; it cannot
be promoted to a physical particle, a general marked-source selection
rule, or a K48 Hilbert-energy normalization.
\end{proposition}
\begin{proof}
The stated orthogonal conjugations give sign and rotational invariance.
The block determinant gives the characteristic polynomial and the
three eigenvalues. The second variation is the divided difference of
$g'$ at $a,b$, multiplied by the fixed normalized octet trace. No
new constituent or coefficient is assumed.
\end{proof}


\begin{remark}[Residual-coordinate firewall]
\label{p3c:rem:v192-k48-residual-firewall}
The absolute primitive residual $D_{48}=P_W/7$ and its normalized compression
$\Sigma_{48}=P_T/2$ have different target sorts from a normalized
successor-cell return coordinate.  In particular the number $1/7$ may not be
inserted into a K40--K48 or K48--K56 Cayley law unless a theorem identifies the
same owner, normalization and return process.  Likewise the exact quarter above
may not be promoted through $1-1/4=3/4$ to an $H$ response without the complete
same-action boundary reduction.
\end{remark}

\subsection{K48 homogeneous $H$--$T_5$ Schur typing}
\label{p3c:sec:v193-k48-h-t5-schur-typing}

\begin{corollary}[No unbroken $C_3$ mixed $H$--$T_5$ quadratic block]
\label{p3c:cor:v193-k48-h-t5-mixed-zero}
On the physical homogeneous K48 claim domain, the retained $H$ response
coordinate is a real scalar and hence carries the trivial $C_3$ representation,
whereas $T_5$ is the irreducible nontrivial real two-dimensional $C_3$ rotation
module of Proposition~\ref{p3c:prop:v192-k48-complete-invariant-return}.
Therefore every $C_3$-equivariant quadratic mixed block obeys
\[
 \boxed{
 B_{H,T_5}\in\operatorname{Hom}_{C_3}(T_5,H)
 \simeq\operatorname{Hom}_{C_3}(T_5,\mathbb R)=0.}
\]
Thus the unbroken terminal plane can contribute invariant quadratic returns and
higher source jets, but it cannot generate a linear quadratic Hessian bridge to
the homogeneous scalar $H$ coordinate.
\end{corollary}

\begin{remark}[Quadratic decoupling does not delete higher K48 response jets]
\label{p3c:rem:v193-k48-higher-jet-firewall}
The vanishing of $B_{H,T_5}$ is a statement about the second variation on the
unbroken branch.  It does not set the already generated invariant quadratic
packet $(G_{48}^{\rm edge},\Sigma_{48})$ to zero and does not forbid a
source--$H$--$H$ third variation obtained after the complete terminal relation
is consumed.  In particular the exact edge quarter remains a quadratic/source
normalization rather than a finite homogeneous Schur decrement.
\end{remark}
\begin{proposition}[One quadratic coefficient, with cubic C3 terms retained]
\label{p3c:prop:v257-k48-invariant-jets}
Let $h$ be a real trivial $C_3$ coordinate and $v\in T_5$ its nontrivial
real plane. For a real $C^4$ invariant same-action function, set
$z=v_1+iv_2$. Through degree four in $v$,
\[
 S(h,v)=S_0(h)+\tfrac12\kappa_{48}(h)|v|^2
       +\operatorname{Re}(c(h)z^3)+\tfrac14\lambda_{48}(h)|v|^4
       +o(|v|^4).
\]
Consequently the quadratic mixed block $D_hD_vS(h,0)$ vanishes, while
\[
 D_v^2S(h,0)=\kappa_{48}(h)I_T,\qquad
 \operatorname{Tr}_T(\Sigma_{48}D_v^2S)=\kappa_{48}(h),
 \qquad\Sigma_{48}=\tfrac12I_T.
\]
The contracted first and second $h$ jets equal $\kappa_{48}'$ and
$\kappa_{48}''$ on their common derivative domain. An interpretation
as a physical packet-source derivative requires its actual source map
and same-action incidence certificate.
\end{proposition}
\begin{proof}
The monomial $z^p\bar z^q$ is invariant exactly when $p-q\equiv0$
modulo three. At degrees one through four the invariant pairs are
respectively none, $(1,1)$, $(3,0),(0,3)$ and $(2,2)$.
Differentiate this Taylor polynomial and contract with $I_T/2$.
The cubic terms show why $C_3$ cannot be replaced by the accidental
$O(2)$ symmetry of the isolated uninserted block.
\end{proof}
The divided-difference coefficient of that isolated block is already
Proposition~\ref{p3c:prop:v1101-terminal-edge-all-orders}. The actual
$H$-boundary calculation must retain its source, metric, localizer,
cutoff and stationary-reduction contacts. This proposition neither
supplies a missing physical pullback nor equates the edge quarter to
a homogeneous response decrement.


\section{Berger semigroups, normalized visibility, and fixed-seed heat}
\label{p3c:sec:v2216-collapse-visibility}

\begin{theorem}[Sharp Berger semigroup and resolvent convergence]
\label{p3c:thm:v2216-berger-norms}
For $0<\varepsilon\le1$, identify the Berger scalar $L^2$ spaces with
normalized Haar $L^2(S^3)$ by the constant-density unitary.  Let
$\Pi_{\mathrm{bas}}$ be the Hopf-circle average and let
$\Delta_{\mathrm{bas}}$ be the base Laplacian on $S^2(\Rmom/2)$,
pulled back to its range.  Put
$\Lambda_\varepsilon=(2+\varepsilon^{-2})/\Rmom^2$.
For heat parameter $t>0$ and resolvent parameter $z>0$ in their
respective physical units,
\begin{align}
 \|e^{-t\Delta_\varepsilon}
       -e^{-t\Delta_{\mathrm{bas}}}\Pi_{\mathrm{bas}}\|
 &=e^{-t\Lambda_\varepsilon},\nonumber\\
 \|(\Delta_\varepsilon+z)^{-1}
       -(\Delta_{\mathrm{bas}}+z)^{-1}\Pi_{\mathrm{bas}}\|
 &=\frac1{z+\Lambda_\varepsilon}.
 \label{p3c:eq:v2216-berger-norms}
\end{align}
On a finite complete charge-invariant band these are upper bounds,
with equality exactly when a minimizing $k=1,q=\pm1$ mode is retained.
\end{theorem}
\begin{proof}
The inherited commuting resolution has eigenvalues
$[k(k+2)+(\varepsilon^{-2}-1)q^2]/\Rmom^2$.
The $q=0$ sector is unchanged and is the base spectrum.  For $q\ne0$,
$k(k+2)-q^2\ge2|q|\ge2$ and
$\varepsilon^{-2}q^2\ge\varepsilon^{-2}$, with equality exactly
at $k=1,q=\pm1$.  Both difference operators vanish on the basic
sector and are diagonal on its complement.  Taking their spectral
suprema proves \eqref{p3c:eq:v2216-berger-norms} and the finite-band
statement.
\end{proof}

\begin{corollary}[Exact basic exclusion and its scope]
\label{p3c:cor:v2216-berger-cutoff}
Under the closed cutoff $\lambda\le\Lcut^2$, all retained scalar
modes are basic whenever $2+\varepsilon^{-2}>\Rmom^2\Lcut^2$.
The strict inequality respects the included threshold.  The projected
resolvent limit acts on the surviving subspace and is not a resolvent
of a densely defined operator on the whole original Hilbert space.
In particular $\Pi_{\mathrm{bas}}$ is infinite rank before cutoff;
it is not the protected finite projector $P_0$.  No charge-to-$W$
intertwiner or physical sequence of seven events follows from this limit.
\end{corollary}
\begin{proof}
Every nonbasic eigenvalue is at least $\Lambda_\varepsilon$.
The projected limit vanishes on the nonbasic complement and hence has
a nontrivial kernel, whereas a resolvent on the whole space is injective.
The remaining distinctions follow from the stated ranges and inherited
owner domains.
\end{proof}

\begin{theorem}[Normalized visibility has no fixed finite linear generator]
\label{p3c:thm:v2216-visibility-generator}
On the complexification of the protected $1\oplus7$ carrier put
\[
 V_\tau=P_0+e^{-\tau}P_W,\qquad
 \rho_\tau=\frac{V_\tau}{1+7e^{-\tau}},\qquad \tau\ge0.
\]
There is no time-independent linear endomorphism $\mathcal L$ of its
finite matrix space such that $e^{\tau\mathcal L}(I/8)=\rho_\tau$
for every $\tau\ge0$.  Thus no homogeneous CPTP semigroup generates
this precise normalized curve with a clock affine in $\tau$.
\end{theorem}
\begin{proof}
The scalar matrix element
$\Tr(P_0e^{z\mathcal L}(I/8))$ is entire in $z\in\mathbb C$.
On the positive real axis it would equal $(1+7e^{-z})^{-1}$,
which has genuine poles at $z=\log7+(2r+1)\pi i$.
The identity theorem on their connected complement would force equality
there, contradicting entirety at a pole.  A constant positive clock
rescaling leaves this obstruction intact.
\end{proof}

\begin{proposition}[Filter and two CPTP realizations of the same state curve]
\label{p3c:prop:v2216-visibility-realizations}
The single-Kraus maps $\mathcal F_\tau(A)=K_\tau A K_\tau$ with
$K_\tau=P_0+e^{-\tau/2}P_W$ form a CP trace-nonincreasing
semigroup.  On $I/8$ their success probability is
$(1+7e^{-\tau})/8$ and their normalized state is $\rho_\tau$.
Alternatively define
\[
 \mathcal A(\rho)=\sum_{r=1}^7|0\rangle\langle r|\rho
                    |r\rangle\langle0|-\tfrac12\{P_W,\rho\}.
\]
The nonautonomous Lindblad equation
$\rho'=\gamma(\tau)\mathcal A(\rho)$,
$\gamma(\tau)=(1+7e^{-\tau})^{-1}>0$, sends $I/8$ exactly to
$\rho_\tau$ and has CP-divisible CPTP propagators.  The constant-rate
semigroup $e^{gt\mathcal A}$, $g>0$, gives the same curve under
\[
 \tau(t)=gt+\log(8-7e^{-gt}),\qquad
 t(\tau)=\frac{\tau+\log(1+7e^{-\tau})-\log8}{g}.
\]
These constructions select no physical clock or environmental coupling.
\end{proposition}
\begin{proof}
$K_\tau^*K_\tau=V_\tau\le I$ proves the filter claim and the
commuting projections give its semigroup law.  Normalization produces
the displayed state, but is nonlinear on general inputs because success
depends on the input.  For amplitude damping the invariant population
$p=\Tr(P_0\rho)$ obeys $p'=\gamma(1-p)$.
The target $p=(1+7e^{-\tau})^{-1}$ satisfies $p'=p(1-p)$,
so the positive rate proves the second realization.  Its accumulated
rate is $\tau+\log(1+7e^{-\tau})-\log8\ge0$, proving CPTP
propagation.  At constant rate the excited eigenvalue is $e^{-gt}/8$;
equating it to $e^{-\tau}/(1+7e^{-\tau})$ proves the clock formulas.
Their derivative $8g/(8-7e^{-gt})$ decreases from $8g$ to $g$, so this
is not the native linear clock calibration.  The state curve therefore
does not choose among a success-conditioned filter, a time-dependent
rate, and a reparameterized homogeneous channel.
\end{proof}

\subsection{Exact scalar depth on the fixed momentum seed}
\label{p3c:sec:v2218-fixed-momentum-depth}

The microscopic \emph{generated seed} used in the following model is
\(\mathcal H_1\) at level one, with fixed radius \(\Rmom\). A change of spectral resolution or response on this carrier
does not change either its dimension or its constituted geometry.
The following construction derives a scalar interpolation directly from
the finite heat spectrum and the Born role weight. It does not require
the middle profile of the protected visibility process to be reproduced.

\begin{proposition}[Exact finite seed heat spectra and the square owner]
\label{p3c:prop:v2218-fixed-seed-heat-spectra}
In the sector order
\[
 (0,0),\quad(\tfrac12,\tfrac12),\quad(1,0),\quad(0,1),
 \quad(\tfrac12,0),\quad(0,\tfrac12),
\]
the fixed carrier has complex dimensions \((1,4,3,3,2,2)\).
I-A's Casimir \(K_1=C_L+C_R\) and precise retained Levi--Civita
square \(K_{\rm LC}=J_{\Rmom}^*\mathscr D_{\rm LC}^{\,2}J_{\Rmom}\)
obey
\begin{align}
 K_1&\sim(0,3/2,2,2,3/4,3/4),\nonumber\\
 \Rmom^2K_{\rm LC}&\sim(0,3,4,4,9/4,9/4).
 \label{p3c:eq:v2218-fixed-seed-spectra}
\end{align}
For the dimensionless heat parameter \(\theta\ge0\), their exact
traces in the full-module and Levi--Civita-square conventions are
\begin{align}
 Z_{\rm full}(\theta)
 &:=\Tr_{\mathcal H_1}e^{-2\theta K_1}
 =1+4e^{-3\theta/2}+4e^{-3\theta}+6e^{-4\theta},\nonumber\\
 Z_H(\theta)
 &:=\Tr_{\mathcal H_1}e^{-\theta\Rmom^2K_{\rm LC}}
 =1+4e^{-9\theta/4}+4e^{-3\theta}+6e^{-4\theta}.
 \label{p3c:eq:v2218-fixed-seed-heat-traces}
\end{align}
The compressed first-order operator
\(D_F=J_{\Rmom}^*\mathscr D_{\rm LC}J_{\Rmom}\) instead has
\[
 D_F^2\sim\Rmom^{-2}(0,0,0,0,9/4,9/4),\qquad
 \Tr e^{-\theta\Rmom^2D_F^2}=11+4e^{-9\theta/4}.
\]
Thus \(D_F^2\ne K_{\rm LC}\): their kernels have dimensions eleven
and one respectively. The statements about \(Z_H\) below use the
specified positive square \(K_{\rm LC}\), not an inference from its
principal symbol or from compressing the first-order operator first.
\end{proposition}
\begin{proof}
I-A Theorem~\ref{p1:thm:s3one-casimir-scale-space} gives the Casimir
eigenvalues and multiplicities. Corollary~\ref{p1:cor:canonical-retained-dirac-square}
gives the second row, and summing the exponentials with the displayed
sector dimensions proves both traces. The first-order compression and
its leakage are I-A Theorem~\ref{p1:thm:v214-first-order-retained-closure}.
Its twenty-five-dimensional invariant smooth closure is a comparison
construction; it does not replace the physical fifteen-dimensional seed.
\end{proof}

\begin{theorem}[Born-weighted heat defect and its finite crossings]
\label{p3c:thm:v2218-born-heat-depth}
On the specified finite heat branch define the scalar readout
\begin{equation}
 \boxed{q_H(\theta):=\frac12\Tr_{\mathcal H_1}
       (I-e^{-\theta\Rmom^2K_{\rm LC}})
 =7-2e^{-9\theta/4}-2e^{-3\theta}-3e^{-4\theta}.}
 \label{p3c:eq:v2218-born-heat-depth}
\end{equation}
The coefficient \(1/2\) is the Born role weight in III-B's
Definition~\ref{p3:def:physical_spectral_action}; it is not a conversion
of a complex state into a real protected channel. This scalar is
strictly increasing from \(q_H(0)=0\) to the terminal limit \(7\).
Each \(j=1,\ldots,6\) has exactly one finite crossing \(\theta_j\),
with
\begin{equation}
 \frac14\log\frac7{7-j}\le\theta_j
 \le\frac49\log\frac7{7-j}.
 \label{p3c:eq:v2218-born-depth-crossing-bounds}
\end{equation}
The associated scalar trace clock
\[
 \tau_H(\theta):=-\log\frac{Z_H(\theta)-1}{14}
\]
satisfies \(q_H=7(1-e^{-\tau_H})\), but is not additive in
\(\theta\). In particular
\[
 \tau_H'(0)=45/14,\qquad \tau_H''(0)=-213/392.
\]
These conclusions use fixed \(m_p\) and \(\Rmom\) throughout.
\end{theorem}
\begin{proof}
There are fourteen excited states in the specified square. The Born
weight gives terminal defect \(14/2=7\), and the trace formula gives
\[
 q_H'=\tfrac92e^{-9\theta/4}+6e^{-3\theta}
                         +12e^{-4\theta}>0.
\]
The endpoint limits, continuity and strict monotonicity prove existence
and uniqueness of every finite crossing. The three positive gaps lie
between \(9/4\) and \(4\), so
\[
 7e^{-4\theta}\le7-q_H(\theta)\le7e^{-9\theta/4}.
\]
Evaluation at a crossing gives the bounds. Finally
\(\tau_H'\) is the weighted mean of those gaps with positive weights
proportional to their heat multiplicities, and
\(\tau_H''\) is minus their variance. At zero these give the
displayed values; the nonzero variance excludes additivity.
\end{proof}

\begin{corollary}[Neutral hierarchy pullback of the fixed-seed scalar]
\label{p3c:cor:v2218-born-depth-neutral-character}
Fix the positive determinant-shell member
\(\alpha_\sigma=e^{L_1^{(\sigma)}}>1\) and \(w>0\).
The existing neutral depth character admits the explicit scalar pullback
\begin{equation}
 D_H(\theta):=\alpha_\sigma^{q_H(\theta)},\qquad
 \boxed{\Xi_{H,w}(\theta):=D_H(\theta)^{-w}
                         =\alpha_\sigma^{-wq_H(\theta)}.}
 \label{p3c:eq:v2218-born-depth-neutral-character}
\end{equation}
It is strictly decreasing from \(1\) to
\(\alpha_\sigma^{-7w}\), and
\(\Xi_{H,w}(\theta_j)=\alpha_\sigma^{-wj}\).
No fitted individual gap, threshold weight, or interpolation function
is introduced in this scalar construction. For a reference \(X_0>0\),
its neutral scaled readout is \(X_0\Xi_{H,w}\); at \(w=1\) and
\(X_0=U\) its endpoint values are \(U\) and \(U\alpha_\sigma^{-7}\).
This is a neutral scalar continuation. Its physical action-coefficient
use still consumes Corollary~\ref{p3c:cor:owner-typed-physical-dual-scale}
and the actual-action certificate, rather than following from the
endpoint values alone.
\end{corollary}
\begin{proof}
The determinant-shell logarithm, scalar multiplication, exponentiation
and positive dualization are the constructors already used in
Theorem~\ref{p3c:thm:generated-dual-depth-hierarchy}. Pull them back
along the explicitly constructed scalar \(q_H\). The signs and values
follow from the preceding theorem. This forms a scalar descendant,
without identifying its source Hilbert space with the protected
Clifford operator space or placing an unverified physical coefficient
in the determinant owner.
\end{proof}

\begin{proposition}[Carrier heat and protected operator visibility]
\label{p3c:prop:v2218-protected-heat-commutant}
On \((\mathrm{ND}_{\rm sp})\), represent the protected real even
algebra \(A_{\rm ev}\cong\Cl(0,3)\) on
\(\mathcal S_{\rm prot}=(\tfrac12,0)\oplus(0,\tfrac12)\).
Both operators in \eqref{p3c:eq:v2218-fixed-seed-heat-traces} are
scalar on that entire four-complex-dimensional block. Consequently
their unitary adjoint evolution fixes every protected operator, while
their symmetric heat sandwich attenuates every protected operator,
including its identity, by one common factor. Neither operation is the
unital visibility \(V_\tau=P_0+e^{-\tau}P_W\) on
\(\mathbb R1\oplus W\).
\end{proposition}
\begin{proof}
The restrictions are
\(\Rmom^2K_{\rm LC}|_{\mathcal S_{\rm prot}}=(9/4)I\) and
\(2K_1|_{\mathcal S_{\rm prot}}=(3/2)I\).
Thus \([K,A]=0\) for every represented \(A\in A_{\rm ev}\), and
\(e^{-\theta K/2}Ae^{-\theta K/2}=e^{-\theta\lambda}A\)
for the corresponding dimensionless generator and eigenvalue
\(\lambda\). On \(W\) the real superoperator trace is
\(7e^{-\theta\lambda}\), but the scalar is attenuated as well;
on the full octet the trace is \(8e^{-\theta\lambda}\), rather
than \(1+7e^{-\tau}\). Normalizing the sandwich to fix its identity
instead makes the entire protected algebra stationary. This excludes
these particular induced-operator identifications, not a differently
constructed unital response map or the scalar pullback above.
\end{proof}

\begin{remark}[Fixed-seed scalar, protected depth and physical action]
\label{p3c:rem:v2218-fixed-seed-role-status}
The canonical scalar heat ground projector on \(\mathcal H_1\) has
complex rank one. The identity of the represented protected algebra is
instead the rank-four projector onto \(\mathcal S_{\rm prot}\).
The protected \(1\oplus7\) is an operator-space split, not a
Hilbert-state split of the fifteen-space. Equal endpoint scalar counts
therefore do not identify \(q_H\) with \(q_{\rm sp}\), their
operator processes or their clocks. Conversely, an operator isomorphism
is not required to define the scalar \(q_H\) and its neutral
character; its distinct middle profile is not an obstruction to that
construction.

On a physical finite action branch that uses \(D_p^2=K_{\rm LC}\),
the Born momentum summand has the exact value
\[
 S_p(\theta)=\tfrac12Z_H(\theta)=\tfrac{15}{2}-q_H(\theta).
\]
If its neutral curve is denoted by
\(R_H=U\alpha_\sigma^{-q_H}\), and
\(C=U\alpha_\sigma^{-7}\), elimination of the parameter gives
\begin{equation}
 \boxed{S_p(R_H)=\frac12+
                  \frac{\log(R_H/C)}{\log\alpha_\sigma}.}
 \label{p3c:eq:v2218-fixed-heat-log-value}
\end{equation}
This is a total finite action value, not already a local metric
coefficient. At fixed momentum operator, trace, radius and heat
parameter its direct reciprocal-position metric derivative is zero;
the complete action can still respond through its derived position,
joint-coupling or actual density ports. III-B's scalar-action compiler
and Theorem~\ref{p3b:thm:v2211-action-contact-certificate} test that
coefficient transport. PQTI likewise compares an actual normalized
second-variation port; it is not a conversion of half a complex-state
count into seven real Clifford channels.

Finally, the Cartan-twisted \(D_u\) on the retained smooth round
\emph{spatial} representative in III-F
Section~\ref{p3f:sec:m8td-spatial-operators} has its own infinite
Hilbert space and radius role. Its protected-zero and resolvent results
are not substituted for the fixed finite momentum spectra above.
No momentum-radius flow, new momentum carrier, or identification of
that spatial radius with \(\Rmom\) is made here.
\end{remark}

\subsection{Coordinate-word exports to the protected hierarchy}
\label{p3c:sec:v226-hierarchy-exports}
\begin{corollary}[The protected hierarchy uses coordinate-derived chiral projectors]
\label{p3c:cor:v226-hierarchy-projectors}
On the fixed intrinsic momentum fibre the separate protected chiral
projectors and geometric duality projectors are the degree-six words of
Theorem~\ref{p1:thm:v226-chirality}. Existing protected-corner group averages,
finite characters and determinant matrices therefore retain their exact
operator ancestry under these projectors. Any actually native matrix can
be tested by the dual coefficients of
Theorem~\ref{p1b:thm:v226-dual-coefficients}.
\end{corollary}
\begin{proof}
Insert the reconstructed projectors in the already licensed finite words.
They equal the existing projectors, so no character, corner rank or determinant
changes. The dual inverse is an identity on that same algebra.
\end{proof}
Coordinate-word degree six, the 137-dimensional native operation envelope,
and the protected seven-step hierarchy are different counts on different
sorts. The new word bound makes no identification between their indices or
physical scales. The hierarchy phase character, endpoint and response data
retain their original definitions. Optimal-radius coordinate readout alone
cannot determine protected chiral populations
(Theorem~\ref{p2c:thm:v226-chirality-radius}); an actual source/state provider
or another licensed measurement may supply them on its own domain.



\subsection{Finite protected hierarchy measurement-resource certificate}
\label{p3c:sec:v266-protected-resource}
\begin{corollary}[Protected projectors can be tested without complete tomography]
\label{p3c:cor:v266-protected-projector-resource}
On the retained native finite coordinate carrier, the central protected
projector $P_S$ and orientation-sensitive protected chirality
$\widetilde\chi=P_S\chi$ can each be reconstructed together with all
four exact intrinsic position moments by six outcomes at any assigned
radius $r\ge2\Rmom$; seven outcomes reconstruct them both in a single
setting. The required effects are explicit and positive, and their
outcome numbers are sharp relative to the indicated target operator
systems. Five exact-position outcomes are blind to both added targets.
For the normalized protected chiral family, the sharp position-compatible
chirality-information bound can be attained by seven outcomes for every
$r\ge\sqrt3\Rmom$ and by six for
$r\ge\sqrt{34}\Rmom/3$.
\end{corollary}
\begin{proof}
Apply Theorems~\ref{p2c:thm:v266-minimal-protected-readout} and
\ref{p2c:thm:v266-sharp-chirality-support} to the already constructed
protected Clifford and chirality operators. The position-label first
moment and positivity checks are internal finite mathematical results.
\end{proof}
No hierarchy mass, chiral source occupation, physical phase assignment,
protected-step transition rate, or same-action detector map follows
from this readout-feasibility certificate. In particular the abstract
mathematical ability to test a protected hierarchy projection is
logically distinct from deriving which hierarchy source populates it.



\section{Finite exceptional comparisons and protected source responses}
\label{p3c:sec:v196-finite-comparisons}
\subsection{Petersen two-core polar compatibility with the balanced Clifford transition}
\label{p3c:sec:v319-petersen-protected}
Only the actual two-dimensional Hopf tangent
$V_u=u^\perp\cap\operatorname{Im}\mathbb H$ may enter the existing
Clifford transition $S_u:V_u\to\operatorname{End}_{\mathbb R}\Cl(0,3)$
of Theorem~\ref{p3c:thm:v179-clifford-projector-tangent}.  A
conjugate $\mathbf3'$ finite representation is \emph{not} a new
Clifford generator or an occupied additional carrier.

\begin{theorem}[Exact two-core reciprocal congruence and preserved midpoint]
\label{p3c:thm:v319-petersen-midpoint}
For the two tetrahedral cores $i,j$ meeting at $u=u_{ij}$ let
$T_{ij}=J_i+J_j:\mathbf3\to\mathbf3'$.  On
$\operatorname{Im}\mathbb H$ set $Q_4=uu^{\mathsf T}$,
$Q_5=I_3-Q_4$ and
\[
  T_{ij}^{\mathsf T}T_{ij}=4Q_4+\frac14Q_5,\qquad
  R_{ij}:=(T_{ij}^{\mathsf T}T_{ij})^{1/2}
             =2Q_4+\frac12Q_5.
\]
For every $v\in V_u$, the off-diagonal tangent source
$S_u(v)$ obeys (on this block)
\[
 R_{ij} S_u(v)R_{ij}=S_u(v).
\]
The positive carrier grading restricted to these blocks transforms
by the algebraic congruence
\[
 R_{ij}(E_{32}Q_4+E_{40}Q_5)R_{ij}
      =4E_{32}Q_4+\frac14E_{40}Q_5,
\]
so the balanced transition factor remains
$\sqrt{(4E_{32})(E_{40}/4)}=E_{36}$.
\end{theorem}
\begin{proof}
The polar spectrum is the proved pair-subset statement of Paper~I-B,
Theorem~\ref{p1:thm:petersen-subset-polar}.  A transverse source
has only $(4,5)$ and $(5,4)$ matrix blocks; the two factors act
there by $2$ and $1/2$, so their product is one.  The carrier identity
and geometric mean are block multiplication.  This proof uses no
observable mass, field occupation or modified physical action.
\end{proof}

\begin{remark}[Congruence is not equivalence of physical actions]
\label{p3c:rem:v319-not-a-gauge-change}
$R_{ij}$ is positive and nonunitary.  The displayed transformation is
not a unitary similarity, not an allowed change of physical source
normalization by itself, and not an invariance theorem for the complete
Clifford or matter action.  Only the specified off-diagonal transition
factor is protected.  The source metric
$\tau_8(S_u(v)^\dagger S_u(w))=\langle v,w\rangle/4$
and the existing $E_{36}$ response stay as in
Theorem~\ref{p3c:thm:v179-balanced-hierarchy-response}.
\end{remark}

\begin{proposition}[Two- versus three-core auxiliary mirror block]
\label{p3c:prop:v319-petersen-complement-spectrum}
The selected two-core map $T=T_{ij}$ has singular values
$2,\tfrac12,\tfrac12$ and determinant $+\tfrac12$ in the
matched oriented quaternionic representations.  Its complementary
three-core map is $-T$ and has determinant $-\tfrac12$.
The \emph{auxiliary} self-adjoint six-dimensional block
\[
 \mathfrak D_T=\begin{pmatrix}0&T^{\mathsf T}\\T&0\end{pmatrix}
\]
has eigenvalues $\{\pm2,\pm\tfrac12,\pm\tfrac12\}$, with the
last pair each occurring twice.  Passing to $-T$ preserves the even
spectrum but reverses complement-odd quantities.  No fermion or
hadron Dirac operator is thereby constructed.
\end{proposition}
\begin{proof}
The polar identity gives the singular values.  Direct evaluation of
the five $J_m\in SO(3)$ gives $\det(J_i+J_j)=1/2$ and
$J_{\{1,\ldots,5\}\setminus\{i,j\}}=-T$.  The standard
symmetric off-diagonal block has eigenvalues equal to plus and minus
the singular values of $T$.
\end{proof}


\subsection{Octahedral-complement geometry of the balanced transition}
\label{p3c:sec:v189-octahedral-transition}
\begin{theorem}[Six maximizers of the existing balanced Clifford source]
\label{p3c:thm:v189-octahedral-e36-selector}
Fix an occurring local $H\cong2T$ and one of its four active axes $u$.
Use its ambient normalizer $2O$ and its complement $X=2O\setminus H$.
Set $P_u$ equal to orthogonal projection in
$\operatorname{Im}\mathbb H$ onto $u^\perp$ and define, on $X$,
\[
 \mathcal J_u(x):=E_{36}S_u(P_u\operatorname{Im}x).
\]
For $u=(i+j+k)/\sqrt3$ the six-element normalizer coset
$B_6(u)$ of the sixfold normalizer proposition in Paper~I-B
obeys
\[
 \boxed{\|\mathcal J_u(x)\|_{\tau_8}^2=
 \begin{cases}E_{36}^2/4,&x\in B_6(u),\\
 E_{36}^2/12,&x\in X\setminus B_6(u).
 \end{cases}}
\]
Thus $B_6(u)$ is exactly the set of six maximizers of this
source norm among the $24$ complementary vertices; the other
$18$ have one third the squared response.  Moreover
\[
 \sum_{x\in B_6(u)}(P_u\operatorname{Im}x)^{\otimes2}
 =\sum_{x\in X\setminus B_6(u)}(P_u\operatorname{Im}x)^{\otimes2}
 =3P_u.
\]
The two sectors have equal total transverse second-moment tensors.
\end{theorem}
\begin{proof}
The normalized source identity
$\|S_u(v)\|_{\tau_8}^2=\|v\|^2/4$ is
Theorem~\ref{p3c:thm:v179-clifford-projector-tangent}.
In the quaternion list $X=\{(\pm e_a\pm e_b)/\sqrt2:a<b\}$,
six pure-imaginary members lie in $u^\perp$ and have projection
norm squared one.  Every other member has projection norm squared
$1/3$.  The six maximizers are exactly $B_6(u)$; direct summation
of their squared coordinate vectors gives $3P_u$, and summation
of the remaining eighteen does the same.  The construction is
normalizer-covariant, so the result holds at each transported axis.
\end{proof}
\begin{remark}[Carrier energies, transition energies and physical status]
\label{p3c:rem:v189-octahedral-clock-scope}
The active carrier still has the previously derived spectrum
$\{E_0,E_{32},E_{40},E_{56}\}$, while
$E_{36}^2=E_{32}E_{40}$ weights its off-diagonal source.
Consequently the $6+12+6$ complementary octahedral partition is a
candidate \emph{particle response support} at the balanced transition,
not twenty-four new eigenvalues of $H_{\rm car}$.
The attached $C_6$ action and its six outer normalizer points define
the $E_{40}$ local tangent geometry; they are not six extra elements
outside $2O$.  Particle labels and mass/pole placements still require
the downstream generated physical source and kinetic normalization.
\end{remark}



\subsection{Sixfold moment and local harmonic comparisons}


\subsubsection{Exact fourth- and sixth-order resolution of the sixfold complement}
\label{p3c:subsec:v190-sixfold-higher-moment-bridge}

\begin{proposition}[First orientation-sensitive moment of the balanced sixfold source]
\label{p3c:prop:v190-sixfold-moment-hierarchy}
Fix the mathematical data and axis $u=(i+j+k)/\sqrt3$ of
Theorem~\ref{p3c:thm:v189-octahedral-e36-selector}. Put
$B=B_6(u)$ and $R=(2O\setminus 2T)\setminus B$; these have cardinalities
$6$ and $18$. For $x\in B\sqcup R$ let
$b_x=P_u\operatorname{Im}x\in u^\perp$, and define, for
$w\in u^\perp$ and integer $k\ge1$,
\[
 M_{2k}^{B}(w)=\sum_{x\in B}\langle b_x,w\rangle^{2k},
 \qquad
 M_{2k}^{R}(w)=\sum_{x\in R}\langle b_x,w\rangle^{2k}.
\]
Then the two sectors obey the exact identities
\[
\boxed{
 M_2^B=M_2^R=3\|w\|^2,
 \qquad
 M_4^B=\frac94\|w\|^4,
 \quad
 M_4^R=\frac34\|w\|^4.}
\]
Choose the positively oriented orthonormal coordinates
$e=(i-j)/\sqrt2$ and $f=(i+j-2k)/\sqrt6$ on $u^\perp$
(here $k$ in $f$ denotes the quaternionic unit), and write
$w=ze+tf$.  With $r^2=z^2+t^2$, the sixth moments are
\[
\boxed{\begin{aligned}
 M_6^B(w)&=\frac{15}{8}r^6+
      \frac{3}{16}\operatorname{Re}(z+it)^6,\\
 M_6^R(w)&=\frac5{24}r^6-
      \frac1{48}\operatorname{Re}(z+it)^6.
\end{aligned}}
\]
In particular,
\[
\boxed{
 M_2^B-M_2^R=0,\quad
 M_4^B-3M_4^R=0,\quad
 M_6^B-9M_6^R=\frac38\operatorname{Re}(z+it)^6.}
\]
The last term is a nonzero, sixth-order angular harmonic of the
$C_6$ normalizer interface. The identities are directional moments
of the finite comparison vectors, \emph{not} moments of a physical
meson, baryon, or decay-rate distribution.
\end{proposition}
\begin{proof}
The explicit $24$ quaternion list in the octahedral
proposition gives $\|b_x\|^2=1$ on $B$ and $1/3$ on $R$.
Both finite sets, regarded in $u^\perp$, are invariant under rotation
through $2\pi/3$ and under inversion. Thus their even moment
polynomials of degrees two and four are radial, since there are no
nonconstant planar Fourier modes of these degrees divisible by three.
The angular averages of $\langle b,w\rangle^2$ and
$\langle b,w\rangle^4$ are $\|b\|^2\|w\|^2/2$ and
$3\|b\|^4\|w\|^4/8$. Inserting the sector sizes and norms gives
$3,3$ and $9/4,3/4$ respectively.
At degree six the first nonradial invariant is
$\operatorname{Re}(z+it)^6$. Direct expansion of the same explicit
quaternion list in the displayed $(e,f)$ basis gives the two sixth
moments; subtracting nine times the second eliminates their radial
part and leaves the stated hexagonal harmonic. All formulas transport
covariantly to any conjugate axis.
\end{proof}

\begin{remark}[Scope and relationship to the five-frame Petersen sextic]
\label{p3c:rem:v190-sixfold-owner-frontier}
The local angular degree-six response is not automatically the
$A_5$-invariant global Petersen sextic
$H_6^P$ of Proposition~\ref{p3c:prop:v1188-fiveframe-quartic-cancellation}.
Identifying those two sextics requires an explicitly generated
axis-to-five-core aggregation/intertwiner, including its sign and
orientation choices. The present equal-second-moment identity also
does not imply equal decay branching fractions: dynamical matrix
elements, pole residues, phase space and channel completeness are
separate physical inputs.
\end{remark}


\subsection{Four-core chirality contrasts and a second factorisation of the ten-axis $E_{36}$ Gram}
\label{p3c:subsec:cc12-chiral-gram-factorization}
\begin{corollary}[Four-core chirality contrasts factor the averaged ten-axis Clifford Gram]
\label{p3c:cor:cc12-e36-chiral-factorization}
Let $\mathcal J(V)=\operatorname{Herm}_3(\mathbb C)$ be the
\emph{auxiliary matrix space} associated to the generated
$V=\operatorname{Im}\mathbb H\simeq\mathbb R^3$, with its real
Hilbert--Schmidt metric and canonical $\mathbf1\oplus\mathbf3\oplus\mathbf5$
spatial decomposition. No downstream hadronic source is imported.
On its imaginary-skew summand write $B(v)=iA(v)$, with
$A(v)w=v\times w$; let $T:\mathcal J(V)\to V$ be
the isometric triplet extraction $T(aI+Q+iA(v))=\sqrt2v$,
and let $\Pi_3=T^*T$ denote the orthogonal projector
onto that summand. Select one of the five four-axis tetrahedral cores
in the already existing $2I$ packet and its four directed unit
vectors $n_j$. Define the \emph{centered circular contrasts}
\[
 C_j:=\tfrac12iA(n_j)\in\operatorname{Herm}_3(\mathbb C).
\]
Their finite Hilbert--Schmidt frame is
\[
 \boxed{\sum_{j=1}^4 |C_j\rangle\langle C_j|
       =\tfrac23\Pi_3.}
\]
Independently, the normalized Clifford source of
Theorem~\ref{p3c:thm:v179-clifford-projector-tangent}, applied
to the ten projective axes and extracted isometrically on the
matching three-space, has averaged comparison Gram
\[
 G^{(9)}_{36}:=\frac1{10}\sum_{a=1}^{10}
  \frac{E_{36}^2}{4}\,T^*(I-u_au_a^{\mathsf T})T
  =\frac{E_{36}^2}{6}\Pi_3.
\] Consequently the \emph{same} rank-three
comparison Gram has two exact factorizations:
\[
 \boxed{G^{(9)}_{36}=\frac{E_{36}^2}{4}
  \sum_{j=1}^4|C_j\rangle\langle C_j|.}
\]
The quadratic source metric is independent of which tetrahedral core
is marked and of the simultaneous reversal $n_j\mapsto-n_j$.
Linear orientation-sensitive readout coefficients, however, reverse
under this global hand change; an actual chirality detector would
need its own orientation/source reference.
\end{corollary}
\begin{proof}
The four local unit vectors form a regular tetrahedron, so
$\sum_jn_jn_j^T=\tfrac43I_3$. On $B(v)$ one has
$\operatorname{Tr}(B(n)B(v))=2n\cdot v$, and therefore
\[
 \sum_j\langle C_j,B(v)\rangle C_j
 =\tfrac12 B\!\left(\sum_j(n_j\cdot v)n_j\right)
 =\tfrac23B(v).
\]
Every $C_j$ is orthogonal to the real symmetric $\mathbf1\oplus\mathbf5$,
proving the first identity. The Clifford trace normalization
$\tau_8(S_u(v)^\dagger S_u(w))=\tfrac14v\cdot w$ and the ten-axis
resolution $\sum_au_au_a^T=\tfrac{10}{3}I_3$ imply
$\tfrac1{10}\sum_a(I-u_au_a^T)=\tfrac23I_3$.
This yields the second Gram identity and its factorization.
\end{proof}

\begin{remark}[No same-action identification]
The equality compares two positive quadratic forms on the already generated
mathematical $V$. Paper VII-B may later transport it through an
independently supplied nonet isometry; the present result does not
consume a later-paper theorem as a premise. It does not identify the auxiliary qutrit
measurement with the occurring $Q_4$--$Q_5$ physical transition,
with an actual meson source, or with a decay branching fraction.
\end{remark}

\subsection{Exact Hopf/Petersen-to-icosian quadratic comparison}
\label{cc19:sec:petersen-hopf-complement}
\begin{theorem}[Petersen sextic as a right-averaged Hopf lift]
\label{cc19:thm:petersen-hopf-lift}
Let $H_6^P$ denote the source's degree-six $A_5$-invariant Petersen
spherical harmonic on the unit $S^2$. Mark the aligned fivefold
icosahedral vertex $n$ for which $H_6^P(n)=-32/189$, and the Hopf map
$\pi_n(q)=q n q^{-1}$. With $P_I^R$ the Haar-normalized right
$2I$ average,
\[
\boxed{P_I^R(H_6^P\circ\pi_n)=-\frac{32}{189}C_{12}.}
\]
Moreover, if $B_8=\sum_{a=1}^5Q_a^2$ is the already-established
five-octahedral-frame squared-quartic comparison on the unit $S^2$,
\[
\boxed{P_I^R(B_8\circ\pi_n)=\frac{16}{105}(1-C_{12}),}
\]
\[
\boxed{\rho_6+\frac{315}{4}P_I^R(B_8\circ\pi_n)=28.}
\]
These are mathematical comparison identities, not physical
$E_{36}$ action or conservation equations.
\end{theorem}
\begin{proof}
The Hopf lift of the scalar $S^2$ degree-six harmonic is a degree-twelve
round scalar harmonic on $S^3$ since its eigenvalue is
$4\cdot6\cdot7=12\cdot14$. It is left-$2I$ invariant. Right-$2I$
averaging therefore gives a bi-$2I$-invariant degree-twelve scalar
harmonic and hence a multiple of the unique normalized line $C_{12}$.
Evaluate at $q=1$: the average of $H_6^P(gng^{-1})$ is $H_6^P(n)$
because the sextic is invariant under the $A_5$ rotation subgroup.
Paper III-C's exact quartic identity
$B_8=16/105+(9/10)H_6^P$ gives the second formula.
Substitute $\rho_6=16+12C_{12}$ for the third.
\end{proof}



\subsection{Exceptional nearest-shell Gram on protected Clifford response edges}
\label{p3c:sec:et-shell-gram}

\begin{proposition}[Exact exceptional-shell Gram comparison; source-edge envelope]
\label{p3c:prop:et-shell-edge-gram}
Fix a unit marked imaginary Hurwitz direction $u$, its horizontal plane
$H_u=u^\perp$, and the complete \emph{unnormalized} nearest-shell unit
vectors $S_\Gamma$ for $\Gamma\in\{2T,2O,2I\}$.
Let $J_{45,u}(v)$ and $J_{57,u}(v)$ be the self-adjoint protected
$Q_4\leftrightarrow Q_5$ and $Q_5\leftrightarrow Q_7$ Clifford
response edges, respectively, on $v\in H_u$, with the original
$\tau_8=\operatorname{Tr}_{\mathbb R}/8$ source normalization.  Define
\[
 \mathcal N_\Gamma^e:=\sum_{a\in S_\Gamma}
 \tau_8\bigl(J_{e,u}(P_{H_u}a)^*J_{e,u}(P_{H_u}a)\bigr),
 \qquad e\in\{45,57\}.
\]
For either edge, as an exact finite mathematical comparison,
\[
 \boxed{\mathcal N_{2T}^e=\frac43,\quad
 \mathcal N_{2O}^e=1,\quad \mathcal N_{2I}^e=2.}
\]
Their \emph{balanced squared-edge} versions carry the separate factors
$E_{36}^2$ and $E_{48}^2$, respectively; they are not themselves the
physical gravitational $R^2$ coefficients.  In particular, the
unbroken terminal $C_3$ branch does not acquire a linear $E_{48}$
source from this quadratic comparison.
\end{proposition}
\begin{proof}
The protected Clifford edge Gram is
$\tau_8(J_{e,u}(v)^*J_{e,u}(w))=g_E(v,w)/4$ for each edge.
Paper~I-B, Theorem~\ref{p1b:thm:v11-three-shell-frames}, yields
$\sum_{a\in S_\Gamma}P_Ha\otimes P_Ha=(n_\Gamma/3)I_H$,
where $(n_T,n_O,n_I)=(8,6,12)$ and $\dim H=2$.
Consequently
$\mathcal N_\Gamma^e=\tfrac14\operatorname{tr}_H((n_\Gamma/3)I_H)
=n_\Gamma/6$.
A balanced source edge is multiplied by $E_{36}$ or $E_{48}$,
so its squared Gram has the corresponding energy-square weight.
\end{proof}

\begin{remark}[Antipodal incidence and the failed direct energy identification]
\label{p3c:rem:et-edge-incidence-firewall}
For each complete antipodal shell, linearity gives
$\sum_{a\in S_\Gamma}J_{e,u}(P_Ha)=0$, whereas the preceding
quadratic shell sum is strictly positive.  Hence a coherent
first-moment source cannot be replaced by the second-moment covariance
without specifying the actual incidence before source elimination.
Moreover, assigning the raw balanced $E_{36}^2,E_{48}^2$ source
norms directly to the proposed UV/IR coefficients, with reciprocal
$\delta_{\rm loc}^{\mp1}$ factors, gives
\[
 \frac{\delta_{\rm loc}^{-1}(4/3)E_{36}^2}
 {\delta_{\rm loc}(2)E_{48}^2}
 =48\alpha_{\rm hier}^{3/4},
\]
not $64\alpha_{\rm hier}^{-3}$.  The octet reflection
$K\mapsto64-K$ and its $E_8\leftrightarrow E_{56}$ identity are
mathematical hierarchy relations, not an automatic physical action
rebase.  No gravitational coupling is promoted by this result.
\end{remark}


\subsection{Octahedral-frame decomposition of the homogeneous Milnor shear}
\label{bh:berger:v117:p3c:fiveframe-shear}
\begin{proposition}[Five-frame quintet cancellation]
\label{bh:berger:v117:p3c:projectors}
On $V_5=\operatorname{Sym}^2_0(\mathbb R^3)$ the real octahedral
restriction decomposes as $\mathbf5|_O=\mathbf2\oplus\mathbf3$.
Let $P_2^{(i)},P_3^{(i)}$ be its orthogonal isotypic projectors for
the five already fixed conjugate octahedral frames in the marked
icosahedral comparison. Then
\[
 \sum_{i=1}^5 P_2^{(i)}=2I_5,\qquad
 \sum_{i=1}^5 P_3^{(i)}=3I_5.
\]
For one \emph{actually selected} octahedral-invariant matter state,
its shape Hessian has the symmetry-permitted form
$K=k_2P_2+k_3P_3$; the values $k_2,k_3$ are not fixed by
representation theory alone.  Equally weighted frame averages are
proportional to the identity and cannot lock a direction quadratically.
\end{proposition}
\begin{proof}
Conjugating a rank-two projector by the five coset representatives
and summing produces an $A_5$-invariant endomorphism of the real
irreducible quintet.  Schur's lemma makes it scalar. The trace
is $5\cdot2=10$, so the scalar is $2$; the complementary projector
sum is $5I_5-2I_5=3I_5$.
\end{proof}
\begin{remark}[Quartic versus sextic source information]
\label{bh:berger:v117:p3c:source-scope}
In a specified local frame the difference between the two quadratic
shape responses has the established octahedral quartic angular
profile. For the complete five-frame comparison those profiles cancel,
leaving the icosahedral sixth-order shape as the earliest possible
collective orientation comparison.  Neither a nonzero difference
$k_2-k_3$ nor a physically populated frame follows from this algebra.
\end{remark}




\subsection{Protected $E_{36}$ response composed with the finite Petersen source}
\noindent\emph{Scope.} The statements below retain their declared finite geometric, source-chart, and action-relative hypotheses. They do not assert physical particle decay or unconditional source occupation.
\begin{corollary}[Source-visible protected $E_{36}$ comparison response]
\label{p3c:cor:v368-petersen-chamber-e36-response}
With the source isometry $C_E$ of Theorem~\ref{p1b:thm:v368-petersen-chamber-rotation},
and under the existing flux identity $KK^*=5I_3$ and the normalized
protected Clifford source
$\mathcal J_u=-\frac23E_{36}S_u(I-uu^*)K$, define
$\mathcal D_u=\mathcal J_uC_E^*$, on the common selected
Petersen source basis. Then
\[
 \mathcal D_u(z-A)^{-1}\mathcal D_u^*
 =\frac{z+2}{z^2+z-5}\mathcal J_u\mathcal J_u^*,
 \quad
 \mathcal J_u^*\mathcal J_u=\frac{E_{36}^2}{9}K^*(I-uu^*)K.
\]
The Gram has rank two and both nonzero eigenvalues $5E_{36}^2/9$.
For the existing ten-axis tight frame $\sum_a u_au_a^*=10I_3/3$,
\[
 \frac1{10}\sum_a\mathcal J_{u_a}^*\mathcal J_{u_a}
 =\frac{2E_{36}^2}{27}K^*K,
 \quad\operatorname{Spec}=\{(10E_{36}^2/27)^{(3)},0^{(12)}\}.
\]
This is a comparison response, not an action-owned vertex, single-spinor
state, hadron propagator, or physical decay rate.
\end{corollary}

\begin{proposition}[Orbit derivative and balanced Clifford response are distinct]
\label{p3c:prop:v368-orbit-balanced-distinction}
On the previously constructed positive carrier grading $H_{\mathrm{car}}$
and its protected $E_{32}$--$E_{40}$ transverse source $S_u(v)$,
\[
 [J_v,H_{\mathrm{car}}]=(E_{32}-E_{40})S_u(v),\qquad
 \mathscr E_H(S_u(v))=E_{36}S_u(v),\quad E_{36}^2=E_{32}E_{40}.
\]
Consequently, on the declared seven-step hierarchy branch,
\[
 \frac{E_{32}-E_{40}}{E_{36}}
 =\alpha_{\mathrm{hier}}^{1/8}-\alpha_{\mathrm{hier}}^{-1/8}.
\]
The first coefficient is an infinitesimal orbit derivative, while the
second belongs to the balanced positive-energy map. Their comparison
does not constitute equality of a geometric tangent with a physical
weak gauge-field source derivative.
\end{proposition}
\begin{proof}
In the two common/relative grading blocks the commutator
of an off-diagonal generator with the diagonal carrier energy
multiplies the off-diagonal matrix by the difference of the two
corresponding energies. The existing balanced map multiplies the
same source by their geometric mean. Using the inherited adjacent
hierarchy powers yields the displayed ratio.
\end{proof}

\section{Hierarchy outputs and physical incidence limits}
\label{paper3c:sec:organized-conclusion}
\subsection{Conclusion}
\label{p3c:sec:conclusion}
Paper~III-C has computed the normalized principal unit-Chern Hopf/contact residue, the
\((\mathrm{REG})\)-specific smooth determinant correction, the generated
active-chain orbit and rank-seven endpoint, and the exact finite-window Dirac formulas.
The normalized contact operator is
\(A_{\rm H}=8\Delta_{b,0}^{-2}\).  The generated level-one momentum seed \(\mathcal H_1\) is imported
from Paper~II-B and is not reselected here; universal physical-endpoint
multiplicity-one is not an antecedent of the protected hierarchy.

The protected physical Clifford formulas require only the nondegenerate
stratum \((\mathrm{ND}_{\rm sp})\).  On an occurring Hurwitz-even
Hopf/contact branch, equivariant finite functional calculus derives the
complete orbit of support chains \(8\to4\to3\to1\) at \(q=0,4,5,7\).
A particular diagonal-dependent chain remains marked, while all members have
the same ranks, spectra, and terminal scalar Haar projector.  Thus
\((\mathrm{IR}_{O(3)})\) is a certification, not
an additional extraction assumption.  The endpoint is
\(\Cl(0,3)=\mathbb R\mathbf1\oplus W\) with \(\dim_{\mathbb R}W=7\).  Its
choice-free common-rate semigroup has one sevenfold-degenerate non-scalar
block.  Its physical process role requires no separate occurrence because the
explicit forgetful constructor satisfies the current G0 witness-transport,
essential-image, and zero-joint-defect clauses.  It remains distinct from each
marked multi-rate response.  Its natural parameter is dimensionless;
identification with physical time uses the separate calibration
\((\mathrm{CAL}_{\rm vis})\).  Its scalar depth coordinate and dual
determinant character derive seven neutral ratios.  They are decreasing on
the exact positive one-step scheme domain, including \(\mathrm{REG}_0\);
seven successive rank-one projector events require
\((\mathrm{FLAG})\).  The unique invariant additive defect on the common
\(1\oplus7\) owner is the trace defect; it restricts to projector codimension
on the hard chain and to \(7(1-e^{-\tau})\) on the smooth flow.
Corollary~\ref{p3c:cor:visibility-universal-continuation-instance} identifies
this as the first complete universal-continuation instance: its admissible
generator cone is the one ray \(\mathbb R_{\ge0}(-P_W)\), the transported
decay witness excludes zero, and only the declared visibility-rate
calibration remains.  The same architecture does not identify this process
with a different sector parent or rate.
The upper
complementary trace is \(N_W=7e^{-\tau}\), so its unique primitive finite
residual occurs at \(N_W=1\), equivalently
\(\tau=\ln7\), \(q_{\rm sp}=6\), and \(\Ksh=48\).  This neutral trace result
does not create a current, source, transfer history, or cosmological event.
The upper
spectral resolution of \(\mathcal Q_{\rm act}\) therefore gives a canonical
dimensionless, orbit-equivariant step filtration and four exact defect
anchors without \((\mathrm{ACT}_{\rm Cl})\).  It is not an
operator-continuous rank-changing path, does not choose a diagonal, and does
not actualize physical transitions.  An alternative or external
scale-dependent projector law retains the additional
\((\mathrm{ACT}_{\rm Cl})\) interface.  The visibility flow's extensive value
\(1+7e^{-\tau}\) is a closed heat loop, whereas its unit-trace normalization
is a boundary/state-conditioned object.
The resolved trace and that full loop satisfy
\(q_{\rm sp}+Z_W=8\).  Hence the generated scalar levels carry the involution
\(j\mapsto8-j\), and every already owner-typed dual scale obeys the midpoint
product law \(X_{w,j}X_{w,8-j}=X_{w,4}^2\).  In shell notation this is
\(K\mapsto64-K\), with fixed point \(K=32\).  It neither reverses the
contractive flow nor supplies physical event predicates.

The smooth value
\[
        \ddef=\delta_{\rm loc}=\frac1{6\sqrt2}
\]
is unique only inside each fixed member of the displayed determinant-shell
family \(\mathfrak A_{\rm REG}\); the occurring Hurwitz/Hopf branch uses the derived \((\mathrm{REG}_0)\) of Theorem~\ref{p3c:thm:occurring-hurwitz-reg-origin}. Only the unanchored branch retains an unresolved affine coordinate.
The finite-Hilbert class \(\Dcl\) is not used to prove that smooth zeta/determinant
uniqueness.  The neutral character, its clocked dual-line ratios, and its
unique positive associated fourth root are generated outputs.  Once a typed
occurrence supplies this character as the determinant of Paper~III-A's causal
generator bundle, the rank-four action \(H^{1/4}I_4\) and one-quarter
connection are forced; the occurrence of that cross-role ancestry is not.
Their action on a physical scale is inherited only after
a generated same-owner ancestry map is proved; on that branch no second
normalization is possible.  Thus \((\mathrm{HT})\) is only a
certification and cannot substitute for ancestry.  The
associated-character atlas \(\{L^q:q\in\mathbb Z\}\) likewise remains
unselected upstream; the present primitive oriented associated-line comparison
role derives \(q=+1\) only as a chronological later specialization and makes
no universal matter-charge claim.  The unconditional native export is the
neutral ratio \(\widehat\chi_{\sigma,w}(j)=\alpha_\sigma^{-wj}\).  The physical ratio
\(X_{w,j}/X_0\) is the owner-typed corollary; neither the dimensionful
reference \(X_0\) nor a cosmological/particle event label is selected here.

For the generated level-one momentum-seed comparison, the associated window \(\mathcal S_1\) has only
\(k=0\), so there is a single associated-spinor Dirac threshold.  The generic
\(\lfloor m-\tfrac12\rfloor\) staircase is an auxiliary finite-window family
recorded in Paper~III-G.  The \(E_6\) denominator identity belongs to that companion, and the Paper-I-B
\(56\)-dimensional resolution-data module remains a supplementary compatibility
readout in its separately declared domain.  Physical reciprocal-role
statements continue to consume the premise-minimal Paper-II-C/GP2
reciprocal-role closure; no \(\mathsf P_{\rm mc}\) or independent
\( (\mathrm{BR}_{\rm ind}) \)
or \( (\mathrm{IND}_{xp}) \) assumption remains in Paper~III-C.
All actual extractions, refinements, process constructors, owner transfers,
and physical readouts are typed under G0, while \((\mathrm{REG})\) remains an
explicit analytic domain rather than a consequence of geometric primacy.

The fixed-shell scheme family is an affine \((\mathbb R,+)\)-torsor, so no
translation-natural rule selects \(\mathrm{REG}_0\).  Its canonical content is
the exact translation-groupoid continuation and flat determinant-weight
connection of
Theorem~\ref{p3c:thm:exact-scheme-continuation-quotient}.  A total-weight-zero
scalar descendant is invariant on the scheme quotient.  More generally, a
nonzero-weight associated descendant has canonical covariant continuation
over the torsor through its flat determinant-weight connection.  Neither
statement selects a physical member.  A raw numerical evaluation still
requires a renormalization boundary or owner comparison.  On the full scheme--depth
base, all positive weights remain flat and gauge-equivalent; only a morphism
induced by the same fundamental associated determinant representation fixes
weight one.  This criterion does not itself identify an independently
constituted physical owner.



\printbibliography
\end{document}
